% La gestion du budget et des projets — Allemand Tle A — Cours signé OnBuch+
% Programme officiel MINESEC. Compiler avec : tectonic AL06-la-gestion-du-budget-et-des-projets.tex
\newcommand{\DOCMATIERE}{Allemand}
\newcommand{\DOCNIVEAU}{Tle A}
\newcommand{\DOCMODULE}{Chapitre 2 : Vie économique}
\newcommand{\DOCLECON}{AL06}
\newcommand{\DOCTITRE}{La gestion du budget et des projets}
% OnBuch+ — préambule « cours signés » ENRICHI (compilé avec tectonic / XeLaTeX)
% Système visuel « Pop » d'OnBuch+ : fond crème (jamais blanc), bordures encre
% épaisses, ombres « dures » décalées, titres Archivo Black en capitales.
% Version MATHÉMATIQUES : graphes (pgfplots/tikz), boîtes pédagogiques
% (définition, propriété, méthode, attention, le savais-tu, exercice résolu,
% exercice, TP, bilan), page de garde et signature OnBuch+ ancrée en bas.
%
% Macros attendues dans main.tex AVANT \input{preamble} :
%   \DOCMATIERE  \DOCNIVEAU  \DOCMODULE  \DOCLECON (numéro)  \DOCTITRE
\documentclass[11pt]{article}

\usepackage{fontspec}
\usepackage[ngerman,french]{babel} % français = langue principale ; l'allemand via \all{..} / textebox
\usepackage[a4paper,margin=2cm,top=3.4cm,bottom=2.8cm,headheight=1.4cm,footskip=1.3cm]{geometry}
\usepackage{amsmath,amssymb,amsfonts}
\usepackage{mathtools} % \xmapsto, \coloneqq... souvent utilisés par les modèles
\usepackage{cancel} % simplifications barrées (bilans de forces)
% \abs{z} (module, valeur absolue) et \norm{u} (norme), utilisés par les modèles
\providecommand{\abs}{}\let\abs\relax\DeclarePairedDelimiter{\abs}{\lvert}{\rvert}
\providecommand{\norm}{}\let\norm\relax\DeclarePairedDelimiter{\norm}{\lVert}{\rVert}
\usepackage[table]{xcolor} % [table] : fournit aussi \rowcolors (alternance de lignes)
\usepackage{pagecolor}
\usepackage{tikz}
\usetikzlibrary{arrows.meta,positioning,calc,shapes.geometric,shapes.misc,shapes.arrows,decorations.pathmorphing,decorations.pathreplacing,decorations.markings,patterns,shadows,mindmap,trees,fit,backgrounds,matrix,intersections,angles,quotes}
\usepackage{pgfplots}
\usepackage[version=4]{mhchem} % équations nucléaires \ce{^{235}_{92}U}
\usepackage[european,siunitx]{circuitikz} % schémas électriques
\pgfplotsset{compat=1.18}
\usepgfplotslibrary{fillbetween,groupplots} % aires entre courbes (calcul intégral)
\usepackage{enumitem}
\usepackage{fancyhdr}
% [most] charge toutes les bibliothèques tcolorbox (listings, minted, raster,
% documentation...) et gonfle inutilement la mémoire TeX sur un document long
% (~300 boîtes) : on ne charge que ce qui est réellement utilisé.
\usepackage[skins,breakable]{tcolorbox}
\usepackage{graphicx}
\usepackage{colortbl}
\usepackage{booktabs}
\usepackage{tabularx}
\usepackage{diagbox} % case d'en-tête diagonale des tableaux à double entrée
% Colonnes extensibles centrée / gauche / droite (C, L, R), souvent utilisées par les modèles
\newcolumntype{C}{>{\centering\arraybackslash}X}
\newcolumntype{L}{>{\raggedright\arraybackslash}X}
\newcolumntype{R}{>{\raggedleft\arraybackslash}X}
\usepackage{array}
\usepackage{multicol}
\usepackage{multirow}
\usepackage{siunitx}
% parse-numbers=false : accepte aussi \SI{2,5 \times 10^{-3}}{...}
\sisetup{locale=FR,output-decimal-marker={,},group-minimum-digits=4,parse-numbers=false}
\DeclareSIUnit\ohm{\ensuremath{\symup{\Omega}}} % U+2126 absent de la police
\DeclareSIUnit\division{div} % réglages d'oscilloscope (ms/div, V/div)
\DeclareSIUnit\tr{tr} % tours (vitesse de rotation, ex. \SI{1500}{\tr\per\minute})
\DeclareSIUnit\molar{M} % concentration molaire (ex. \SI{3}{\molar}, \SI{1}{\micro\molar})
\DeclareSIUnit\an{an} % année (le modèle confond parfois avec \year, un compteur TeX) — ex. \SI{5}{\kilo\meter\per\an}
\DeclareSIUnit\FCFA{FCFA} % franc CFA (ex. \SI{500}{\FCFA\per\kilo\gram})
\DeclareSIUnit\degreeBrix{\textdegree Brix} % teneur en sucre d'un sirop/jus (réfractométrie)
\DeclareSIUnit\month{mois} % le modèle utilise parfois \month, un compteur TeX, comme unité de durée
\usepackage{qrcode}
% \degree : commande fréquemment utilisée par les modèles pour un angle ou une
% température en dehors de \SI{}{} (non fournie par les paquets chargés).
\providecommand{\degree}{^{\circ}}
% \qed (fin de démonstration), utilisé par les modèles sans amsthm
\providecommand{\qed}{\hfill$\square$}
% \barre : double barre des valeurs interdites dans un tableau de variations (tkz-tab)
\providecommand{\barre}{\ensuremath{\|}}
% environnement proof (amsthm non chargé) utilisé par les modèles
\ifdefined\proof\else\newenvironment{proof}[1][Démonstration]{\par\noindent\textit{#1.}\ }{\qed\par}\fi
\usepackage{lastpage}
\usepackage{hyperref}
\hypersetup{hidelinks}

% ============================================================
% Palette « Pop » OnBuch+ — mêmes hex que la classe P (plus_ui.dart)
% ============================================================
\definecolor{poporange}{HTML}{F59321}
\definecolor{popdark}{HTML}{E07A0C}
\definecolor{popcream}{HTML}{FFF6E8}
\definecolor{popink}{HTML}{1C1714}
\definecolor{poppurple}{HTML}{7A5AE0}
\definecolor{popgreen}{HTML}{2E9E6A}
\definecolor{poppink}{HTML}{E0457B}
\definecolor{popblue}{HTML}{2F7DE1}
\definecolor{popgold}{HTML}{C98A12}
\definecolor{popmuted}{HTML}{8A7F76}
\definecolor{popline}{HTML}{EADFD2}
\definecolor{popgray}{HTML}{8A7F76}
\colorlet{popgrey}{popgray}
% Alias tolérants (le modèle nomme parfois les couleurs différemment)
\colorlet{popwhite}{white}
\colorlet{popblack}{popink}
\colorlet{popred}{poppink}
\colorlet{popyellow}{popgold}
% Teintes claires (fonds de tableaux/encadrés) — le modèle nomme parfois
% les couleurs avec un suffixe L (ex. popblueL) sans les avoir définies.
\colorlet{poporangeL}{poporange!20}
\colorlet{popdarkL}{popdark!20}
\colorlet{poppurpleL}{poppurple!20}
\colorlet{popgreenL}{popgreen!20}
\colorlet{poppinkL}{poppink!20}
\colorlet{popblueL}{popblue!20}
\colorlet{popgoldL}{popgold!20}
\colorlet{popredL}{popred!20}
\colorlet{popgrayL}{popgray!20}
% Teintes claires pour fonds de tableaux / zones
\colorlet{popblueL}{popblue!10!white}
\colorlet{poporangeL}{poporange!14!white}
\colorlet{popgreenL}{popgreen!12!white}
\colorlet{poppurpleL}{poppurple!12!white}
\colorlet{poppinkL}{poppink!12!white}
\colorlet{popgoldL}{popgold!14!white}

\pagecolor{popcream}

% ============================================================
% Typographie
% ============================================================
\defaultfontfeatures{Ligatures=TeX}
\setmainfont{PlusJakartaSans-Medium.ttf}[Path=../../../fonts/,BoldFont=PlusJakartaSans-Bold.ttf]
% Symboles, API et emojis absents de la police principale : repli sur DejaVu Sans (caractères actifs)
\newfontface\symfont{DejaVuSans.ttf}[Path=../../../fonts/]
\makeatletter
\newcommand\symactive[1]{\catcode#1=13 \begingroup\lccode`\~=#1 \lowercase{\endgroup\def~}{{\symfont\char#1}}}
\newcommand\symblank[1]{\catcode#1=13 \begingroup\lccode`\~=#1 \lowercase{\endgroup\def~}{}}
\newcommand\symhyph[1]{\catcode#1=13 \begingroup\lccode`\~=#1 \lowercase{\endgroup\def~}{\mbox{-}}}
\newcount\symc
\newcommand\symrange[2]{\symc=#1 \@whilenum\symc<#2 \do{\expandafter\symactive\expandafter{\the\symc}\advance\symc by 1 }}
\newcommand\symblankrange[2]{\symc=#1 \@whilenum\symc<#2 \do{\expandafter\symblank\expandafter{\the\symc}\advance\symc by 1 }}
\makeatother
\symrange{"0250}{"02FF}\symactive{"2713}\symactive{"2714}\symactive{"2717}\symactive{"2718}\symactive{"2610}\symactive{"2611}\symactive{"2612}
\symhyph{"2011}\symhyph{"2E1A}\symblankrange{"1F300}{"1FAFF}\symblank{"FE0F}
% Maths en sans-serif (Fira Math) avec chiffres et lettres droites de Plus
% Jakarta, pour que les formules se fondent dans le texte.
\usepackage[math-style=ISO,bold-style=ISO]{unicode-math}
\setmathfont{FiraMath-Regular.otf}[Path=../../../fonts/]
\setmathfont{PlusJakartaSans-Medium.ttf}[Path=../../../fonts/,range={up/num,up/{Latin,latin},\mathup}]
\usepackage{icomma} % virgule décimale sans espace en mode math
% Caractères mathématiques Unicode absents des polices de texte (Plus Jakarta,
% Archivo Black) : rendus par leur équivalent mathématique, en texte comme en
% formule — sinon ils s'affichent en carré vide (ex. « Arithmétique dans ℤ »).
\usepackage{newunicodechar}
\newunicodechar{ℝ}{\ensuremath{\mathbb{R}}}
\newunicodechar{ℤ}{\ensuremath{\mathbb{Z}}}
\newunicodechar{ℂ}{\ensuremath{\mathbb{C}}}
\newunicodechar{ℕ}{\ensuremath{\mathbb{N}}}
\newunicodechar{ℚ}{\ensuremath{\mathbb{Q}}}
\newunicodechar{𝔻}{\ensuremath{\mathbb{D}}}
\newunicodechar{α}{\ensuremath{\alpha}}
\newunicodechar{β}{\ensuremath{\beta}}
\newunicodechar{γ}{\ensuremath{\gamma}}
\newunicodechar{δ}{\ensuremath{\delta}}
\newunicodechar{ε}{\ensuremath{\varepsilon}}
\newunicodechar{θ}{\ensuremath{\theta}}
\newunicodechar{λ}{\ensuremath{\lambda}}
\newunicodechar{μ}{\ensuremath{\mu}}
\newunicodechar{σ}{\ensuremath{\sigma}}
\newunicodechar{τ}{\ensuremath{\tau}}
\newunicodechar{φ}{\ensuremath{\varphi}}
\newunicodechar{ψ}{\ensuremath{\psi}}
\newunicodechar{ω}{\ensuremath{\omega}}
\newunicodechar{Σ}{\ensuremath{\Sigma}}
% majuscules produites par \MakeUppercase dans les titres (α-aminés -> Α)
\newunicodechar{Α}{\ensuremath{\alpha}}
\newunicodechar{Β}{\ensuremath{\beta}}
\newunicodechar{Γ}{\ensuremath{\Gamma}}
\newunicodechar{Δ}{\ensuremath{\Delta}}
\newunicodechar{Ε}{\ensuremath{\varepsilon}}
\newunicodechar{Θ}{\ensuremath{\Theta}}
\newunicodechar{Λ}{\ensuremath{\Lambda}}
\newunicodechar{Μ}{\ensuremath{\mu}}
\newunicodechar{Τ}{\ensuremath{\tau}}
\newunicodechar{Φ}{\ensuremath{\Phi}}
\newunicodechar{Ψ}{\ensuremath{\Psi}}
\newunicodechar{Ω}{\ensuremath{\Omega}}
\newunicodechar{↦}{\ensuremath{\mapsto}}
\newunicodechar{∈}{\ensuremath{\in}}
\newunicodechar{∉}{\ensuremath{\notin}}
\newunicodechar{∧}{\ensuremath{\wedge}}
\newunicodechar{⇔}{\ensuremath{\Leftrightarrow}}
\newunicodechar{⟺}{\ensuremath{\Longleftrightarrow}}
\newunicodechar{⇒}{\ensuremath{\Rightarrow}}
\newunicodechar{⟹}{\ensuremath{\Longrightarrow}}
\newunicodechar{∘}{\ensuremath{\circ}}
\newunicodechar{∪}{\ensuremath{\cup}}
\newunicodechar{∩}{\ensuremath{\cap}}
\newunicodechar{≡}{\ensuremath{\equiv}}
\newunicodechar{⊂}{\ensuremath{\subset}}
\newunicodechar{⊥}{\ensuremath{\perp}}
\newunicodechar{∃}{\ensuremath{\exists}}
\newunicodechar{∀}{\ensuremath{\forall}}
\newunicodechar{∛}{\ensuremath{\sqrt[3]{\,}}}
\newunicodechar{∜}{\ensuremath{\sqrt[4]{\,}}}
\newunicodechar{⃗}{\ensuremath{{}^{\to}}}
\newunicodechar{ⁿ}{\ifmmode^{n}\else\textsuperscript{\ensuremath{n}}\fi}
\newunicodechar{ˣ}{\ifmmode^{x}\else\textsuperscript{\ensuremath{x}}\fi}
\newunicodechar{ᵇ}{\ifmmode^{b}\else\textsuperscript{\ensuremath{b}}\fi}
\newunicodechar{ʳ}{\ifmmode^{r}\else\textsuperscript{\ensuremath{r}}\fi}
\newunicodechar{ᵘ}{\ifmmode^{u}\else\textsuperscript{\ensuremath{u}}\fi}
\newunicodechar{ʸ}{\ifmmode^{y}\else\textsuperscript{\ensuremath{y}}\fi}
\newunicodechar{ᵃ}{\ifmmode^{a}\else\textsuperscript{\ensuremath{a}}\fi}
\newunicodechar{ᵉ}{\ifmmode^{e}\else\textsuperscript{\ensuremath{e}}\fi}
\newunicodechar{ᵏ}{\ifmmode^{k}\else\textsuperscript{\ensuremath{k}}\fi}
\newunicodechar{ᵖ}{\ifmmode^{p}\else\textsuperscript{\ensuremath{p}}\fi}
\newunicodechar{ᴺ}{\ifmmode^{N}\else\textsuperscript{\ensuremath{N}}\fi}
\newunicodechar{ᶜ}{\ifmmode^{c}\else\textsuperscript{\ensuremath{c}}\fi}
\newunicodechar{ʰ}{\ifmmode^{h}\else\textsuperscript{\ensuremath{h}}\fi}
\newunicodechar{ᵠ}{\ifmmode^{q}\else\textsuperscript{\ensuremath{q}}\fi}
\newunicodechar{ₙ}{\ifmmode_{n}\else\textsubscript{\ensuremath{n}}\fi}
\newunicodechar{ₐ}{\ifmmode_{a}\else\textsubscript{\ensuremath{a}}\fi}
\newunicodechar{ᵢ}{\ifmmode_{i}\else\textsubscript{\ensuremath{i}}\fi}
\newunicodechar{ᵣ}{\ifmmode_{r}\else\textsubscript{\ensuremath{r}}\fi}
\newunicodechar{ₖ}{\ifmmode_{k}\else\textsubscript{\ensuremath{k}}\fi}
\newunicodechar{ᵦ}{\ifmmode_{\beta}\else\textsubscript{\ensuremath{\beta}}\fi}
\newunicodechar{ₓ}{\ifmmode_{x}\else\textsubscript{\ensuremath{x}}\fi}
\newfontfamily\popdisplay{ArchivoBlack-Regular.ttf}[Path=../../../fonts/]
\newfontfamily\poplogo{SpaceGrotesk-Bold.ttf}[Path=../../../fonts/]
\newfontfamily\popsemi{PlusJakartaSans-SemiBold.ttf}[Path=../../../fonts/]
\newfontfamily\popxbold{PlusJakartaSans-ExtraBold.ttf}[Path=../../../fonts/]
\newfontfamily\popxboldls{PlusJakartaSans-ExtraBold.ttf}[Path=../../../fonts/,LetterSpace=3.0]

\setlist[enumerate]{leftmargin=1.7em,itemsep=5pt,topsep=4pt}
\setlist[itemize]{leftmargin=1.7em,itemsep=4pt,topsep=4pt,label={\color{poporange}\textbullet}}
\setlength{\parindent}{0pt}
\setlength{\parskip}{5pt}
\renewcommand{\arraystretch}{1.25}

% pgfplots : style Pop par défaut
\pgfplotsset{
  every axis/.append style={axis line style={popink,line width=1pt},tick style={popink},
    grid style={popline,line width=0.5pt},label style={font=\small},tick label style={font=\footnotesize}},
  popcurve/.style={poporange,line width=1.8pt,no markers,smooth},
}
% Même style disponible dans un \draw TikZ simple (hors pgfplots)
\tikzset{popcurve/.style={poporange,line width=1.8pt}}
\tikzset{popgrid/.style={popline,very thin}}
\colorlet{popcurve}{poporange} % le nom du style est parfois utilisé comme couleur

% ============================================================
% Pill / squiggle
% ============================================================
\newcommand{\poppill}[3]{%
  \tikz[baseline=-0.55ex]\node[draw=popink,line width=1.2pt,rounded corners=2.6pt,%
    fill=#2,inner xsep=6.5pt,inner ysep=3pt]{%
    \popxboldls\fontsize{7}{7}\selectfont\color{#3}\MakeUppercase{#1}};%
}
\newcommand{\popsquiggle}[1]{%
  \tikz[baseline=-0.4ex]{
    \draw[#1,line width=2.6pt,line cap=round]
      plot [smooth] coordinates {(0,0.09) (0.34,0.01) (0.68,0.21) (1.02,0.01) (1.36,0.10)};
  }%
}
% Mot-clé surligné (marqueur orange sous le texte)
\newcommand{\cle}[1]{{\popxbold\color{popdark}#1}}

% ============================================================
% En-tête / pied
% ============================================================
\pagestyle{fancy}
\fancyhf{}
\renewcommand{\headrulewidth}{0pt}
\renewcommand{\footrulewidth}{0pt}
\lhead{\raisebox{-0.15ex}{\poplogo\fontsize{13}{13}\selectfont\color{popink}OnBuch\color{poporange}+}}
\rhead{\poppill{\DOCMATIERE\ \textbullet\ \DOCNIVEAU\ \textbullet\ Leçon \DOCLECON}{white}{popink}}
\lfoot{\popxbold\fontsize{7.6}{7.6}\selectfont\color{popmuted}\MakeUppercase{Cours signé OnBuch+}\ \textbullet\ {\popsemi\DOCTITRE}}
\rfoot{\tikz[baseline=-0.6ex]\node[rounded corners=8.5pt,fill=popink,minimum height=17pt,minimum width=17pt,inner xsep=5pt,inner sep=0]{\popxbold\fontsize{8}{8}\selectfont\color{popcream}\thepage\,/\,\pageref*{LastPage}};}

% ============================================================
% Titres
% ============================================================
\newcommand{\chaptitre}[2]{%
  \begin{tcolorbox}[enhanced,colback=poporange,colframe=popink,boxrule=2.6pt,arc=11pt,
      drop shadow=popink,left=15pt,right=15pt,top=12pt,bottom=11pt,before skip=3pt,after skip=18pt]
    \poppill{#2}{popink}{popcream}\par\vspace{9pt}
    {\raggedright\hyphenpenalty=10000\exhyphenpenalty=10000\popdisplay\fontsize{22}{24}\selectfont\color{popcream}\MakeUppercase{#1}\par}
  \end{tcolorbox}
}
% Section numérotée : pastille orange avec le numéro + titre Archivo
\newcounter{popsec}
\newcommand{\coursec}[1]{%
  \stepcounter{popsec}\setcounter{popsub}{0}%
  \par\vspace{16pt}\noindent
  \begin{minipage}{\linewidth}
  \tikz[baseline=-0.7ex]\node[draw=popink,line width=1.6pt,fill=poporange,rounded corners=4pt,minimum size=22pt,inner sep=0]{\popdisplay\fontsize{12}{12}\selectfont\color{popcream}\thepopsec};%
  \hspace{8pt}{\popdisplay\fontsize{15}{17}\selectfont\color{popink}\MakeUppercase{#1}}\par
  \vspace{4pt}\noindent{\color{poporange}\rule{36pt}{3.6pt}}
  \end{minipage}\par\nopagebreak\vspace{8pt}
}
\newcounter{popsub}
\newcommand{\courssub}[1]{%
  \stepcounter{popsub}%
  \par\vspace{9pt}\noindent{\popxbold\fontsize{11.5}{13}\selectfont\color{popink}{\color{poporange}\thepopsec.\thepopsub}\hspace{6pt}#1}\par\nopagebreak\vspace{4pt}
}

% ============================================================
% Boîtes pédagogiques — même recette : fond blanc, bordure épaisse colorée,
% ombre dure encre, badge plein. Titre optionnel [#1].
% ============================================================
\tcbset{popbase/.style={breakable,enhanced,colback=white,boxrule=2.3pt,arc=9pt,
  shadow={3pt}{-3pt}{0pt}{popink},left=14pt,right=14pt,top=11pt,bottom=11pt,before skip=11pt,after skip=15pt,
  fontupper=\popsemi\fontsize{10.6}{14.5}\selectfont\color{popink},
  fontlower=\popsemi\fontsize{10.6}{14.5}\selectfont\color{popink}}}
% Coche dessinée (le glyphe ✓ n'existe pas dans les polices du document)
\newcommand{\popcheck}[1]{\tikz[baseline=-0.5ex]\draw[#1,line width=1.6pt,line cap=round,line join=round](-0.13,0.0)--(-0.04,-0.09)--(0.13,0.1);}
\providecommand{\coche}{\popcheck{popgreen}}
\providecommand{\resultat}[1]{\par\smallskip\noindent\fcolorbox{popgreen}{popgreenL}{\parbox{\dimexpr\linewidth-2\fboxsep-2\fboxrule\relax}{\textbf{Résultat.} #1}}\par\smallskip}
\newcommand{\popboxhead}[3]{\poppill{#1}{#2}{white}\ifx\relax#3\relax\else\hspace{6pt}{\popxbold\fontsize{10.6}{12}\selectfont\color{popink}#3}\fi\par\vspace{7pt}}

\newtcolorbox{aretenir}[1][]{popbase,colframe=popgreen,before upper={\popboxhead{À retenir}{popgreen}{#1}}}
% Texte en allemand (document de lecture, dialogue, extrait) : cadre
% d'étude. Un titre optionnel (source, auteur) s'affiche dans l'en-tête.
\newtcolorbox{textebox}[1][]{popbase,colframe=popblue,before upper={\popboxhead{Text}{popblue}{#1}\begin{otherlanguage}{ngerman}},after upper={\end{otherlanguage}}}
% Marques ✓ / ✗ (forme correcte / fautive) souvent utilisées par les modèles
\providecommand{\cross}{\textcolor{poppink}{\ensuremath{\times}}}
\providecommand{\xmark}{\cross}\providecommand{\crossmark}{\cross}\providecommand{\cmark}{\ensuremath{\checkmark}}
% Ligne de réponse / justification à compléter (inventée par les modèles)
\providecommand{\justification}[1]{\par\noindent\textit{Justification~:} \dotfill\par}
% \textipa (package tipa non chargé) : la police ne couvre pas l'API -> simple passage
\providecommand{\textipa}[1]{#1}
% Soulignement (ulem non chargé) : \uline, \ul
\providecommand{\uline}[1]{\underline{#1}}\providecommand{\ul}[1]{\underline{#1}}
% \glossaire{\item ...} inventée par les modèles : liste de glossaire
\providecommand{\glossaire}[1]{\begin{itemize}#1\end{itemize}}
% \alert (beamer) inventée par les modèles : texte mis en évidence
\providecommand{\alert}[1]{\textbf{\textcolor{poppink}{#1}}}
% variantes ASCII d'ordinaux écrites en macro
\providecommand{\iere}{\textsuperscript{re}}\providecommand{\ieme}{\textsuperscript{e}}\providecommand{\ere}{\textsuperscript{re}}\providecommand{\eme}{\textsuperscript{e}}\providecommand{\er}{\textsuperscript{er}}\providecommand{\ie}{\textsuperscript{e}}
% Ordinaux français écrits en macro par les modèles : 1\ière, 3\ième
\newcommand{\ière}{\textsuperscript{re}}\newcommand{\ième}{\textsuperscript{e}}
% \exemple (inventée par les modèles) : étiquette « Exemple : »
\providecommand{\exemple}{\textbf{Exemple~:}\ }
% \square utilisable aussi hors mode math (cases à cocher)
\AtBeginDocument{\let\squareMath\square\renewcommand{\square}{\ensuremath{\squareMath}}}
% Mot ou phrase en allemand à l'intérieur d'un paragraphe français
\newcommand{\all}[1]{\ifmmode\text{\textit{#1}}\else\begingroup\selectlanguage{ngerman}\itshape #1\endgroup\fi}
\let\esp\all % alias toléré
\newtcolorbox{exemplebox}[1][]{popbase,colframe=poporange,before upper={\popboxhead{Exemple}{poporange}{#1}}}
\newtcolorbox{definition}[1][]{popbase,colframe=popblue,before upper={\popboxhead{Définition}{popblue}{#1}}}
\newtcolorbox{propriete}[1][]{popbase,colframe=poppurple,before upper={\popboxhead{Propriété}{poppurple}{#1}}}
\newtcolorbox{methode}[1][]{popbase,colframe=popgold,before upper={\popboxhead{Méthode}{popgold}{#1}}}
\newtcolorbox{attention}[1][]{popbase,colframe=poppink,before upper={\popboxhead{Attention}{poppink}{#1}}}
\newtcolorbox{experience}[1][]{popbase,colframe=popblue,colback=popblueL,before upper={\popboxhead{Expérience}{popblue}{#1}}}
% « Le savais-tu ? » : carte violette pleine (contexte, culture, Cameroun)
\newtcolorbox{savaistu}[1][]{popbase,colback=poppurple,colframe=popink,
  fontupper=\popsemi\fontsize{10.4}{14.2}\selectfont\color{white},
  before upper={\poppill{Le savais-tu\,?}{white}{poppurple}\ifx\relax#1\relax\else\hspace{6pt}{\popxbold\color{white}#1}\fi\par\vspace{7pt}}}
% Exercice résolu : énoncé en haut, \tcblower puis la solution sur fond crème
\newcounter{popexo}\newcounter{popexr}
\newtcolorbox{exoresolu}[1][]{popbase,colframe=poporange,colbacklower=popcream,
  segmentation style={popink,line width=1.2pt,dashed},
  before upper={\stepcounter{popexr}\popboxhead{Exercice résolu \thepopexr}{poporange}{#1}},
  before lower={\poppill{Solution}{popink}{popcream}\par\vspace{6pt}}}
% Exercice (non résolu en place) — la correction est dans la partie Corrigés
\newtcolorbox{exercice}[1][]{popbase,colframe=popink,
  before upper={\stepcounter{popexo}\popboxhead{Exercice \thepopexo}{popink}{#1}}}
% Corrigé
\newtcolorbox{corrige}[1][]{popbase,colframe=popgreen,colback=popgreenL,
  before upper={\popboxhead{Corrigé}{popgreen}{#1}}}
% Objectifs / prérequis / bilan
\newtcolorbox{objectifs}{popbase,colframe=popink,colback=poporangeL,
  before upper={\popboxhead{Objectifs de la leçon}{popink}{}}}
\newtcolorbox{prerequis}{popbase,colframe=popink,colback=popblueL,
  before upper={\popboxhead{Prérequis}{popblue}{}}}
\newtcolorbox{bilan}{popbase,colframe=popink,colback=white,boxrule=2.8pt,
  before upper={\popboxhead{Fiche bilan}{poporange}{L'essentiel à connaître pour l'examen}}}
% Figure encadrée avec légende
\newtcolorbox{popfigure}[1][]{enhanced,breakable=false,colback=white,colframe=popline,boxrule=1.4pt,arc=8pt,
  left=10pt,right=10pt,top=10pt,bottom=8pt,before skip=10pt,after skip=12pt,halign=center,
  lower separated=false,halign lower=center,
  fontlower=\popsemi\fontsize{9}{11.5}\selectfont\color{popmuted},
  #1}
\newcommand{\legende}[1]{\par\vspace{6pt}{\popsemi\fontsize{9}{11.5}\selectfont\color{popmuted}#1}}

% ============================================================
% Signature OnBuch+
% ============================================================
\newcommand{\poplogoinline}[1]{{\poplogo\fontsize{#1}{#1}\selectfont\color{popink}OnBuch\color{poporange}+}}
\newcommand{\signatureonbuch}{%
  \begin{tcolorbox}[enhanced,colback=popink,colframe=popink,boxrule=2.2pt,arc=10pt,
      drop shadow=poporange,left=14pt,right=14pt,top=10pt,bottom=10pt,before skip=0pt,after skip=0pt]
    \tikz[baseline=-0.7ex]\node[circle,fill=popgreen,draw=popcream,line width=1.3pt,minimum size=22pt,inner sep=0]{\popcheck{white}};%
    \hspace{9pt}\begin{minipage}[c]{0.62\linewidth}
      {\popxbold\fontsize{12}{13}\selectfont\color{popcream}Signé\ \textbullet\ {\poplogo OnBuch\color{poporange}+}}\par\vspace{2pt}
      {\raggedright\popsemi\fontsize{8}{10.5}\selectfont\color{popline}\MakeUppercase{Équipe pédagogique OnBuch+}\\\MakeUppercase{Conforme au programme officiel MINESEC}\par}
    \end{minipage}\hfill
    {\poplogo\fontsize{22}{22}\selectfont\color{popcream}OnBuch\color{poporange}+}
  \end{tcolorbox}%
}

% ============================================================
% Page de garde
% #1 titre · #2 matière · #3 niveau · #4 module · #5 n° leçon · #6 durée ·
% #7 description / objectifs (texte ou itemize)
% ============================================================
\newcommand{\pagegarde}[7]{%
  \thispagestyle{empty}
  \newgeometry{a4paper,margin=1.7cm,top=1.6cm,bottom=1.4cm}%
  \begin{tikzpicture}[remember picture,overlay]
    \fill[poporange!18] ([xshift=-3.2cm,yshift=-3.6cm]current page.north east) circle (4.6cm);
    \fill[poppurple!14] ([xshift=2.4cm,yshift=7.6cm]current page.south west) circle (3.8cm);
  \end{tikzpicture}%
  {\poplogo\fontsize{30}{30}\selectfont\color{popink}OnBuch\color{poporange}+}\hfill
  \raisebox{0.6ex}{\poppill{Cours signé \textbullet\ Premium}{popink}{popcream}}\par
  \vspace{1pt}\popsquiggle{poppurple}\par
  \vspace{8pt}
  \begin{tcolorbox}[enhanced,colback=poporange,colframe=popink,boxrule=2.8pt,arc=13pt,
      drop shadow=popink,left=18pt,right=18pt,top=12pt,bottom=13pt,after skip=10pt]
    \poppill{#2 \textbullet\ #3}{popink}{popcream}\hspace{5pt}\poppill{#4}{white}{popink}\par\vspace{8pt}
    {\popdisplay\fontsize{14}{14}\selectfont\color{popink}LEÇON #5}\par\vspace{4pt}
    {\raggedright\hyphenpenalty=10000\exhyphenpenalty=10000\popdisplay\fontsize{25}{27}\selectfont\color{popcream}\MakeUppercase{#1}\par}
    \vspace{8pt}
    {\popxbold\fontsize{9.5}{9.5}\selectfont\color{popink}\MakeUppercase{Durée conseillée : #6}}
  \end{tcolorbox}
  \begin{tcolorbox}[enhanced,colback=white,colframe=popink,boxrule=2.2pt,arc=10pt,
      drop shadow=popink,left=16pt,right=16pt,top=10pt,bottom=10pt,after skip=8pt]
    \poppill{Dans cette leçon}{poppurple}{white}\par\vspace{5pt}
    \popsemi\fontsize{9.3}{12.6}\selectfont\color{popink}#7
  \end{tcolorbox}
  \noindent\begin{minipage}[t]{0.487\linewidth}
    \begin{tcolorbox}[enhanced,colback=popgreen,colframe=popink,boxrule=2.2pt,arc=10pt,
        drop shadow=popink,left=11pt,right=11pt,top=10pt,bottom=10pt,height=3.1cm,valign=center]
      \tcbox[colback=white,colframe=popink,boxrule=1.3pt,arc=5pt,boxsep=0pt,nobeforeafter,
        left=3pt,right=3pt,top=3pt,bottom=3pt,valign=center]{\qrcode[height=1.9cm]{https://play.google.com/store/apps/details?id=com.luvvix.onbuch&hl=fr}}%
      \hspace{8pt}\begin{minipage}[c]{0.5\linewidth}
        {\popdisplay\fontsize{11}{12.5}\selectfont\color{white}\MakeUppercase{Télécharge OnBuch}}\par\vspace{4pt}
        {\raggedright\popsemi\fontsize{8.4}{11}\selectfont\color{white}Gratuit \textbullet\ Google Play\\Tuteur IA, annales, concours, résultats\par}
      \end{minipage}
    \end{tcolorbox}
  \end{minipage}\hfill
  \begin{minipage}[t]{0.487\linewidth}
    \begin{tcolorbox}[enhanced,colback=poppurple,colframe=popink,boxrule=2.2pt,arc=10pt,
        drop shadow=popink,left=11pt,right=11pt,top=10pt,bottom=10pt,height=3.1cm,valign=center]
      \tikz[baseline=-0.7ex]\node[circle,fill=white,draw=popink,line width=1.3pt,minimum size=1.6cm,inner sep=0]{\popdisplay\fontsize{24}{24}\selectfont\color{poppurple}+};%
      \hspace{8pt}\begin{minipage}[c]{0.6\linewidth}
        {\popdisplay\fontsize{11}{12.5}\selectfont\color{white}\MakeUppercase{Passe à OnBuch+}}\par\vspace{4pt}
        {\raggedright\popsemi\fontsize{8.4}{11}\selectfont\color{white}Tous les cours signés, Tuteur IA illimité, fiches PDF — onglet OnBuch+ de l'app\par}
      \end{minipage}
    \end{tcolorbox}
  \end{minipage}
  \vspace{8pt}
  \popcontenu
  \vfill
  \signatureonbuch
  \restoregeometry
  \newpage
}

% Rangée « Ce cours contient » (page de garde)
\newcommand{\poptile}[3]{% #1 couleur #2 gros texte #3 légende
  \begin{tcolorbox}[enhanced,colback=#1,colframe=popink,boxrule=2pt,arc=9pt,shadow={2.5pt}{-2.5pt}{0pt}{popink},
      left=6pt,right=6pt,top=9pt,bottom=9pt,halign=center,valign=center,height=1.75cm,width=\linewidth,nobeforeafter]
    {\popdisplay\fontsize{12.5}{13}\selectfont\color{white}\MakeUppercase{#2}}\par\vspace{3pt}
    {\centering\popsemi\fontsize{7.8}{9.5}\selectfont\color{white}#3\par}
  \end{tcolorbox}}
\newcommand{\popcontenu}{%
  \par\noindent{\popxboldls\fontsize{7.5}{7.5}\selectfont\color{popmuted}CE COURS SIGNÉ CONTIENT}\par\vspace{7pt}
  \noindent\begin{minipage}[t]{0.235\linewidth}\poptile{popblue}{Cours}{complet, illustré, pas à pas}\end{minipage}\hfill
  \begin{minipage}[t]{0.235\linewidth}\poptile{poporange}{Méthodes}{et exercices résolus}\end{minipage}\hfill
  \begin{minipage}[t]{0.235\linewidth}\poptile{poppink}{Exercices}{type examen + corrigés}\end{minipage}\hfill
  \begin{minipage}[t]{0.235\linewidth}\poptile{popgreen}{Fiche bilan}{l'essentiel à retenir}\end{minipage}\par}

% Sommaire
\newtcolorbox{sommaire}{enhanced,colback=white,colframe=popink,boxrule=2.2pt,arc=10pt,
  drop shadow=popink,left=18pt,right=18pt,top=13pt,bottom=13pt,before skip=6pt,after skip=20pt,
  before upper={\poppill{Sommaire}{poppurple}{white}\par\vspace{9pt}\popsemi\fontsize{10.6}{14.5}\selectfont\color{popink}}}

% Signature de fin de cours — poussée tout en bas de la dernière page
\newcommand{\findecours}{%
  \par\vfill\nopagebreak
  \begin{center}\popsquiggle{poporange}\end{center}\vspace{4pt}
  \signatureonbuch
}
\AtBeginDocument{\def\llbracket{\mathopen{[\mkern-3.5mu[}}\def\rrbracket{\mathclose{]\mkern-3.5mu]}}} % intervalles d'entiers (glyphe ⟦ absent de la police)
% Glyphes absents de Fira Math : redéfinis à partir de glyphes disponibles
\AtBeginDocument{%
  \def\setminus{\mathbin{\backslash}}\let\smallsetminus\setminus
  \def\vdots{\mathord{\vbox{\baselineskip3.2pt\lineskiplimit0pt\kern4pt\hbox{.}\hbox{.}\hbox{.}}}}%
  \def\ddots{\mathinner{\mkern1mu\raise6.5pt\hbox{.}\mkern2mu\raise3.6pt\hbox{.}\mkern2mu\raise0.7pt\hbox{.}\mkern1mu}}%
  \def\triangle{\mathord{\scalebox{0.9}{$\symup{\Delta}$}}}%
  \def\nabla{\mathord{\raisebox{\depth}{\scalebox{1}[-1]{$\symup{\Delta}$}}}}%
  \def\checkmark{\popcheck{popdark}}%
  \def\star{\mathbin{\ast}}\def\bigstar{\mathord{\ast}}%
  \def\diamond{\mathbin{\scalebox{0.6}{$\Diamond$}}}%
  \def\bigcup{\mathop{\mathchoice{\raisebox{-0.3em}{\scalebox{1.7}{$\cup$}}}{\scalebox{1.25}{$\cup$}}{\cup}{\cup}}\displaylimits}%
  \def\bigcap{\mathop{\mathchoice{\raisebox{-0.3em}{\scalebox{1.7}{$\cap$}}}{\scalebox{1.25}{$\cap$}}{\cap}{\cap}}\displaylimits}%
  \def\lesseqqgtr{\mathrel{\lessgtr}}\def\bigoplus{\mathop{\oplus}\displaylimits}%
}
\providecommand{\ssi}{si et seulement si} % abréviation fréquente
\AtBeginDocument{\providecommand{\wideparen}{\overparen}} % arc de cercle

%% FIGFIX-BEGIN v3 — correctifs globaux des figures (docs/onbuch-generation/tools/figfix)
\usepackage{adjustbox}\usepackage{etoolbox}
% toute figure TikZ/pgfplots plus large que la ligne est réduite à la largeur disponible
\BeforeBeginEnvironment{tikzpicture}{\begin{adjustbox}{max width=\linewidth}}
\AfterEndEnvironment{tikzpicture}{\end{adjustbox}}
% légendes pgfplots lisibles par-dessus les courbes
\pgfplotsset{every axis/.append style={legend style={fill=white,fill opacity=0.92,text opacity=1,draw=black!15,inner sep=3pt}}}
% étiquettes d'arêtes (syntaxe edge["..."]) avec fond blanc
\tikzset{every edge quotes/.append style={fill=white,inner sep=1.2pt}}
% bulles de cartes mentales : texte à la ligne dans la bulle, bulle un peu plus grande
\tikzset{every concept/.append style={text width=2.0cm,align=center,minimum size=2.3cm,inner sep=2pt}}
%% FIGFIX-END
\begin{document}

\pagegarde{La gestion du budget et des projets}{Allemand}{Tle A}{Chapitre 2 : Vie économique}{AL06}{2 h}{Cette leçon de type « texte » propose la lecture et le commentaire d'un texte allemand sur la gestion du budget et des projets, avec un lexique thématique riche (Geld, Budget, Einnahmen, Ausgaben, sparen, Projekt, Plan). Elle consolide trois points de grammaire du programme : le comparatif et le superlatif, les subordonnées conditionnelles avec « wenn » et le Futur I, avant une production guidée.
\vspace{4pt}
\begin{multicols}{2}\small
\begin{itemize}
\item À la fin de cette leçon, je sais comprendre globalement puis en détail un texte allemand sur la gestion d'un budget et d'un projet.
\item À la fin de cette leçon, je sais employer correctement le vocabulaire thématique (das Budget, die Einnahmen, die Ausgaben, sparen, das Projekt, der Plan, das Geld) avec articles et pluriels.
\item À la fin de cette leçon, je sais former et utiliser le comparatif et le superlatif des adjectifs et adverbes allemands.
\item À la fin de cette leçon, je sais construire une subordonnée conditionnelle avec « wenn » en respectant la place du verbe.
\item À la fin de cette leçon, je sais conjuguer et employer le Futur I pour annoncer un projet ou faire une prévision.
\item À la fin de cette leçon, je sais répondre par écrit à des questions de compréhension en phrases complètes.
\end{itemize}\end{multicols}}

\begin{sommaire}
\begin{multicols}{2}
\begin{enumerate}
\item Texte support : « Lernen, mit Geld umzugehen »
\item Compréhension du texte
\item Lexique thématique : Geld, Budget und Projekte
\item Grammaire 1 : comparatif et superlatif
\item Grammaire 2 : conditionnelles (wenn) et Futur I
\item Méthode du commentaire et commentaire modèle
\item Production guidée : Mein Budget, mein Projekt
\item Activité d'intégration
\item Méthodes et exercices résolus
\item Exercices
\item Corrigés des exercices
\item Fiche bilan
\end{enumerate}
\end{multicols}
\end{sommaire}

\begin{objectifs}
\begin{itemize}
\item À la fin de cette leçon, je sais comprendre globalement puis en détail un texte allemand sur la gestion d'un budget et d'un projet.
\item À la fin de cette leçon, je sais employer correctement le vocabulaire thématique (das Budget, die Einnahmen, die Ausgaben, sparen, das Projekt, der Plan, das Geld) avec articles et pluriels.
\item À la fin de cette leçon, je sais former et utiliser le comparatif et le superlatif des adjectifs et adverbes allemands.
\item À la fin de cette leçon, je sais construire une subordonnée conditionnelle avec « wenn » en respectant la place du verbe.
\item À la fin de cette leçon, je sais conjuguer et employer le Futur I pour annoncer un projet ou faire une prévision.
\item À la fin de cette leçon, je sais répondre par écrit à des questions de compréhension en phrases complètes.
\item À la fin de cette leçon, je sais résumer un texte en quelques phrases avec des connecteurs simples.
\item À la fin de cette leçon, je sais rédiger un court paragraphe guidé présentant mon budget ou mon projet.
\end{itemize}
\end{objectifs}

\begin{prerequis}
\begin{itemize}
\item La conjugaison des verbes au présent (verbes réguliers et irréguliers : haben, sein, können, wollen) vue en Première.
\item La place du verbe en position 2 dans la phrase principale et en fin de proposition subordonnée (weil, dass).
\item Le vocabulaire de base de l'argent et des achats (kaufen, kosten, der Preis, der Euro) vu en Première.
\item Les articles définis et indéfinis et le pluriel des noms courants.
\item Les adjectifs qualificatifs courants (gut, viel, teuer, billig, wichtig).
\end{itemize}
\end{prerequis}

\newpage

\coursec{Texte support : « Lernen, mit Geld umzugehen »}

\courssub{Situation d'accroche et problématique}

Aminatou, élève de Terminale à Yaoundé, vient de recevoir 20 000 FCFA de sa mère pour ses fournitures scolaires. Devant les étals du marché Mokolo, une question lui trotte dans la tête : doit-elle tout dépenser tout de suite, ou garder une partie de l'argent pour plus tard ? Son petit frère, lui, rêve d'acheter un ballon de football ; sa grande sœur, elle, met de côté chaque semaine un peu d'argent pour s'offrir un téléphone. Toute la famille parle déjà, sans le savoir, de \cle{budget} et de \cle{projet}.

En Allemagne, en Autriche et en Suisse, les jeunes apprennent très tôt à gérer leur \all{Taschengeld} (argent de poche) et à planifier des projets : un voyage, une fête d'anniversaire, un nouveau vélo. Les magazines pour jeunes publient régulièrement des articles sur ce sujet. C'est exactement le type de texte que tu vas lire aujourd'hui.

Problématique de la leçon : comment un jeune Allemand parle-t-il de son argent, de ses dépenses et de ses projets ? Comment dit-on en allemand « plus cher », « le meilleur plan » ? Que se passera-t-il « si » (\all{wenn}) Aminatou économise chaque mois ? Et comment annoncer ce qu'elle fera plus tard ? Pour répondre, nous allons lire le portrait de Jonas, 17 ans, qui gère son budget comme un chef.

\begin{methode}[Lire un texte allemand en 3 étapes]
\begin{enumerate}
\item \textbf{Lire le titre et la première phrase} pour deviner le thème et le type de texte (article, dialogue, lettre, récit). Ici : \all{Lernen, mit Geld umzugehen} annonce un texte sur la gestion de l'argent.
\item \textbf{Lire le texte en entier sans dictionnaire} et souligner les mots connus ou transparents (\all{Budget}, \all{Projekt}, \all{Plan} ressemblent au français). Comprendre le sens global : qui ? quoi ? où ? pourquoi ?
\item \textbf{Relire paragraphe par paragraphe} avec le glossaire, puis répondre aux questions de compréhension détaillée.
\end{enumerate}
\end{methode}

\courssub{Lecture globale : première approche du texte}

Avant de lire, observe le titre : \all{Lernen, mit Geld umzugehen}. Le verbe \all{umgehen mit} (+ datif) signifie « gérer, se comporter avec » ; c'est un verbe à particule \emph{inséparable} : \all{er geht mit Geld um}. Le titre annonce donc un texte qui apprend aux jeunes à \emph{gérer l'argent}. Le type de texte : un article de magazine jeunesse, avec un portrait (Jonas) et des citations au discours direct. Le locuteur est un journaliste qui raconte la vie d'un adolescent ; Jonas parle lui-même à la première et à la dernière partie.

\begin{attention}[Ne traduis pas mot à mot !]
Le piège classique du francophone : traduire chaque mot dans l'ordre. \all{mit Geld umgehen} ne veut PAS dire « aller avec l'argent » mais « gérer l'argent ». De même, \all{Taschengeld} (littéralement « argent de poche ») se traduit par « argent de poche » en français — ici le mot à mot fonctionne, mais c'est une chance ! Cherche toujours d'abord le \cle{sens global} de la phrase, puis précise avec le glossaire.
\end{attention}

\courssub{Le texte : « Lernen, mit Geld umzugehen »}

\begin{textebox}[Lernen, mit Geld umzugehen — Jugendmagazin « Klartext »]
Jonas ist 17 Jahre alt und wohnt in München. Wie viele Jugendliche bekommt er jeden Monat Taschengeld von seinen Eltern: 80 Euro. Manchmal verdient er mehr Geld, weil er am Wochenende im Garten der Nachbarn hilft oder Zeitungen austrägt. „Meine Einnahmen sind nicht groß, aber sie reichen“, sagt er.

Seine größten Ausgaben sind sein Handy und Kleidung. Auch für die Freizeit — Kino, Sport, Eis mit Freunden — gibt er Geld aus. Früher war er unvorsichtiger mit Geld: Am Ende des Monats war sein Portemonnaie oft leer. Aber jetzt plant er besser. Er schreibt alle Einnahmen und Ausgaben in ein kleines Heft.

Jonas hat ein Ziel: eine Sprachreise nach Österreich. Wenn er jeden Monat 20 Euro spart, wird er im Sommer nach Wien reisen. Das ist sein wichtigstes Projekt. „Wenn ich genug Geld habe, werde ich mit meinen Freunden eine Woche in Wien verbringen“, sagt er stolz. Sein Plan ist einfach: weniger ausgeben, mehr sparen! Und wenn das Geld nicht reicht, wird er in den Ferien arbeiten. So lernt Jonas, mit Geld umzugehen — und sein Traum wird vielleicht wahr.
\end{textebox}

\textbf{Glossaire (mots difficiles) :}

\begin{center}
\begin{tabularx}{\linewidth}{|X|X|}\hline
\rowcolor{popblueL} \textbf{Deutsch} & \textbf{Français} \\ \hline
umgehen mit (+ Dat.) & gérer, se comporter avec \\ \hline
das Taschengeld, -er & l'argent de poche \\ \hline
verdienen & gagner (de l'argent) \\ \hline
die Einnahme, -n & la recette, le revenu \\ \hline
reichen & suffire \\ \hline
die Ausgabe, -n & la dépense \\ \hline
ausgeben (gibt aus, gab aus, hat ausgegeben) & dépenser \\ \hline
unvorsichtig & imprudent \\ \hline
das Portemonnaie, -s & le porte-monnaie \\ \hline
das Heft, -e & le cahier \\ \hline
das Ziel, -e & le but, l'objectif \\ \hline
die Sprachreise, -n & le séjour linguistique \\ \hline
sparen & épargner, économiser \\ \hline
stolz & fier \\ \hline
wahr werden & devenir réalité, se réaliser \\ \hline
\end{tabularx}
\end{center}

\courssub{Lecture détaillée, paragraphe par paragraphe}

\textbf{Paragraphe 1 : les revenus de Jonas (die Einnahmen).}\par
Jonas a deux sources de revenus : son \all{Taschengeld} (80 euros par mois, versé par ses parents) et de petits jobs : il aide dans le jardin des voisins et distribue des journaux (\all{Zeitungen austragen}). Phrase clé : \all{Meine Einnahmen sind nicht groß, aber sie reichen.} « Mes revenus ne sont pas élevés, mais ils suffisent. »

\textbf{Paragraphe 2 : ses dépenses (die Ausgaben).}\par
Ses plus grandes dépenses : le téléphone portable (\all{das Handy}) et les vêtements (\all{die Kleidung}), plus les loisirs. Avant (\all{früher}), il était imprudent ; maintenant (\all{jetzt}), il note tout dans un cahier. C'est la naissance de son budget.

\textbf{Paragraphe 3 : son projet et son plan d'épargne (das Projekt, der Plan).}\par
Son but : un séjour linguistique en Autriche. Sa stratégie : épargner 20 euros par mois ; si l'argent ne suffit pas, il travaillera pendant les vacances. Sa devise : \all{weniger ausgeben, mehr sparen!} « dépenser moins, épargner plus ! »

\begin{exemplebox}[Phrases clés du texte, traduites]
\all{Er verdient mehr Geld.} « Il gagne plus d'argent. »

\all{Wenn er spart, wird er reisen.} « S'il économise, il voyagera. »

\all{Sein wichtigstes Projekt} « son projet le plus important »

\all{Früher war er unvorsichtiger mit Geld.} « Avant, il était plus imprudent avec l'argent. »

\all{Wenn das Geld nicht reicht, wird er in den Ferien arbeiten.} « Si l'argent ne suffit pas, il travaillera pendant les vacances. »
\end{exemplebox}

\courssub{Relevé des structures grammaticales de la leçon}

Le texte contient exactement les structures que nous étudierons dans les sections 4 et 5. Repérons-les dès maintenant :

\begin{center}
\begin{tabularx}{\linewidth}{|X|X|X|}\hline
\rowcolor{popblueL} \textbf{Structure} & \textbf{Exemple du texte} & \textbf{Traduction} \\ \hline
Comparatif & \all{er verdient mehr Geld} & il gagne plus d'argent \\ \hline
Comparatif & \all{früher war er unvorsichtiger} & avant, il était plus imprudent \\ \hline
Superlatif & \all{sein wichtigstes Projekt} & son projet le plus important \\ \hline
Conditionnelle (wenn) & \all{Wenn er jeden Monat 20 Euro spart, ...} & S'il épargne 20 euros par mois, ... \\ \hline
Conditionnelle (wenn) & \all{Wenn das Geld nicht reicht, ...} & Si l'argent ne suffit pas, ... \\ \hline
Futur I & \all{wird er im Sommer nach Wien reisen} & il voyagera à Vienne en été \\ \hline
Futur I & \all{werde ich eine Woche in Wien verbringen} & je passerai une semaine à Vienne \\ \hline
\end{tabularx}
\end{center}

\begin{aretenir}[Ce qu'il faut observer dès maintenant]
\begin{itemize}
\item Dans la subordonnée avec \all{wenn}, le verbe conjugué va \textbf{à la fin} : \all{Wenn er jeden Monat 20 Euro \textbf{spart}, ...}
\item Au Futur I, on utilise \all{werden} + infinitif \textbf{à la fin} : \all{er wird nach Wien \textbf{reisen}}.
\item Le comparatif ajoute \textbf{-er} à l'adjectif (\all{unvorsichtig $\to$ unvorsichtiger}), le superlatif \textbf{-st-} (\all{wichtig $\to$ am wichtigsten}).
\end{itemize}
\end{aretenir}

\courssub{Carte mentale : le budget de Jonas}

\begin{popfigure}
\begin{center}
\begin{tikzpicture}[
  mindmap,
  grow cyclic,
  every node/.style={concept, align=center, text=white},
  root concept/.append style={concept color=popblue, font=\bfseries},
  level 1/.append style={level distance=3.4cm, sibling angle=120},
  level 2/.append style={level distance=2.2cm, sibling angle=50, font=\small},
]
\node[root concept, align=center]{Jonas'\\Budget}
  child[concept color=popgreen]{ node[align=center]{Einnahmen\\(recettes)}
    child { node[align=center]{Taschengeld\\80 Euro} }
    child { node[align=center]{Jobs:\\Garten,\\Zeitungen} }
  }
  child[concept color=poporange]{ node[align=center]{Ausgaben\\(dépenses)}
    child { node {Handy} }
    child { node {Kleidung} }
    child { node {Freizeit} }
  }
  child[concept color=poppurple]{ node[align=center]{Projekt\\und Plan}
    child { node[align=center]{Sprachreise\\Wien} }
    child { node[align=center]{20 Euro\\pro Monat\\sparen} }
  };
\end{tikzpicture}
\end{center}
\legende{Carte mentale : les trois parties du budget de Jonas — recettes, dépenses, projet.}
\end{popfigure}

\begin{popfigure}
\begin{center}
\begin{tikzpicture}[node distance=0.55cm and 0.5cm,
  box/.style={draw=popblue, fill=popblueL, rounded corners, align=center, minimum height=1cm, inner sep=6pt, font=\small},
  arr/.style={-{Stealth[length=3mm]}, thick, popdark}]
\node[box, align=center] (a){Geld kommt\\(Einnahmen)};
\node[box, right=of a, align=center] (b){Budget wird\\geplant};
\node[box, right=of b, align=center] (c){Ausgaben werden\\kontrolliert};
\node[box, right=of c, align=center] (d){Geld wird\\gespart};
\node[box, right=of d, align=center] (e){Projekt wird\\verwirklicht};
\draw[arr] (a) -- (b);
\draw[arr] (b) -- (c);
\draw[arr] (c) -- (d);
\draw[arr] (d) -- (e);
\end{tikzpicture}
\end{center}
\legende{Le cycle du budget : de la recette au projet réalisé.}
\end{popfigure}

\courssub{Questions de compréhension}

\begin{exercice}[Verstehen des Textes]
Réponds en allemand, par des phrases complètes :
\begin{enumerate}
\item Wie viel Taschengeld bekommt Jonas jeden Monat?
\item Wie verdient er manchmal mehr Geld?
\item Was sind seine größten Ausgaben?
\item Was macht er jetzt mit seinen Einnahmen und Ausgaben?
\item Was ist sein Ziel?
\item Was wird er machen, wenn das Geld nicht reicht?
\end{enumerate}
\end{exercice}

\begin{corrige}[Verstehen des Textes]
\begin{enumerate}
\item \all{Er bekommt jeden Monat 80 Euro Taschengeld von seinen Eltern.} « Il reçoit 80 euros d'argent de poche par mois de ses parents. »
\item \all{Er verdient mehr Geld, weil er am Wochenende im Garten der Nachbarn hilft oder Zeitungen austrägt.} « Il gagne plus d'argent parce qu'il aide dans le jardin des voisins ou distribue des journaux le week-end. »
\item \all{Seine größten Ausgaben sind sein Handy und Kleidung.} « Ses plus grandes dépenses sont son portable et les vêtements. »
\item \all{Er schreibt alle Einnahmen und Ausgaben in ein kleines Heft.} « Il note toutes les recettes et les dépenses dans un petit cahier. »
\item \all{Sein Ziel ist eine Sprachreise nach Österreich.} « Son but est un séjour linguistique en Autriche. »
\item \all{Wenn das Geld nicht reicht, wird er in den Ferien arbeiten.} « Si l'argent ne suffit pas, il travaillera pendant les vacances. »
\end{enumerate}
\end{corrige}

\begin{exoresolu}[Compréhension détaillée : question modèle]
Énoncé : Was ist Jonas' wichtigstes Projekt, und wie will er es verwirklichen? (Réponds en deux phrases complètes.)
\tcblower
\textbf{Solution détaillée :}
\begin{enumerate}
\item On repère dans le texte le mot \all{wichtigstes Projekt} : il se trouve au paragraphe 3.
\item On recopie la structure de la phrase en changeant le point de vue (on répond à la troisième personne).
\end{enumerate}
Réponse : \all{Sein wichtigstes Projekt ist eine Sprachreise nach Österreich. Er will jeden Monat 20 Euro sparen, und wenn das Geld nicht reicht, wird er in den Ferien arbeiten.}

Traduction : « Son projet le plus important est un séjour linguistique en Autriche. Il veut épargner 20 euros par mois, et si l'argent ne suffit pas, il travaillera pendant les vacances. »
\end{exoresolu}

\courssub{Reformulation orale : die Zusammenfassung}

Après la lecture détaillée, tu dois être capable de raconter le texte \emph{avec tes propres mots}, sans regarder l'article. Voici un modèle de résumé oral (environ 5 phrases) :

\all{Jonas ist 17 und bekommt 80 Euro Taschengeld pro Monat. Er verdient auch Geld mit kleinen Jobs. Sein Geld gibt er für Handy, Kleidung und Freizeit aus. Jetzt plant er besser und spart jeden Monat 20 Euro. Sein Traum ist eine Reise nach Wien mit seinen Freunden.}

« Jonas a 17 ans et reçoit 80 euros d'argent de poche par mois. Il gagne aussi de l'argent avec de petits jobs. Il dépense son argent pour le portable, les vêtements et les loisirs. Maintenant, il planifie mieux et épargne 20 euros par mois. Son rêve est un voyage à Vienne avec ses amis. »

\begin{experience}[À toi de parler !]
Par groupes de deux : l'élève A résume le texte avec ses propres mots (sans le regarder), l'élève B vérifie avec le texte et corrige. Puis on échange les rôles. Utilise obligatoirement trois mots du glossaire : \all{Einnahmen}, \all{Ausgaben}, \all{sparen}.
\end{experience}

\begin{savaistu}[Le Taschengeld en Allemagne]
En Allemagne, beaucoup d'enfants et d'adolescents reçoivent un \all{Taschengeld} régulier, souvent versé chaque semaine ou chaque mois. Il existe même une \all{Taschengeldtabelle} officielle, publiée par les services de la jeunesse, qui conseille aux parents le montant adapté à l'âge : par exemple quelques euros par semaine pour un enfant de 10 ans, et davantage par mois pour un jeune de 16 ou 17 ans. L'objectif est pédagogique : apprendre très tôt à gérer un budget, à choisir ses dépenses et à épargner pour un projet — exactement comme Jonas dans notre texte !
\end{savaistu}

\begin{aretenir}[À retenir de cette section]
\begin{itemize}
\item Le texte est un \textbf{article de magazine jeunesse} : portrait d'un adolescent + citations au discours direct.
\item Trois idées = trois paragraphes : \all{Einnahmen} (recettes), \all{Ausgaben} (dépenses), \all{Projekt/Plan} (projet et plan d'épargne).
\item Stratégie de lecture : titre $\to$ sens global sans dictionnaire $\to$ lecture détaillée avec glossaire.
\item Les structures grammaticales de la leçon (comparatif, superlatif, \all{wenn}, Futur I) sont déjà repérées dans le texte : nous allons maintenant les étudier en détail.
\end{itemize}
\end{aretenir}

\coursec{Compréhension du texte}

Après avoir lu et écouté le texte support \all{Lernen, mit Geld umzugehen} (section 1), nous passons maintenant à son exploitation méthodique. Comprendre un texte allemand au baccalauréat, ce n'est pas le traduire mot à mot : c'est savoir répondre à des questions précises, en allemand, en phrases complètes, en s'appuyant sur des citations du texte. Cette section vous donne la méthode complète, puis l'applique pas à pas au texte de Jonas.

\courssub{Rappel du texte support}

Pour travailler efficacement, gardons sous les yeux un extrait du texte lu dans la section précédente.

\begin{textebox}[Lernen, mit Geld umzugehen (Auszug)]
Jonas ist siebzehn Jahre alt und besucht das Gymnasium in einer kleinen Stadt in Bayern. Jeden Monat bekommt er 80 Euro Taschengeld von seinen Eltern. Früher gab er sein Geld sofort aus: für Süßigkeiten, für Handy-Spiele und für Kinobesuche mit Freunden. Am Ende des Monats war sein Portemonnaie immer leer. „Ich weiß nicht, wohin mein Geld verschwindet“, sagte er oft.

Seit drei Monaten führt Jonas ein Haushaltsbuch. Dort schreibt er alle Einnahmen und Ausgaben auf. Seine größten Ausgaben sind das Mittagessen in der Mensa (etwa 25 Euro im Monat) und sein Handy-Vertrag (15 Euro). Jetzt plant er ein Projekt: er möchte im Sommer mit dem Zug nach Österreich fahren und dort eine Woche wandern. Die Reise kostet ungefähr 240 Euro. Wenn er jeden Monat 20 Euro spart, wird er in einem Jahr genug Geld für die Reise haben. „Sparen ist leichter, als ich gedacht habe“, sagt Jonas heute. „Man muss nur einen Plan haben.“
\end{textebox}

\textbf{Glossaire de l'extrait}

\begin{center}
\begin{tabularx}{\linewidth}{|l|X|}
\hline
\rowcolor{popblueL} \textbf{Deutsch} & \textbf{Français} \\
\hline
das Taschengeld, -er & l'argent de poche \\
\hline
ausgeben (gibt aus, gab aus, hat ausgegeben) & dépenser (de l'argent) \\
\hline
die Süßigkeit, -en & la friandise, la douceur \\
\hline
das Portemonnaie, -s & le porte-monnaie \\
\hline
das Haushaltsbuch, -bücher & le carnet de comptes \\
\hline
die Einnahme, -n & la recette, le revenu \\
\hline
die Ausgabe, -n & la dépense \\
\hline
die Mensa, Mensen & la cantine (scolaire ou universitaire) \\
\hline
der Handy-Vertrag, -verträge & le forfait de téléphone portable \\
\hline
sparen & économiser, épargner \\
\hline
der Plan, -pläne & le plan, le projet \\
\hline
ungefähr & environ, à peu près \\
\hline
\end{tabularx}
\end{center}

\courssub{La méthode de réponse à une question de compréhension}

Au baccalauréat, les questions de compréhension (\all{Fragen zum Text}) sont notées sur la pertinence de la réponse ET sur la qualité de la langue. Une bonne réponse suit toujours quatre étapes.

\begin{methode}[Répondre à une question de compréhension au Bac]
\begin{enumerate}
\item \textbf{Souligner le mot interrogatif} : \all{Wer?} (qui ?), \all{Was?} (quoi ?), \all{Wie viel?} / \all{Wie viele?} (combien ?), \all{Warum?} (pourquoi ?), \all{Wann?} (quand ?), \all{Wohin?} (où, direction ?). Le mot interrogatif détermine la nature de la réponse attendue.
\item \textbf{Localiser le passage} du texte qui contient l'information : repérer dans le texte les mots-clés de la question (ou leurs synonymes).
\item \textbf{Rédiger une phrase complète en allemand}, en reprenant les mots de la question, sans recopier tout le paragraphe. Une question avec \all{Wie viel} exige un chiffre dans la réponse ; une question avec \all{Warum} exige une cause (souvent avec \all{weil} ou \all{denn}).
\item \textbf{Citer le texte} (le \all{Beleg}) entre guillemets allemands „…“ pour justifier la réponse, surtout dans les questions vrai/faux.
\end{enumerate}
\end{methode}

\begin{popfigure}
\begin{center}
\begin{tikzpicture}[node distance=0.55cm,
etape/.style={rectangle, rounded corners, draw=popblue, fill=popblueL, thick, align=center, minimum width=4.6cm, minimum height=0.9cm, font=\small},
fleche/.style={-{Stealth[length=3mm]}, thick, draw=poporange}]
\node[etape, align=center] (e1){1. \textbf{Frage lesen}\\ lire la question};
\node[etape, below=of e1, align=center] (e2){2. \textbf{Stelle im Text finden}\\ localiser le passage};
\node[etape, below=of e2, align=center] (e3){3. \textbf{Satz bilden}\\ rédiger une phrase complète};
\node[etape, below=of e3, align=center] (e4){4. \textbf{Beleg zitieren}\\ citer le texte};
\draw[fleche] (e1) -- (e2);
\draw[fleche] (e2) -- (e3);
\draw[fleche] (e3) -- (e4);
\end{tikzpicture}
\end{center}
\legende{Organigramme de la méthode de réponse à une question de compréhension}
\end{popfigure}

\begin{exemplebox}[Question et réponse modèles]
\textbf{Frage :} \all{Wie viel Taschengeld bekommt Jonas?}

\textbf{Antwort :} \all{Jonas bekommt jeden Monat 80 Euro Taschengeld.}

Traduction : « Jonas reçoit 80 euros d'argent de poche par mois. »

Observez : la réponse reprend les mots de la question (\all{bekommen}, \all{Taschengeld}), donne le chiffre demandé et forme une phrase complète avec sujet, verbe en position 2 et compléments.
\end{exemplebox}

\begin{propriete}[La place du verbe dans la réponse]
Dans une réponse, le verbe conjugué reste toujours en \cle{position 2}. Si la réponse commence par un complément (et non par le sujet), le sujet passe après le verbe : c'est l'inversion.

\all{Jonas bekommt jeden Monat 80 Euro.} (sujet -- verbe)

\all{Jeden Monat bekommt Jonas 80 Euro.} (complément -- verbe -- sujet)

« Chaque mois, Jonas reçoit 80 euros. »
\end{propriete}

\begin{attention}[\all{Wie viel} ou \all{Wie viele} ?]
\all{Wie viel} s'emploie devant un nom indénombrable (qu'on ne peut pas compter) : \all{Wie viel Geld?} (combien d'argent ?), \all{Wie viel Zeit?} (combien de temps ?).

\all{Wie viele} s'emploie devant un nom dénombrable au pluriel : \all{Wie viele Euro?} (combien d'euros ?), \all{Wie viele Freunde?} (combien d'amis ?).

Piège du francophone : en français, « combien de » convient partout ; en allemand, il faut choisir. On dit donc \all{Wie viel Geld bekommt Jonas?} mais \all{Wie viele Euro bekommt Jonas?} — les deux sont correctes, mais jamais \all{wie viel Euro}.
\end{attention}

\courssub{Compréhension globale du texte}

Les questions globales portent sur le sens général du texte. Voici les questions types et leurs réponses modèles.

\begin{enumerate}
\item \all{Worum geht es in dem Text?} (De quoi parle le texte ?)

Réponse modèle : \all{Der Text handelt von einem Jugendlichen, der lernt, mit seinem Geld umzugehen.} « Le texte parle d'un adolescent qui apprend à gérer son argent. »

\item \all{Wer ist Jonas?} (Qui est Jonas ?)

Réponse modèle : \all{Jonas ist ein siebzehnjähriger Schüler. Er besucht das Gymnasium in einer kleinen Stadt in Bayern.} « Jonas est un élève de dix-sept ans. Il fréquente le lycée dans une petite ville de Bavière. »

\item \all{Welches Problem hat Jonas, und welche Lösung findet er?} (Quel est son problème et quelle solution trouve-t-il ?)

Réponse modèle : \all{Sein Problem ist, dass er sein Geld sofort ausgibt und am Ende des Monats nichts mehr hat. Die Lösung ist ein Haushaltsbuch: er schreibt alle Einnahmen und Ausgaben auf.} « Son problème est qu'il dépense son argent immédiatement et n'a plus rien à la fin du mois. La solution est un carnet de comptes : il note toutes ses recettes et ses dépenses. »
\end{enumerate}

\courssub{Compréhension détaillée du texte}

Les questions détaillées exigent des informations précises : chiffres, noms, montants.

\begin{enumerate}
\item \all{Wie viel Taschengeld bekommt Jonas im Monat?}

\all{Er bekommt 80 Euro Taschengeld im Monat.} « Il reçoit 80 euros d'argent de poche par mois. »

\item \all{Was sind seine größten Ausgaben?}

\all{Seine größten Ausgaben sind das Mittagessen in der Mensa (etwa 25 Euro) und sein Handy-Vertrag (15 Euro).} « Ses plus grandes dépenses sont le déjeuner à la cantine (environ 25 euros) et son forfait de téléphone (15 euros). »

\item \all{Was wird Jonas machen, wenn er jeden Monat 20 Euro spart?}

\all{Wenn er jeden Monat 20 Euro spart, wird er in einem Jahr genug Geld für die Reise nach Österreich haben.} « S'il économise 20 euros par mois, il aura dans un an assez d'argent pour le voyage en Autriche. »

Observez la structure : la subordonnée avec \all{wenn} renvoie le verbe conjugué à la fin (\all{spart}), et la principale commence par le verbe (\all{wird er}). Nous étudierons cette construction en détail dans la section grammaire.
\end{enumerate}

\courssub{Vrai ou faux : justifier par une citation}

L'exercice \all{Richtig oder falsch?} est classique au baccalauréat. La règle d'or : une réponse sans citation (\all{Beleg}) ne rapporte que la moitié des points.

\begin{exoresolu}[Richtig oder falsch ? Begründen Sie mit einem Beleg]
Énoncé. Dites si les phrases suivantes sont vraies ou fausses et justifiez par une citation du texte.

a) \all{Jonas bekommt 100 Euro Taschengeld im Monat.}

b) \all{Früher gab Jonas sein Geld sofort aus.}

c) \all{Die Reise nach Österreich kostet 400 Euro.}
\tcblower
Solution détaillée.

a) \textbf{Falsch.} Beleg : \all{„Jeden Monat bekommt er 80 Euro Taschengeld von seinen Eltern.“} Jonas reçoit 80 euros, et non 100 euros.

b) \textbf{Richtig.} Beleg : \all{„Früher gab er sein Geld sofort aus.“} La phrase du texte confirme l'affirmation mot pour mot.

c) \textbf{Falsch.} Beleg : \all{„Die Reise kostet ungefähr 240 Euro.“} Le voyage coûte environ 240 euros, et non 400 euros.

Méthode : on écrit d'abord \all{Richtig} ou \all{Falsch}, puis le mot \all{Beleg} suivi de la citation exacte entre guillemets allemands. Si la phrase est fausse, on ajoute la correction en une phrase.
\end{exoresolu}

\courssub{Vocabulaire en contexte}

Les questions de vocabulaire demandent de retrouver dans le texte un mot synonyme ou contraire d'un mot donné. Il faut chercher dans le texte, pas dans son dictionnaire mental.

\begin{enumerate}
\item Trouvez dans le texte le contraire de \all{sparen}.

Réponse : \all{ausgeben}. Belege : \all{„Früher gab er sein Geld sofort aus“} s'oppose à \all{„Wenn er jeden Monat 20 Euro spart …“}. Le texte construit l'opposition entre le Jonas d'avant (il dépense) et le Jonas d'aujourd'hui (il économise).

\item Trouvez dans le texte une expression qui signifie « gérer son argent ».

Réponse : \all{mit Geld umgehen}. Beleg : le titre même du texte, \all{„Lernen, mit Geld umzugehen“}. C'est un verbe à particule séparable : \all{Jonas geht mit seinem Geld um.}

\item Trouvez dans le texte un mot de la même famille que \all{die Ausgabe}.

Réponse : le verbe \all{ausgeben} (dépenser). La famille de mots est transparente : \all{ausgeben} (verbe) donne \all{die Ausgabe} (nom féminin, la dépense) ; de même \all{einnehmen} (encaisser) donne \all{die Einnahme} (la recette).
\end{enumerate}

\begin{center}
\begin{tabularx}{\linewidth}{|X|X|X|}
\hline
\rowcolor{popblueL} \textbf{Français} & \textbf{Deutsch} & \textbf{Remarque} \\
\hline
dépenser de l'argent & \all{Geld ausgeben} & verbe à particule séparable : \all{er gibt aus} \\
\hline
économiser, épargner & \all{sparen} & contraire de \all{ausgeben} \\
\hline
gérer son argent & \all{mit Geld umgehen} & verbe à particule : \all{er geht damit um} \\
\hline
noter, inscrire & \all{aufschreiben} & séparable : \all{er schreibt auf} \\
\hline
recevoir de l'argent de poche & \all{Taschengeld bekommen} & \all{bekommen} = recevoir, pas « devenir » \\
\hline
planifier un projet & \all{ein Projekt planen} & \all{das Projekt, -e} ; \all{planen} est régulier \\
\hline
\end{tabularx}
\end{center}

\begin{attention}[Faux ami : \all{bekommen}]
\all{bekommen} signifie « recevoir, obtenir », jamais « devenir ». « Devenir » se dit \all{werden}. \all{Jonas bekommt 80 Euro.} = « Jonas reçoit 80 euros » ; \all{Jonas wird reich.} = « Jonas devient riche ». Autre faux ami du texte : \all{das Gymnasium} désigne le lycée (voie générale), pas le « gymnase ».
\end{attention}

\courssub{Question personnelle : donner son opinion}

La dernière question invite à parler de soi : elle prépare la production guidée de la fin de la leçon. On répond à la première personne, en phrases simples, avec le vocabulaire du texte.

\textbf{Frage :} \all{Und du? Wie gehst du mit deinem Geld um?} « Et toi, comment gères-tu ton argent ? »

Réponse modèle : \all{Ich bekomme jede Woche etwas Geld von meinen Eltern. Ich gebe es für Transport und Essen aus, aber ich spare auch einen Teil. Wenn ich genug Geld habe, werde ich mir ein neues Handy kaufen.} « Je reçois chaque semaine un peu d'argent de mes parents. Je le dépense pour le transport et la nourriture, mais j'économise aussi une partie. Quand j'aurai assez d'argent, je m'achèterai un nouveau téléphone portable. »

\begin{experience}[À vous de répondre]
Répondez par écrit aux cinq questions suivantes, en phrases complètes, puis comparez vos réponses avec celles de votre voisin.

1. \all{Wie alt ist Jonas?}

2. \all{Wofür gab Jonas früher sein Geld aus?}

3. \all{Was macht Jonas seit drei Monaten?}

4. Richtig oder falsch (mit Beleg) : \all{Jonas möchte im Sommer nach Österreich fahren.}

5. \all{Und du? Sparst du Geld? Wofür?}
\end{experience}

\begin{corrige}[Exercice : questions de compréhension]
1. \all{Jonas ist siebzehn Jahre alt.} Beleg : \all{„Jonas ist siebzehn Jahre alt.“}

2. \all{Früher gab er sein Geld für Süßigkeiten, für Handy-Spiele und für Kinobesuche mit Freunden aus.} Attention : le verbe séparable \all{ausgeben} place sa particule \all{aus} à la fin de la phrase.

3. \all{Seit drei Monaten führt er ein Haushaltsbuch. Dort schreibt er alle Einnahmen und Ausgaben auf.} Inversion après \all{dort} : le sujet \all{er} passe après le verbe.

4. \textbf{Richtig.} Beleg : \all{„er möchte im Sommer mit dem Zug nach Österreich fahren.“}

5. Réponse libre, par exemple : \all{Ja, ich spare Geld. Ich möchte mir ein Fahrrad kaufen.} / \all{Nein, ich gebe mein Geld sofort aus, wie Jonas früher.}
\end{corrige}

\begin{aretenir}[L'essentiel de la méthode]
Pour toute question de compréhension : souligner le mot interrogatif, localiser le passage, répondre en une phrase complète avec le verbe en position 2, et citer le texte entre guillemets „…“ pour justifier. Distinguer \all{wie viel} (indénombrable) et \all{wie viele} (dénombrable). Ces réflexes, acquis sur le texte de Jonas, serviront pour tous les textes de l'année et à l'épreuve du baccalauréat.
\end{aretenir}

\coursec{Lexique thématique : Geld, Budget und Projekte}

Après avoir lu et compris le texte \all{Lernen, mit Geld umzugehen}, nous avons rencontré de nombreux mots nouveaux : \all{das Budget}, \all{die Ausgaben}, \all{sparen}, \all{das Projekt}... Pour bien parler de la gestion de l'argent et des projets — à l'écrit comme à l'oral — il faut maintenant organiser ce vocabulaire, mémoriser les articles et les pluriels, et apprendre les constructions correctes. C'est l'objectif de cette section : faire de vous de véritables experts du lexique économique allemand.

\courssub{Observation : le vocabulaire en contexte}

Lisons d'abord un court texte original qui rassemble presque tout le vocabulaire du chapitre. Observez les mots en gras : nous les étudierons ensuite champ par champ.

\begin{textebox}[Ahmed plant sein erstes Projekt]
Ahmed ist 17 Jahre alt und besucht ein Gymnasium in Yaoundé. Am Ende des Schuljahres möchte er mit zwei Freunden ein kleines \textbf{Projekt} verwirklichen: Sie wollen auf dem Schulhof Snacks und Fruchtsaft verkaufen. Aber zuerst brauchen sie einen guten \textbf{Plan} und ein realistisches \textbf{Budget}.

Ahmed bekommt jeden Monat ein \textbf{Taschengeld} von 10 000 Francs. Normalerweise \textbf{gibt} er fast alles für Kleidung und Handykarten \textbf{aus}. Jetzt will er \textbf{sparen}: Er zahlt jeden Monat 4 000 Francs auf sein \textbf{Sparkonto} \textbf{ein}. „Wenn ich drei Monate lang \textbf{sparsam} lebe, habe ich genug \textbf{Geld} für unser Projekt“, sagt er. Seine Freunde machen auch einen \textbf{Sparplan}.

Zusammen schreiben die drei Jugendlichen alle \textbf{Kosten} auf: Früchte, Zucker, Becher, Servietten. Der gesamte \textbf{Betrag} ist nicht sehr hoch. Sie vergleichen ihre \textbf{Einnahmen} und \textbf{Ausgaben} in einer kleinen Tabelle. „Wir dürfen keine \textbf{Schulden machen}“, warnt Ahmed. „Und wir dürfen unser Geld nicht \textbf{verschwenden}!“

Die drei \textbf{organisieren} alles gut: Sie \textbf{planen} die Einkäufe auf dem Markt, \textbf{bereiten} den Verkauf \textbf{vor} und legen einen \textbf{Termin} fest. Ahmeds Vater, der als Lehrer ein bescheidenes \textbf{Gehalt} \textbf{verdient}, ist stolz auf seinen Sohn: „Du lernst, mit Geld umzugehen. Das ist wichtiger als viel Geld.“

Nach zwei Wochen hat das Projekt Erfolg: Die \textbf{Einnahmen} sind höher als die \textbf{Ausgaben}. Ahmed hat sein \textbf{Ziel erreicht} — und er hat gelernt, dass ein guter Plan und ein klares Budget die besten Freunde eines jeden Projekts sind.
\end{textebox}

\textbf{Glossaire :} \all{das Taschengeld} = l'argent de poche ; \all{die Handykarte (-n)} = la recharge téléphonique ; \all{der Becher (-)} = le gobelet ; \all{die Serviette (-n)} = la serviette ; \all{der Einkauf (⸚e)} = l'achat, les courses ; \all{einzahlen (zahlt ein, hat eingezahlt)} = verser (sur un compte) ; \all{festlegen} = fixer ; \all{bescheiden} = modeste ; \all{stolz} = fier.

\textbf{Questions de compréhension rapide :}
\begin{enumerate}
\item \all{Was will Ahmed mit seinen Freunden machen?} (Que veut faire Ahmed avec ses amis ?)
\item \all{Wie viel Geld zahlt er jeden Monat auf sein Sparkonto ein?} (Combien verse-t-il chaque mois sur son compte d'épargne ?)
\item \all{Wovor warnt Ahmed seine Freunde?} (Contre quoi Ahmed met-il ses amis en garde ?)
\item \all{Warum ist Ahmeds Vater stolz?} (Pourquoi son père est-il fier ?)
\end{enumerate}

\courssub{Champ 1 : L'argent et le budget}

\begin{definition}[das Budget]
\all{Das Budget (-s)} désigne l'ensemble des sommes prévues pour les dépenses et les recettes d'une personne, d'une famille, d'une entreprise ou d'un projet. Pour le budget familial, l'allemand emploie aussi \all{der Haushaltsplan (-⸚e)} (littéralement « plan du ménage »). Exemple : \all{Wir machen am Anfang des Monats ein Budget.} « Nous établissons un budget au début du mois. »
\end{definition}

Voici les noms essentiels de ce champ. Mémorisez-les toujours \textbf{avec leur article et leur pluriel} : en allemand, le genre est imprévisible et le pluriel varie.

\begin{center}
\begin{tabularx}{\linewidth}{|l|l|X|}
\hline
\rowcolor{popblueL} Deutsch & Plural & Français \\
\hline
das Geld & die Gelder & l'argent (somme d'argent) \\
\hline
das Budget & die Budgets & le budget \\
\hline
der Euro & die Euros & l'euro \\
\hline
das Taschengeld & die Taschengelder & l'argent de poche \\
\hline
der Betrag & die Beträge & le montant, la somme \\
\hline
die Kosten & (Plur.) & les frais, les coûts \\
\hline
der Haushaltsplan & die Haushaltspläne & le budget familial \\
\hline
\end{tabularx}
\end{center}

\begin{attention}[die Kosten : presque toujours au pluriel]
\all{Die Kosten} s'emploie presque exclusivement au pluriel, comme « les frais » en français : \all{Die Kosten sind zu hoch.} « Les frais sont trop élevés. » Attention aussi au faux ami : \all{das Geld} signifie « l'argent » au sens de monnaie, jamais « l'argent » le métal (qui se dit \all{das Silber}).
\end{attention}

\courssub{Champ 2 : Revenus et dépenses}

\begin{definition}[die Einnahme / die Ausgabe]
\all{Die Einnahme (-n)} = la recette, le revenu (ce qui entre). \all{Die Ausgabe (-n)} = la dépense (ce qui sort). Ces deux mots se construisent sur les verbes \all{einnehmen} (encaisser) et \all{ausgeben} (dépenser).
\end{definition}

\begin{attention}[Einnahmen und Ausgaben : au pluriel dans les textes économiques]
Dans les textes sur le budget, ces deux noms sont presque toujours au \textbf{pluriel} : \all{die Einnahmen} (les recettes) et \all{die Ausgaben} (les dépenses). Phrase modèle à retenir : \all{Die Einnahmen müssen höher sein als die Ausgaben.} « Les recettes doivent être plus élevées que les dépenses. »
\end{attention}

Les verbes de ce champ :

\begin{center}
\begin{tabularx}{\linewidth}{|l|l|X|}
\hline
\rowcolor{popblueL} Verb & Konjugation & Français \\
\hline
verdienen & verdient, hat verdient & gagner (de l'argent) \\
\hline
ausgeben & gibt aus, gab aus, hat ausgegeben & dépenser \\
\hline
kosten & kostet, hat gekostet & coûter \\
\hline
bezahlen & bezahlt, hat bezahlt & payer \\
\hline
\end{tabularx}
\end{center}

\begin{exemplebox}[Exemples traduits]
\all{Er gibt viel Geld für Kleidung aus.} « Il dépense beaucoup d'argent pour les vêtements. »

\all{Mein Vater verdient ein gutes Gehalt.} « Mon père gagne un bon salaire. » (\all{das Gehalt, -⸚er} = le salaire)

\all{Das Handy kostet 80 000 Francs.} « Le téléphone coûte 80 000 francs. »

\all{Sie bezahlt die Rechnung sofort.} « Elle paie la facture tout de suite. »
\end{exemplebox}

\begin{attention}[ausgeben : verbe à particule séparable et verbe fort]
\all{Ausgeben} cumule deux difficultés : c'est un verbe \textbf{fort} (\all{gibt aus, gab aus, hat ausgegeben}) ET un verbe à \textbf{particule séparable}. À l'indicatif présent, la particule va en fin de phrase : \all{Ich gebe jeden Monat 5 000 Francs aus.} Ne dites jamais \all{*Ich ausgebe...}. Au Perfekt : \all{hat ausgegeben} (le \all{ge-} s'insère entre la particule et le radical). Même logique pour \all{einzahlen} : \all{Er zahlt 4 000 Francs ein.} / \all{Er hat 4 000 Francs eingezahlt.}
\end{attention}

\courssub{Champ 3 : Épargne et gestion}

\begin{definition}[sparen]
\all{Sparen (spart, hat gespart)} = économiser, épargner. Vocabulaire associé : \all{das Sparkonto (-en)} = le compte d'épargne ; \all{der Sparplan (-⸚e)} = le plan d'épargne ; \all{sparsam} (adj.) = économe ; \all{verschwenden (hat verschwendet)} = gaspiller ; \all{Schulden machen} = s'endetter, faire des dettes (\all{die Schuld (-en)} = la dette).
\end{definition}

\begin{attention}[sparen + für / auf, jamais « zu »]
Le verbe \all{sparen} se construit avec la préposition \all{für} ou \all{auf} (+ accusatif) : \all{Ich spare für eine Reise.} « J'économise pour un voyage. » \all{Er spart auf ein neues Handy.} « Il économise en vue d'un nouveau téléphone. » Ne dites jamais \all{*sparen zu} : la préposition française « pour » ne se traduit pas ici par \all{zu}.
\end{attention}

\begin{exemplebox}[Exemples traduits]
\all{Wir sparen für ein neues Motorrad.} « Nous économisons pour une nouvelle moto. »

\all{Sie ist sehr sparsam und verschwendet nie Geld.} « Elle est très économe et ne gaspille jamais d'argent. »

\all{Wer Schulden macht, hat später Probleme.} « Qui s'endette aura des problèmes plus tard. »
\end{exemplebox}

\begin{savaistu}[Die Sparkasse und das Sparschwein]
L'épargne est une valeur très ancrée dans les pays germanophones. Les \all{Sparkassen} (caisses d'épargne) existent en Allemagne depuis le début du XIX\up{e} siècle et distribuent encore aujourd'hui aux écoliers la fameuse tirelire en forme de cochon, \all{das Sparschwein} (litt. « le cochon d'épargne »). On y dit volontiers : \all{Spare in der Zeit, so hast du in der Not} — « épargne quand tout va bien, tu auras de quoi faire face dans le besoin ». Une sagesse que comprend aussi tout commerçant du marché Mokolo ou de Deïdo !
\end{savaistu}

\courssub{Champ 4 : Projet et planification}

\begin{definition}[das Projekt und der Plan]
\all{Das Projekt (-e)} = le projet ; \all{der Plan (-⸚e)} = le plan ; \all{das Ziel (-e)} = le but, l'objectif ; \all{der Termin (-e)} = le rendez-vous, la date butoir. Verbes associés : \all{planen} (planifier), \all{vorbereiten} (préparer, particule séparable : \all{hat vorbereitet}), \all{organisieren} (organiser), \all{verwirklichen} (réaliser, concrétiser).
\end{definition}

\begin{exemplebox}[Exemples traduits]
\all{Wir planen ein Projekt.} « Nous planifions un projet. »

\all{Sie hat ihr Ziel erreicht.} « Elle a atteint son but. »

\all{Die Schüler bereiten ihr Projekt gut vor.} « Les élèves préparent bien leur projet. » (attention : \all{vor} se détache et va en fin de phrase !)

\all{Er will seine Idee verwirklichen.} « Il veut concrétiser son idée. »

\all{Wir müssen einen Termin festlegen.} « Nous devons fixer une date. »
\end{exemplebox}

\begin{attention}[Faux amis : der Plan, das Ziel, der Termin]
\all{Der Plan} est un vrai ami (« le plan »), mais attention au genre : \textbf{der} Plan, pas \all{*die Plan}. \all{Das Ziel} signifie « le but » et non « le zèle ». \all{Der Termin} est un rendez-vous ou une date limite, pas un « terme » au sens de mot. Et \all{realisieren} signifie en allemand « se rendre compte » ; pour dire « réaliser un projet », utilisez \all{verwirklichen} ou \all{durchführen}.
\end{attention}

\courssub{Familles de mots : un mot, toute une famille}

L'allemand construit des familles de mots très régulières. Les connaître permet de multiplier votre vocabulaire sans effort :

\begin{center}
\begin{tabularx}{\linewidth}{|l|l|X|}
\hline
\rowcolor{popblueL} Verbe & Noms dérivés & Adjectif / Autre \\
\hline
sparen (épargner) & das Sparen (l'épargne), der Sparer (l'épargnant), die Ersparnisse (les économies) & sparsam (économe) \\
\hline
planen (planifier) & der Plan (le plan), die Planung (la planification) & geplant (prévu) \\
\hline
zahlen (payer) & die Zahlung (le paiement), der Zahler (le payeur) & bezahlen (payer, régler) \\
\hline
kosten (coûter) & die Kosten (les frais) & kostenlos (gratuit) \\
\hline
\end{tabularx}
\end{center}

Remarquez le procédé : le \textbf{radical du verbe} devient un nom (l'infinitif nominalisé est toujours neutre : \all{das Sparen}, \all{das Planen}) ; le suffixe \all{-ung} crée des noms féminins d'action (\all{die Planung}, \all{die Zahlung}) ; le suffixe \all{-er} désigne la personne qui fait l'action (\all{der Sparer}, \all{der Zahler}).

\courssub{Expressions utiles à mémoriser}

\begin{aretenir}[Les cinq expressions clés du chapitre]
\begin{itemize}
\item \all{mit Geld umgehen} = gérer son argent (litt. « se comporter avec l'argent » ; particule séparable : \all{Er geht gut mit Geld um.}) : \all{Man muss lernen, mit Geld umzugehen.}
\item \all{Geld ausgeben für} (+ acc.) = dépenser de l'argent pour : \all{Er gibt Geld für Bücher aus.}
\item \all{Geld sparen für} (+ acc.) = économiser pour : \all{Wir sparen Geld für die Reise.}
\item \all{ein Projekt durchführen} = mener à bien un projet : \all{Die Klasse führt ein Projekt durch.} (particule séparable !)
\item \all{einen Plan machen} = faire un plan : \all{Am Anfang machen wir einen Plan.}
\end{itemize}
\end{aretenir}

\courssub{Carte mentale : le lexique en un coup d'œil}

\begin{popfigure}
\begin{tikzpicture}[mindmap, grow cyclic, every node/.style={concept}, concept color=popblue, text=white,
level 1/.append style={level distance=3.4cm, sibling angle=90},
level 2/.append style={level distance=2.2cm, sibling angle=45}]
\node[concept, font=\large\bfseries] {Geld \& Projekte}
  child[concept color=poporange] { node[concept] {Budget}
    child { node[concept, font=\small] {das Budget} }
    child { node[concept, font=\small] {der Betrag} }
    child { node[concept, font=\small] {die Kosten} }
  }
  child[concept color=popgreen] { node[concept] {Einnahmen \& Ausgaben}
    child { node[concept, font=\small] {verdienen} }
    child { node[concept, font=\small] {ausgeben} }
    child { node[concept, font=\small] {bezahlen} }
  }
  child[concept color=popgold] { node[concept] {Sparen}
    child { node[concept, font=\small] {das Sparkonto} }
    child { node[concept, font=\small] {sparsam} }
    child { node[concept, font=\small] {verschwenden} }
  }
  child[concept color=poppurple] { node[concept] {Projekt \& Plan}
    child { node[concept, font=\small] {das Ziel} }
    child { node[concept, font=\small] {planen} }
    child { node[concept, font=\small] {der Termin} }
  };
\end{tikzpicture}
\legende{Carte mentale du vocabulaire : quatre champs lexicaux autour de « Geld \& Projekte ».}
\end{popfigure}

\courssub{Tableau récapitulatif général}

\begin{center}
\begin{tabularx}{\linewidth}{|l|l|X|}
\hline
\rowcolor{popblueL} Deutsch (art. + Plur.) & Français & Exemple \\
\hline
das Geld (-er) & l'argent & \all{Geld ist nicht alles.} \\
\hline
das Budget (-s) & le budget & \all{Unser Budget ist klein.} \\
\hline
das Taschengeld (-er) & l'argent de poche & \all{Ich bekomme Taschengeld.} \\
\hline
der Betrag (-⸚e) & le montant & \all{Der Betrag ist hoch.} \\
\hline
die Einnahme (-n) & la recette & \all{Die Einnahmen steigen.} \\
\hline
die Ausgabe (-n) & la dépense & \all{Die Ausgaben sinken.} \\
\hline
das Gehalt (-⸚er) & le salaire & \all{Er bekommt sein Gehalt.} \\
\hline
das Sparkonto (-en) & le compte d'épargne & \all{Sie eröffnet ein Sparkonto.} \\
\hline
der Sparplan (-⸚e) & le plan d'épargne & \all{Wir machen einen Sparplan.} \\
\hline
das Projekt (-e) & le projet & \all{Das Projekt beginnt.} \\
\hline
der Plan (-⸚e) & le plan & \all{Der Plan ist gut.} \\
\hline
das Ziel (-e) & le but & \all{Er erreicht sein Ziel.} \\
\hline
der Termin (-e) & le rendez-vous & \all{Der Termin ist am Montag.} \\
\hline
\end{tabularx}
\end{center}

\courssub{Exercice résolu}

\begin{exoresolu}[Complétez avec le mot qui convient]
Complétez les phrases avec : \all{Budget} — \all{Ausgaben} — \all{spart} — \all{verwirklichen} — \all{Ziel} — \all{verschwendet}.

1. Marie \_\_\_\_\_\_\_\_ jeden Monat 5 000 Francs für eine Reise nach Deutschland.
2. Die \_\_\_\_\_\_\_\_ sind diesen Monat höher als die Einnahmen.
3. Unser \_\_\_\_\_\_\_\_ für das Schulprojekt ist 50 000 Francs.
4. Paul \_\_\_\_\_\_\_\_ sein Geld für unnütze Dinge.
5. Ihr \_\_\_\_\_\_\_\_ ist es, Ärztin zu werden.
6. Die Schüler wollen ihr Projekt im Juni \_\_\_\_\_\_\_\_.
\tcblower
\textbf{Solution détaillée :}

1. \all{Marie \textbf{spart} jeden Monat 5 000 Francs für eine Reise nach Deutschland.} — « Marie économise 5 000 francs par mois pour un voyage en Allemagne. » Verbe \all{sparen}, construit ici avec \all{für} + accusatif ; à la 3\up{e} personne du singulier : \all{spart}.

2. \all{Die \textbf{Ausgaben} sind diesen Monat höher als die Einnahmen.} — « Les dépenses sont ce mois-ci plus élevées que les recettes. » Le verbe \all{sind} au pluriel impose le pluriel \all{Ausgaben} (et non le singulier).

3. \all{Unser \textbf{Budget} für das Schulprojekt ist 50 000 Francs.} — « Notre budget pour le projet scolaire est de 50 000 francs. » Nom neutre : \all{unser Budget}.

4. \all{Paul \textbf{verschwendet} sein Geld für unnütze Dinge.} — « Paul gaspille son argent pour des choses inutiles. » Verbe faible régulier : \all{verschwendet}.

5. \all{Ihr \textbf{Ziel} ist es, Ärztin zu werden.} — « Son but est de devenir médecin. » Nom neutre : \all{ihr Ziel}.

6. \all{Die Schüler wollen ihr Projekt im Juni \textbf{verwirklichen}.} — « Les élèves veulent réaliser leur projet en juin. » Infinitif après \all{wollen}, placé en fin de phrase.
\end{exoresolu}

\begin{attention}[Récapitulatif des pièges pour les francophones]
\begin{enumerate}
\item Toujours apprendre le nom \textbf{avec son article} : \all{das Geld}, \all{der Plan}, \all{die Ausgabe} — jamais le mot seul.
\item \all{Einnahmen} et \all{Ausgaben} : au pluriel dans les contextes budgétaires.
\item \all{sparen für/auf}, jamais \all{*sparen zu}.
\item \all{ausgeben}, \all{einzahlen}, \all{vorbereiten}, \all{durchführen}, \all{umgehen} : particule séparable en fin de phrase au présent.
\item « Réaliser un projet » = \all{ein Projekt verwirklichen/durchführen}, pas \all{*realisieren}.
\item Majuscule à tous les noms : \all{das Geld}, \all{der Betrag}, \all{das Ziel}.
\end{enumerate}
\end{attention}

Vous disposez maintenant de tout le lexique nécessaire pour parler d'argent, de budget et de projets. Dans la section suivante, nous verrons comment comparer les dépenses et les revenus grâce au \textbf{comparatif} et au \textbf{superlatif} : \all{mehr}, \all{besser}, \all{am höchsten}...

\coursec{Grammaire 1 : comparatif et superlatif}

Dans la section précédente, nous avons constitué notre vocabulaire autour de \all{Geld, Budget und Projekte} : \all{das Budget}, \all{die Ausgaben}, \all{die Einnahmen}, \all{sparen}, \all{das Projekt}. Or, quand on gère un budget, on compare sans cesse : ce projet coûte \emph{plus cher} que l'autre, cette dépense est \emph{la plus importante}, ce plan est \emph{meilleur}. Comparer, c'est exactement ce que permettent le \cle{comparatif} (\all{der Komparativ}) et le \cle{superlatif} (\all{der Superlativ}), deux formes de l'adjectif que nous allons maintenant étudier en profondeur.

\courssub{Observation : les degrés de l'adjectif dans le texte}

Reprenons trois phrases tirées de notre texte support \all{Lernen, mit Geld umzugehen} :

\begin{exemplebox}[Phrases du texte à observer]
\all{Früher war er unvorsichtiger, jetzt plant er besser.}\\
« Avant, il était plus imprudent ; maintenant, il planifie mieux. »

\all{Sein wichtigstes Projekt ist ein eigenes Fahrrad.}\\
« Son projet le plus important est un vélo à lui. »

\all{Seine größten Ausgaben sind das Essen und das Handy.}\\
« Ses plus grandes dépenses sont la nourriture et le téléphone portable. »
\end{exemplebox}

Observons les formes en gras implicite : \all{unvorsichtig} devient \all{unvorsichtig\textbf{er}} ; \all{gut} devient \all{besser} ; \all{wichtig} devient \all{wichtig\textbf{st}es} ; \all{groß} devient \all{größ\textbf{st}en}. L'adjectif allemand possède donc trois \cle{degrés} :

\begin{center}
\begin{tabularx}{\linewidth}{|X|X|X|}
\hline
\rowcolor{popblueL} Degré & Allemand & Fonction \\
\hline
Positif & \all{der Positiv} : \all{teuer} (cher) & qualité simple, sans comparaison \\
\hline
Comparatif & \all{der Komparativ} : \all{teurer} & supériorité : « plus cher » \\
\hline
Superlatif & \all{der Superlativ} : \all{am teuersten} & degré maximal : « le plus cher » \\
\hline
\end{tabularx}
\end{center}

\courssub{Le comparatif : formation et emploi}

\begin{propriete}[Formation du comparatif]
Le comparatif se forme en ajoutant le suffixe \cle{-er} à l'adjectif, et le second terme de la comparaison est introduit par \cle{als} (« que ») :

\all{Adjectif + -er ... als}

\all{Das Leben in der Stadt ist teurer als auf dem Land.}\\
« La vie en ville est plus chère qu'à la campagne. »

\all{Das Handy ist teurer als das Buch.}\\
« Le téléphone portable est plus cher que le livre. »
\end{propriete}

Point capital pour les francophones : le comparatif allemand est \cle{synthétique}. Là où le français ajoute un mot (« plus ») devant l'adjectif, l'allemand modifie l'adjectif lui-même par un suffixe :

\begin{center}
\begin{tabularx}{\linewidth}{|X|X|X|}
\hline
\rowcolor{popblueL} Français & Deutsch & Remarque \\
\hline
plus cher que & \all{teurer als} & suffixe \all{-er}, pas de mot ajouté \\
\hline
plus grand que & \all{größer als} & avec Umlaut (voir plus bas) \\
\hline
meilleur que & \all{besser als} & forme irrégulière \\
\hline
plus ... que & \all{...-er als} & \all{als} = « que » de comparaison \\
\hline
aussi ... que & \all{so ... wie} & égalité : \all{so teuer wie} \\
\hline
de plus en plus cher & \all{immer teurer} & \all{immer} + comparatif \\
\hline
\end{tabularx}
\end{center}

\begin{exemplebox}[Exemples traduits]
\all{Ich spare mehr als mein Bruder.} — « J'économise plus que mon frère. »

\all{Ein Fahrrad ist billiger als ein Auto.} — « Un vélo est moins cher (plus bon marché) qu'une voiture. »

\all{In Douala ist das Leben teurer als in einem Dorf.} — « À Douala, la vie est plus chère que dans un village. »

\all{Die Preise werden immer höher.} — « Les prix deviennent de plus en plus élevés. »

\all{Mein Plan ist besser als dein Plan.} — « Mon plan est meilleur que ton plan. »
\end{exemplebox}

Quand le comparatif est \emph{épithète} (placé devant le nom), il prend en plus la désinence d'accord normale de l'adjectif : \all{Ich brauche ein billigeres Handy.} « J'ai besoin d'un téléphone moins cher. » ; \all{Sie hat einen besseren Plan.} « Elle a un meilleur plan. »

\courssub{L'Umlaut au comparatif et au superlatif}

\begin{propriete}[Umlaut des adjectifs monosyllabiques]
Les adjectifs \cle{monosyllabiques} dont la voyelle radicale est \cle{a}, \cle{o} ou \cle{u} prennent généralement l'\cle{Umlaut} (ä, ö, ü) au comparatif et au superlatif :

\all{alt $\rightarrow$ älter $\rightarrow$ am ältesten} (vieux / plus vieux / le plus vieux)

\all{groß $\rightarrow$ größer $\rightarrow$ am größten}

\all{jung $\rightarrow$ jünger $\rightarrow$ am jüngsten}

\all{kalt $\rightarrow$ kälter $\rightarrow$ am kältesten}

\all{stark $\rightarrow$ stärker $\rightarrow$ am stärksten}

\all{lang $\rightarrow$ länger $\rightarrow$ am längsten}
\end{propriete}

\begin{attention}[Exceptions à l'Umlaut]
Quelques adjectifs monosyllabiques en a, o, u ne prennent \textbf{pas} l'Umlaut : \all{klar $\rightarrow$ klarer} (clair), \all{voll $\rightarrow$ voller} (plein), \all{flach $\rightarrow$ flacher} (plat). Les adjectifs \emph{polysyllabiques} ne prennent jamais l'Umlaut : \all{interessant $\rightarrow$ interessanter}, \all{wichtig $\rightarrow$ wichtiger}, \all{praktisch $\rightarrow$ praktischer}.
\end{attention}

\courssub{Le superlatif : deux constructions}

\begin{propriete}[Formation du superlatif]
Le superlatif a \textbf{deux constructions} selon la place de l'adjectif :

\textbf{1. Après un verbe} (adjectif attribut du sujet, en fin de phrase) : \cle{am + adjectif + -(e)sten}

\all{Dieser Plan ist am besten.} — « Ce plan est le meilleur. »

\all{Im Januar sind die Preise am höchsten.} — « En janvier, les prix sont les plus élevés. »

\textbf{2. Devant un nom} (adjectif épithète) : \cle{article + adjectif + -st + désinence d'accord}

\all{der beste Plan} — « le meilleur plan » ; \all{das beste Projekt} — « le meilleur projet » ; \all{die größten Ausgaben} — « les plus grandes dépenses ».

\all{Das ist die größte Ausgabe.} — « C'est la plus grande dépense. »
\end{propriete}

\begin{propriete}[Le suffixe -esten après certaines consonnes]
Après les adjectifs se terminant par \cle{-d}, \cle{-t}, \cle{-s}, \cle{-ß}, \cle{-sch} ou \cle{-z}, on insère un \cle{e} euphonique au superlatif : \all{-(e)sten} devient \all{-esten} :

\all{kalt $\rightarrow$ am kältesten} ; \all{interessant $\rightarrow$ am interessantesten} ; \all{heiß $\rightarrow$ am heißesten} ; \all{frisch $\rightarrow$ am frischesten} ; \all{kurz $\rightarrow$ am kürzesten}.
\end{propriete}

\begin{propriete}[Adjectifs en -el, -er, -en]
Au comparatif, les adjectifs en \cle{-el} perdent le \all{e} du radical : \all{dunkel $\rightarrow$ dunkler} (plus sombre). Les adjectifs en \cle{-er} perdent aussi ce \all{e} : \all{teuer $\rightarrow$ teurer} (et non *\all{teuerer}). Les adjectifs en \cle{-en} le conservent normalement : \all{trocken $\rightarrow$ trockener} (plus sec). Au superlatif, la forme est régulière : \all{am dunkelsten}, \all{am teuersten}, \all{am trockensten}.
\end{propriete}

\begin{exemplebox}[Exemples traduits]
\all{Am liebsten reise ich nach Österreich.} — « C'est en Autriche que je préfère voyager. »

\all{Der älteste Bruder spart am meisten.} — « Le frère aîné économise le plus. »

\all{Yaoundé ist eine der größten Städte Kameruns.} — « Yaoundé est l'une des plus grandes villes du Cameroun. »

\all{Im Dezember sind die Ausgaben am höchsten.} — « En décembre, les dépenses sont les plus élevées. »
\end{exemplebox}

\courssub{Comparatifs et superlatifs irréguliers}

Comme en français (« bon / meilleur / le meilleur »), l'allemand possède des degrés \cle{irréguliers} qu'il faut apprendre par cœur :

\begin{center}
\begin{tabular}{|l|l|l|l|}
\hline
\rowcolor{popblueL} Positiv & Komparativ & Superlativ & Français \\
\hline
\all{gut} & \all{besser} & \all{am besten} & bon / meilleur / le meilleur \\
\hline
\all{viel} & \all{mehr} & \all{am meisten} & beaucoup / plus / le plus \\
\hline
\all{gern} & \all{lieber} & \all{am liebsten} & volontiers / plutôt / de préférence \\
\hline
\all{hoch} & \all{höher} & \all{am höchsten} & haut / plus haut / le plus haut \\
\hline
\all{nah} & \all{näher} & \all{am nächsten} & proche / plus proche / le plus proche \\
\hline
\end{tabular}
\end{center}

Remarques : \all{hoch} perd le \all{c} au comparatif (\all{höher}) ; \all{nah} donne \all{näher} et \all{am nächsten} ; \all{gern} est un adverbe : \all{Ich trinke gern Wasser, aber lieber Saft.} « Je bois volontiers de l'eau, mais plutôt du jus. »

\begin{popfigure}
\begin{center}
\begin{tikzpicture}[
  deg/.style={rectangle, rounded corners=3pt, draw=popline, fill=popcream, align=center, font=\small, minimum width=2.9cm, minimum height=0.85cm},
  hdr/.style={rectangle, rounded corners=3pt, fill=popblue, text=white, align=center, font=\small\bfseries, minimum width=2.9cm, minimum height=0.85cm}
]
\node[hdr] (h1) at (0,0) {Positiv};
\node[hdr] (h2) at (3.6,0) {Komparativ};
\node[hdr] (h3) at (7.2,0) {Superlativ};
\node[deg] at (0,-1.1) {gut};
\node[deg, fill=poporangeL] at (3.6,-1.1) {besser};
\node[deg, fill=popgreenL] at (7.2,-1.1) {am besten};
\node[deg] at (0,-2.1) {viel};
\node[deg, fill=poporangeL] at (3.6,-2.1) {mehr};
\node[deg, fill=popgreenL] at (7.2,-2.1) {am meisten};
\node[deg] at (0,-3.1) {teuer};
\node[deg, fill=poporangeL] at (3.6,-3.1) {teurer};
\node[deg, fill=popgreenL] at (7.2,-3.1) {am teuersten};
\node[deg] at (0,-4.1) {groß};
\node[deg, fill=poporangeL] at (3.6,-4.1) {größer};
\node[deg, fill=popgreenL] at (7.2,-4.1) {am größten};
\end{tikzpicture}
\end{center}
\legende{Les trois degrés de l'adjectif : réguliers (teuer, groß) et irréguliers (gut, viel).}
\end{popfigure}

\begin{popfigure}
\begin{center}
\begin{tikzpicture}[
  stufe/.style={rectangle, rounded corners=4pt, draw=popline, align=center, font=\small, minimum width=3.4cm, minimum height=0.9cm},
  arr/.style={-{Stealth[length=3mm]}, very thick, poporange}
]
\node[stufe, fill=popblueL, align=center] (a) at (0,0){billig\\ \emph{bon marché}};
\node[stufe, fill=poporangeL, align=center] (b) at (0,1.7){billiger\\ \emph{moins cher}};
\node[stufe, fill=popgreenL, align=center] (c) at (0,3.4){am billigsten\\ \emph{le moins cher}};
\draw[arr] (a) -- (b);
\draw[arr] (b) -- (c);
\node[font=\small\itshape, text=popmuted, align=center] at (4.2,1.7) {degré croissant\\ de la qualité};
\end{tikzpicture}
\end{center}
\legende{L'échelle des degrés : du positif au superlatif.}
\end{popfigure}

\courssub{Texte d'application : comparer pour mieux gérer}

\begin{textebox}[Zwei Schüler, zwei Budgets]
Aminu und Béatrice sind Schüler in Yaoundé und lernen, ihr Geld besser zu planen. Aminu ist älter als Béatrice, aber sie ist vorsichtiger mit Geld. „Früher war ich unvorsichtiger“, sagt Aminu. „Meine größten Ausgaben waren Snacks und Spiele. Jetzt plane ich besser: mein wichtigstes Projekt ist ein gebrauchtes Fahrrad. Es ist billiger als ein neues, aber für mich ist es die beste Lösung.“

Béatrice vergleicht gern die Preise. „Auf dem Markt ist das Obst billiger als im Supermarkt, und oft auch frischer“, erklärt sie. „Am liebsten kaufe ich am frühen Morgen, dann sind die Preise am niedrigsten.“ Sie spart mehr als ihr Bruder, weil sie jede Woche einen kleinen Betrag zur Seite legt. „Je mehr ich spare, desto näher komme ich meinem Ziel“, sagt sie stolz.

Ihre Lehrerin findet: „Von allen Schülern plant Béatrice am besten. Ihr Budget ist das klarste, und ihre Projekte sind die realistischsten.“ Am Ende des Monats hat Béatrice am meisten gespart, aber Aminu hat die größte Veränderung gemacht: er ist jetzt viel sparsamer als früher. Beide wissen: Wer besser plant, erreicht seine Ziele schneller.

\textbf{Glossaire :} \all{gebraucht} = d'occasion ; \all{die Lösung, -en} = la solution ; \all{das Obst} = les fruits ; \all{niedrig} = bas ; \all{der Betrag, -ä-e} = le montant ; \all{zur Seite legen} = mettre de côté ; \all{das Ziel, -e} = le but ; \all{stolz} = fier ; \all{realistisch} = réaliste ; \all{die Veränderung, -en} = le changement ; \all{sparsam} = économe ; \all{erreichen} = atteindre.
\end{textebox}

\textbf{Questions de grammaire sur le texte :}
\begin{enumerate}
\item Relève trois comparatifs et trois superlatifs du texte et donne le positif de chacun.
\item Pourquoi \all{älter} et \all{größten} prennent-ils un Umlaut, alors que \all{wichtigstes} n'en prend pas ?
\item Traduis : \all{Je mehr ich spare, desto näher komme ich meinem Ziel.}
\end{enumerate}

\courssub{Exercice résolu et pièges à éviter}

\begin{exoresolu}[Mets l'adjectif au degré demandé]
Énoncé. Complète avec la forme correcte :
\begin{enumerate}
\item Das Fahrrad ist \ldots\ (billig) als das Auto.
\item Von allen Plänen ist dieser \ldots\ (gut).
\item Meine Schwester ist \ldots\ (alt) als ich.
\item Im Sommer sind die Tage am \ldots\ (lang).
\item Das ist die \ldots\ (groß) Ausgabe des Monats.
\end{enumerate}
\tcblower
Solution détaillée :
\begin{enumerate}
\item \all{billiger} : comparatif régulier, adjectif + \all{-er} ; « Le vélo est moins cher que la voiture. »
\item \all{am besten} : superlatif irrégulier de \all{gut} après le verbe \all{sein} ; « De tous les plans, celui-ci est le meilleur. »
\item \all{älter} : comparatif avec Umlaut (\all{alt} est monosyllabique en a) ; « Ma sœur est plus âgée que moi. »
\item \all{längsten} : superlatif avec Umlaut + \all{-sten} ; « En été, les jours sont les plus longs. »
\item \all{größte} : superlatif épithète accordé avec \all{die Ausgabe} (féminin, nominatif) : \all{die größt-e} ; « C'est la plus grande dépense du mois. »
\end{enumerate}
\end{exoresolu}

\begin{attention}[Ne dis jamais « mehr besser » !]
Erreur typique des francophones : traduire « plus » par \all{mehr} devant un adjectif : ✗ \all{mehr teuer}, ✗ \all{mehr gut}. En allemand, le comparatif est synthétique : ✓ \all{teurer}, ✓ \all{besser}. \all{Mehr} s'emploie \textbf{seul}, devant un nom, au sens de « plus de » : \all{mehr Geld} (plus d'argent), \all{mehr Zeit} (plus de temps), \all{mehr Ausgaben} (plus de dépenses).
\end{attention}

\begin{attention}[« als » et non « wie » après un comparatif]
Après un comparatif, on utilise toujours \cle{als} et jamais \all{wie} : ✗ \all{besser wie} $\rightarrow$ ✓ \all{besser als}. \all{Wie} s'emploie dans la comparaison d'\textbf{égalité} : \all{so gut wie} « aussi bon que », \all{Er ist so alt wie ich.} « Il est aussi âgé que moi. » Retiens : \all{so + Positiv + wie} (égalité) / \all{Komparativ + als} (supériorité). N'oublie pas non plus l'accord du superlatif épithète : \all{der beste Plan}, \all{das beste Projekt}, \all{die beste Idee}.
\end{attention}

\begin{exercice}[À toi de jouer]
\begin{enumerate}
\item Traduis en allemand : a) « Mon projet est plus important que ton projet. » b) « C'est la dépense la plus élevée. » c) « La vie devient de plus en plus chère. » d) « J'aime mieux économiser que dépenser. »
\item Mets au comparatif puis au superlatif : \all{schnell, schön, warm, interessant, viel}.
\item Corrige les erreurs : a) \all{Das Buch ist mehr interessant wie der Film.} b) \all{Sie ist die beste Schülerin als ihre Schwester.}
\end{enumerate}
\end{exercice}

\begin{corrige}[Exercice « À toi de jouer »]
1. a) \all{Mein Projekt ist wichtiger als dein Projekt.} b) \all{Das ist die höchste Ausgabe.} c) \all{Das Leben wird immer teurer.} d) \all{Ich spare lieber, als dass ich Geld ausgebe.} (ou, plus simple : \all{Ich spare lieber, statt Geld auszugeben.})

2. \all{schnell $\rightarrow$ schneller $\rightarrow$ am schnellsten} ; \all{schön $\rightarrow$ schöner $\rightarrow$ am schönsten} ; \all{warm $\rightarrow$ wärmer $\rightarrow$ am wärmsten} ; \all{interessant $\rightarrow$ interessanter $\rightarrow$ am interessantesten} ; \all{viel $\rightarrow$ mehr $\rightarrow$ am meisten}.

3. a) ✓ \all{Das Buch ist interessanter als der Film.} (pas de \all{mehr} devant l'adjectif, \all{als} après le comparatif). b) ✓ \all{Sie ist eine bessere Schülerin als ihre Schwester.} (comparatif + \all{als} ; le superlatif \all{die beste} ne se combine pas avec \all{als}).
\end{corrige}

\begin{aretenir}[L'essentiel du comparatif et du superlatif]
\begin{itemize}
\item Comparatif : adjectif + \cle{-er}, second terme introduit par \cle{als} : \all{teurer als}.
\item Superlatif : \cle{am + adjectif + -(e)sten} après un verbe ; \cle{der/die/das ...-ste} devant un nom, avec accord.
\item Umlaut pour les monosyllabes en a, o, u : \all{alt $\rightarrow$ älter}, \all{groß $\rightarrow$ größer}.
\item \all{-esten} après -d, -t, -s, -ß, -sch, -z : \all{am interessantesten}.
\item Irréguliers à mémoriser : \all{gut – besser – am besten} ; \all{viel – mehr – am meisten} ; \all{gern – lieber – am liebsten} ; \all{hoch – höher – am höchsten} ; \all{nah – näher – am nächsten}.
\item « De plus en plus » = \all{immer} + comparatif : \all{immer teurer}.
\end{itemize}
\end{aretenir}

\begin{savaistu}[Le comparatif au Bac]
Au baccalauréat, le comparatif et le superlatif reviennent très souvent dans les questions de grammaire à trous ou de transformation de phrases. Le trio \all{gut – besser – am besten} est le plus fréquent de tous : retiens-le par cœur, avec \all{viel – mehr – am meisten}. Et souviens-toi : les Allemands eux-mêmes plaisantent parfois sur l'erreur \all{mehr besser}, que font les petits enfants en apprenant à parler — preuve que même les natifs doivent dompter cette règle !
\end{savaistu}

\coursec{Grammaire 2 : conditionnelles (wenn) et Futur I}

Dans la section précédente, nous avons appris à comparer les choses : un projet peut être \emph{größer}, \emph{besser} ou \emph{teurer} qu'un autre. Mais gérer un budget, c'est aussi \emph{prévoir} : « Si j'économise assez, je réaliserai mon projet. » Cette phrase contient deux outils grammaticaux essentiels : la subordonnée conditionnelle avec \all{wenn} et le futur avec \all{werden}. C'est exactement ce que nous allons étudier maintenant.

\courssub{Observation : repérons les verbes}

\begin{textebox}[Phrase d'observation]
\all{„Wenn er jeden Monat 20 Euro spart, wird er im Sommer nach Österreich reisen.“}

« S'il économise 20 euros chaque mois, il voyagera en Autriche cet été. »
\end{textebox}

Observons attentivement cette phrase et posons-nous trois questions :

\begin{enumerate}
\item Combien y a-t-il de verbes conjugués ? Deux : \all{spart} (il économise) et \all{wird} (auxiliaire du futur).
\item Où se trouvent-ils ? \all{spart} est à la \textbf{toute fin} de la première proposition ; \all{wird} est en \textbf{première position} de la deuxième proposition, juste après la virgule.
\item Où est l'infinitif ? \all{reisen} (voyager) est à la \textbf{toute fin} de la phrase.
\end{enumerate}

\begin{popfigure}
\begin{tikzpicture}[font=\small]
\node[draw=popblue, fill=popblueL, rounded corners, align=center, minimum width=5.6cm, minimum height=1.1cm] (sub) at (0,0) {\all{Wenn er jeden Monat 20 Euro}\\ \textbf{\all{spart}}};
\node[draw=poporange, fill=poporangeL, rounded corners, align=center, minimum width=5.6cm, minimum height=1.1cm] (haupt) at (7.2,0) {\textbf{\all{wird}} er im Sommer nach Österreich\\ \textbf{\all{reisen}}};
\node[below=0.15cm of sub, text=popblue, align=center] {subordonnée : verbe conjugué\\ à la \textbf{dernière place}};
\node[below=0.15cm of haupt, text=poporange, align=center] {principale : auxiliaire \all{wird} en position 1,\\ infinitif \all{reisen} à la fin};
\draw[-{Stealth[length=3mm]}, very thick, poppurple] (sub.east) -- (haupt.west) node[midway, above, text=popdark] {,};
\draw[-{Stealth[length=2.5mm]}, thick, popgreen] ([yshift=0.9cm]sub.north) .. controls (3.6,1.8) .. ([yshift=0.9cm]haupt.north) node[midway, above, text=popgreen, align=center] {la subordonnée en tête\\ « pousse » le verbe de la principale en position 1};
\end{tikzpicture}
\legende{La phrase à deux verbes : \all{wenn} + verbe en fin de subordonnée, puis verbe en tête de principale.}
\end{popfigure}

\courssub{La subordonnée conditionnelle avec \all{wenn}}

\begin{propriete}[La place du verbe dans la subordonnée avec \all{wenn}]
\all{wenn} (« si », « quand ») est une \cle{conjonction de subordination}. Dans la subordonnée qu'elle introduit, le \textbf{verbe conjugué est rejeté à la dernière place} de la proposition :

\all{Wenn er genug Geld hat, ...} = S'il a assez d'argent...

\all{Wenn ich Geld habe, kaufe ich ein Buch.} = Si j'ai de l'argent, j'achète un livre.
\end{propriete}

\begin{propriete}[L'inversion quand la subordonnée est en tête]
Quand la subordonnée est placée \textbf{au début} de la phrase, elle occupe la position 1. Le verbe conjugué de la proposition principale vient donc \textbf{immédiatement après la virgule}, en position 1 de la principale, et le sujet le suit :

\all{Wenn ich Zeit habe, besuche ich dich.} = Si j'ai le temps, je te rends visite.

On peut ajouter \all{dann} (« alors ») pour insister : \all{Wenn ich Zeit habe, (dann) besuche ich dich.}

Si la subordonnée est placée \textbf{après} la principale, l'ordre de la principale reste normal (verbe en position 2) : \all{Ich besuche dich, wenn ich Zeit habe.}
\end{propriete}

\begin{center}
\begin{tabularx}{\linewidth}{|X|X|}
\hline
\rowcolor{popblueL} Subordonnée en tête & Subordonnée après \\
\hline
\all{Wenn ich Geld spare, werde ich reisen.} & \all{Ich werde reisen, wenn ich Geld spare.} \\
\hline
\all{Wenn es regnet, bleiben wir zu Hause.} & \all{Wir bleiben zu Hause, wenn es regnet.} \\
\hline
\all{Wenn du fleißig lernst, bestehst du die Prüfung.} & \all{Du bestehst die Prüfung, wenn du fleißig lernst.} \\
\hline
\end{tabularx}
\end{center}

\courssub{Ne pas confondre : \all{wenn}, \all{als}, \all{wann}}

\begin{definition}[Trois mots, trois emplois]
\begin{itemize}
\item \cle{wenn} = « si » (condition) ou « quand » (répétition, habitude, présent ou futur) : \all{Wenn ich Zeit habe, helfe ich dir.} / \all{Wenn ich früher Geld hatte, kaufte ich Bücher.} (chaque fois que j'avais de l'argent).
\item \cle{als} = « quand » pour un \textbf{événement unique du passé} : \all{Als ich ein Kind war, lebten wir in Douala.} = Quand j'étais enfant, nous vivions à Douala.
\item \cle{wann} = « quand » dans une \textbf{question} directe ou indirecte : \all{Wann kommst du?} = Quand viens-tu ? / \all{Ich weiß nicht, wann er kommt.} = Je ne sais pas quand il vient.
\end{itemize}
\end{definition}

\begin{attention}[Piège fréquent : \all{wann} ou \all{wenn} ?]
Le francophone traduit souvent « quand » automatiquement par \all{wann}. Mais dans une phrase affirmative qui exprime une condition ou un moment futur, il faut \all{wenn} : « Quand j'aurai le temps, je viendrai » = \all{Wenn ich Zeit habe, werde ich kommen.} (et non \all{wann}). Réservez \all{wann} aux questions.
\end{attention}

\courssub{Le Futur I : formation}

\begin{propriete}[Formation du Futur I]
Le Futur I se forme avec l'auxiliaire \cle{werden} conjugué au présent (en position 2 dans la phrase) + l'\textbf{infinitif} du verbe principal placé à la \textbf{dernière place} :

\all{Sie wird nächstes Jahr ein neues Projekt beginnen.} = Elle commencera un nouveau projet l'année prochaine.

\all{Ich werde sparen.} = J'économiserai. \all{Du wirst reisen.} = Tu voyageras. \all{Wir werden das Projekt verwirklichen.} = Nous réaliserons le projet.
\end{propriete}

\begin{center}
\begin{tabular}{|l|l|l|}
\hline
\rowcolor{popblueL} Personne & \all{werden} au présent & Exemple au Futur I \\
\hline
ich & werde & \all{Ich werde ein Budget machen.} \\
\hline
du & wirst & \all{Du wirst Geld sparen.} \\
\hline
er/sie/es & wird & \all{Er wird nach Österreich reisen.} \\
\hline
wir & werden & \all{Wir werden das Projekt planen.} \\
\hline
ihr & werdet & \all{Ihr werdet die Ausgaben zählen.} \\
\hline
sie/Sie & werden & \all{Sie werden Erfolg haben.} \\
\hline
\end{tabular}
\end{center}

\begin{attention}[Attention à la 2\textsuperscript{e} et 3\textsuperscript{e} personnes du singulier]
\all{werden} est irrégulier au présent : \all{du wir\textbf{st}}, \all{er/sie/es wir\textbf{d}}. Ne dites jamais ✗ \all{du werdest} ou ✗ \all{er werdet}.
\end{attention}

\begin{popfigure}
\begin{tikzpicture}[font=\small]
\draw[-{Stealth[length=3mm]}, very thick, popdark] (0,0) -- (11,0);
\node[circle, fill=popgreen, inner sep=2.5pt, label=below:{\textbf{Gegenwart}}] (g) at (2,0) {};
\node[above=0.25cm of g, align=center, text=popgreen] {présent\\ \all{Ich spare.}};
\node[circle, fill=poporange, inner sep=2.5pt, label=below:{\textbf{Zukunft}}] (z) at (8.5,0) {};
\node[above=0.25cm of z, align=center, text=poporange] {Futur I\\ \all{Ich werde sparen.}};
\draw[-{Stealth[length=2.5mm]}, thick, poppurple] (4.2,0.9) -- (6.8,0.9) node[midway, above, text=poppurple] {\all{werden} + Infinitiv};
\end{tikzpicture}
\legende{Frise des temps : du présent (\all{Gegenwart}) vers le futur (\all{Zukunft}) grâce à \all{werden} + infinitif.}
\end{popfigure}

\courssub{Emplois du Futur I et comparaison avec le français}

Le Futur I exprime :

\begin{itemize}
\item un \textbf{projet} ou une \textbf{intention} : \all{Wir werden ein neues Projekt beginnen.} = Nous commencerons un nouveau projet.
\item une \textbf{prévision} : \all{Es wird morgen regnen.} = Il pleuvra demain.
\item une \textbf{promesse} : \all{Ich werde dir helfen!} = Je t'aiderai !
\end{itemize}

\begin{center}
\begin{tabularx}{\linewidth}{|X|X|X|}
\hline
\rowcolor{popblueL} Français & Deutsch & Remarque \\
\hline
Si j'ai de l'argent, je voyagerai. & \all{Wenn ich Geld habe, werde ich reisen.} & si + présent → futur = \all{wenn} + présent → Futur I \\
\hline
Demain, je vais à Bafoussam. & \all{Morgen fahre ich nach Bafoussam.} & présent + complément de temps = futur \\
\hline
Le mois prochain, le projet commencera. & \all{Nächsten Monat beginnt das Projekt.} & usage très fréquent du présent \\
\hline
Je deviens professeur. & \all{Ich werde Lehrer.} & \all{werden} = devenir (pas d'infinitif !) \\
\hline
\end{tabularx}
\end{center}

\begin{exemplebox}[Exemples traduits]
\all{Wenn ich sparsam bin, werde ich mein Ziel erreichen.} = Si je suis économe, j'atteindrai mon but.

\all{Wenn es regnet, bleiben wir zu Hause.} = S'il pleut, nous restons à la maison.

\all{Wann kommst du? – Wenn ich Zeit habe.} = Quand viens-tu ? – Quand j'aurai le temps.

\all{Wenn wir genug Geld sparen, werden wir ein Motorrad kaufen.} = Si nous économisons assez d'argent, nous achèterons une moto.
\end{exemplebox}

\begin{attention}[Deux \all{werden} différents]
Ne confondez pas \all{werden} \textbf{auxiliaire du futur} et \all{werden} verbe « \textbf{devenir} » :

\all{Ich werde reisen.} = Je voyagerai. (auxiliaire + infinitif \all{reisen})

\all{Ich werde Lehrer.} = Je deviens professeur. (verbe plein, suivi d'un nom, sans infinitif)

Le sens dépend de la présence ou non d'un infinitif en fin de phrase.
\end{attention}

\begin{attention}[Erreur fréquente : l'ordre des mots après \all{wenn}]
Le francophone garde souvent l'ordre de la phrase principale : ✗ \all{Wenn er spart Geld}. Faux ! Le verbe conjugué doit aller en \textbf{fin} de subordonnée : \all{Wenn er Geld spart}. Et dans la principale qui suit, le verbe vient juste après la virgule : ✗ \all{Wenn er Geld spart, er wird reisen} → \all{Wenn er Geld spart, wird er reisen.}
\end{attention}

\begin{savaistu}[Le présent pour parler du futur]
Les Allemands utilisent très souvent le \textbf{présent} pour parler du futur quand un complément de temps est clair : \all{Nächsten Monat beginnt das Projekt.} = Le mois prochain, le projet commence. Le Futur I sert surtout à insister sur la prévision, l'intention ou la promesse. C'est plus simple que le français, où le futur est presque obligatoire !
\end{savaistu}

\courssub{Texte d'application : les projets de la classe}

\begin{textebox}[Unsere Zukunftspläne]
\all{Die Schüler der Terminale sprechen heute über ihre Pläne. „Wenn ich das Abitur bestehe, werde ich an der Universität Yaoundé Wirtschaft studieren“, sagt Amina. „Wenn ich genug Geld spare, werde ich ein kleines Geschäft eröffnen“, erklärt Junior. Er will Handys verkaufen. „Wenn das Geschäft gut läuft, werde ich meinen Eltern helfen“, fügt er hinzu. Die Lehrerin lächelt: „Wenn ihr fleißig arbeitet, werdet ihr eure Ziele erreichen. Aber wenn ihr nicht spart, werdet ihr eure Projekte nicht verwirklichen.“ Alle Schüler versprechen: „Wir werden jeden Monat Geld sparen!“}

\emph{Glossaire :} das Abitur = le baccalauréat ; bestehen (besteht, bestand, bestanden) = réussir (un examen) ; die Wirtschaft = l'économie, les sciences économiques ; das Geschäft, -e = le magasin, le commerce ; eröffnen = ouvrir ; das Handy, -s = le téléphone portable ; laufen (läuft) = marcher, bien fonctionner ; hinzufügen = ajouter ; fleißig = travailleur, assidu ; das Ziel, -e = le but ; erreichen = atteindre ; verwirklichen = réaliser ; versprechen (verspricht) = promettre.
\end{textebox}

\begin{exoresolu}[Exercice résolu : conditionnelle et Futur I]
Énoncé : Mettez les phrases suivantes à la forme correcte, en commençant par la subordonnée avec \all{wenn}.

a) (er / sparen / genug Geld) → (er / kaufen / ein Fahrrad)

b) (wir / haben / Zeit) → (wir / besuchen / unsere Großeltern)

c) Traduisez : « Si tu travailles sérieusement, tu réussiras ton projet. »
\tcblower
Solution détaillée :

a) Subordonnée : verbe conjugué \all{spart} en fin ; principale : auxiliaire \all{wird} après la virgule, infinitif \all{kaufen} en fin.

\all{Wenn er genug Geld spart, wird er ein Fahrrad kaufen.} = S'il économise assez d'argent, il achètera un vélo.

b) \all{Wenn wir Zeit haben, werden wir unsere Großeltern besuchen.} = Si nous avons le temps, nous rendrons visite à nos grands-parents.

c) \all{Wenn du ernsthaft arbeitest, wirst du dein Projekt verwirklichen.} Attention : \all{du} → \all{wirst} (et non ✗ \all{werdest}), et le verbe \all{arbeitest} va à la fin de la subordonnée.
\end{exoresolu}

\begin{exercice}[À vous de jouer]
1. Conjuguez \all{werden} au présent à toutes les personnes, puis formez une phrase au Futur I avec \all{sparen} à la 1\textsuperscript{re} personne du pluriel.

2. Remettez dans l'ordre : \all{wenn / ich / habe / Geld / werde / ich / nach Deutschland / reisen}.

3. Choisissez \all{wenn}, \all{als} ou \all{wann} : a) \all{... ich ein Kind war, wohnten wir in Bafoussam.} b) \all{... kommst du morgen?} c) \all{... es regnet, bleibe ich zu Hause.}

4. Traduisez : « Si nous économisons chaque mois, nous réaliserons notre projet. »
\end{exercice}

\begin{corrige}[Exercice : À vous de jouer]
1. \all{ich werde, du wirst, er/sie/es wird, wir werden, ihr werdet, sie/Sie werden.} Exemple : \all{Wir werden jeden Monat Geld sparen.}

2. \all{Wenn ich Geld habe, werde ich nach Deutschland reisen.}

3. a) \all{Als} (événement unique du passé) ; b) \all{Wann} (question directe) ; c) \all{Wenn} (condition).

4. \all{Wenn wir jeden Monat sparen, werden wir unser Projekt verwirklichen.}
\end{corrige}

\begin{aretenir}[L'essentiel de la section]
\begin{itemize}
\item Après \all{wenn}, le verbe conjugué va à la \textbf{fin} de la subordonnée ; si la subordonnée est en tête, le verbe de la principale vient \textbf{juste après la virgule}.
\item \all{wenn} = si/quand (condition, répétition) ; \all{als} = quand (passé unique) ; \all{wann} = quand (question).
\item Futur I = \all{werden} conjugué (position 2) + infinitif (dernière place) : \all{Ich werde reisen.}
\item \all{werden} + nom = « devenir » ; \all{werden} + infinitif = futur.
\item Le présent + complément de temps remplace souvent le futur : \all{Morgen fahre ich nach Douala.}
\end{itemize}
\end{aretenir}

\coursec{Méthode du commentaire et commentaire modèle}

Après avoir maîtrisé les structures du conditionnel avec \all{Wenn ich jeden Monat 20.000 CFA spare, werde ich nach Österreich reisen} (subordonnée avec \all{wenn} + Futur I dans la principale) et les comparatifs pour évaluer des dépenses (\all{Das Handy ist teurer als die Kleidung.}), nous abordons maintenant l'épreuve reine de l'écrit au Baccalauréat : le commentaire de texte. Au Cameroun comme en Allemagne, cette épreuve exige une méthode rigoureuse : comprendre, structurer, analyser, s'exprimer. Voici la marche à suivre, appliquée à notre texte support \all{Lernen, mit Geld umzugehen}.

\courssub{La méthode du commentaire de texte au Baccalauréat}

\begin{methode}[Les 4 étapes du commentaire au Bac]
\textbf{Étape 1 — Introduction (1 phrase) :} présenter le texte.
\begin{enumerate}
\item Type de texte : article de magazine jeunesse, article de blog, témoignage, interview…
\item Thème : gestion de l'argent de poche, budget, projet d'épargne.
\item Source / personnage principal : ici Jonas, 17 ans, élève à Yaoundé.
\item Accroche : \all{Der Text „Lernen, mit Geld umzugehen“ handelt von Jonas, einem 17-jährigen Schüler, und seiner Budgetplanung.}
\end{enumerate}

\textbf{Étape 2 — Résumé structuré (3–4 phrases, présent de narration) :} dégager la structure en 3 parties logiques (revenus → dépenses → projet d'épargne) et résumer avec des connecteurs temporels et logiques : \all{zuerst}, \all{dann}, \all{danach}, \all{schließlich}, \all{außerdem}. Utiliser \textbf{ses propres mots}, ne pas recopier de longs passages du texte.

\textbf{Étape 3 — Analyse du message et des procédés (2–3 phrases) :} que veut montrer l'auteur ? Relever :
\begin{itemize}
\item L'idée centrale : planifier et épargner permettent de réaliser un projet.
\item Les procédés : exemple concret (Jonas), chiffres précis (20.000 CFA par mois, 450.000 CFA pour le vol), citation directe (\all{„Ich will diszipliniert sein.“}), contraste revenus/dépenses.
\item Le vocabulaire clé : \all{die Einnahmen}, \all{die Ausgaben}, \all{sparen}, \all{der Plan}, \all{das Projekt}.
\end{itemize}

\textbf{Étape 4 — Conclusion avec avis personnel (1 phrase minimum) :} exprimer son opinion en la justifiant avec \all{weil} (verbe à la fin).
\begin{itemize}
\item \all{Meiner Meinung nach ist Jonas' Plan vernünftig, weil er zeigt, wie man mit wenig Geld ein Ziel erreicht.}
\item Ou : \all{Ich finde das Beispiel gut, weil viele Jugendliche nicht planen.}
\end{itemize}
\end{methode}

\begin{aretenir}[Nombre minimal de phrases : 5]
Au Bac, un commentaire réussi contient \textbf{au moins 5 phrases} :
\begin{enumerate}
\item 1 phrase d'introduction (type, thème, personnage).
\item 3 phrases de résumé (une par partie identifiée).
\item 1 phrase d'avis personnel avec \all{weil} (ou \all{da}).
\end{enumerate}
Toute phrase supplémentaire (analyse approfondie, comparaison avec le Cameroun) est un atout pour la note d'expression.
\end{aretenir}

\courssub{Application au texte « Lernen, mit Geld umzugehen »}

\begin{textebox}[Texte support — Lernen, mit Geld umzugehen]
Jonas ist 17 Jahre alt und geht in Yaoundé aufs Gymnasium. Jeden Monat bekommt er 15.000 CFA Taschengeld von seinen Eltern. Außerdem arbeitet er samstags manchmal in der Boutique seines Onkels und verdient dort 5.000 CFA pro Tag. Das sind seine Einnahmen.

Seine größten Ausgaben sind das Handy und die Kleidung. Vor drei Monaten hat er ein Smartphone für 95.000 CFA gekauft. Jeden Monat gibt er etwa 10.000 CFA für Kredit und Internet aus. Für Kleidung und Schuhe braucht er noch einmal 15.000 CFA im Monat. „Am Ende des Monats ist oft nichts mehr übrig“, sagt Jonas.

Aber Jonas hat ein Projekt: Er möchte nächstes Jahr im Sommer nach Österreich reisen, um seine Tante in Wien zu besuchen. Der Flug kostet ungefähr 450.000 CFA. Er hat gerechnet: Wenn er jeden Monat 20.000 CFA spart, hat er nach 22 Monaten genug Geld. „Das ist ein langer Plan“, sagt er, „aber ich will diszipliniert sein.“ Zuerst muss er seine Ausgaben reduzieren: weniger Kredit, keine neuen Schuhe jeden Monat. Dann legt er das Geld auf ein Sparkonto bei der Bank. Schließlich wird er sein Ziel erreichen.
\end{textebox}

\begin{definition}[Glossaire du texte]
\begin{tabular}{|l|l|l|}
\hline
\rowcolor{popblueL}
\textbf{Allemand} & \textbf{Français} & \textbf{Remarque} \\
\hline
das Taschengeld, -er & l'argent de poche & nom composé \\
der Onkel, - & l'oncle & pluriel identique \\
verdienen (verdient, verdiente, hat verdient) & gagner (de l'argent) & verbe \textbf{faible} \\
die Ausgabe, -n & la dépense & \\
die Einnahme, -n & le revenu, la recette & \\
übrig bleiben (bleibt übrig, blieb übrig, ist übrig geblieben) & rester (en surplus) & verbe à particule séparable ; auxiliaire \all{sein} \\
das Projekt, -e & le projet & \\
der Plan, Pläne & le plan & pluriel avec Umlaut \\
sparen (spart, sparte, hat gespart) & épargner & verbe faible \\
diszipliniert & discipliné & adjectif \\
reduzieren (reduziert, reduzierte, hat reduziert) & réduire & verbe faible \\
das Sparkonto, -en & le compte d'épargne & nom composé : \all{sparen} + \all{das Konto} \\
\hline
\end{tabular}
\end{definition}

\begin{exemplebox}[Identification de la structure en 3 parties]
\begin{tabular}{|p{2.5cm}|p{6cm}|p{5cm}|}
\hline
\rowcolor{popblueL}
\textbf{Partie} & \textbf{Contenu} & \textbf{Mots-clés / Connecteurs} \\
\hline
1. Einnahmen & Taschengeld (15.000 CFA), petit job chez l'oncle (5.000 CFA par jour). & \all{außerdem}, \all{das sind seine Einnahmen} \\
\hline
2. Ausgaben & Handy (95.000 CFA + 10.000 CFA/mois), vêtements (15.000 CFA/mois), citation « nichts mehr übrig ». & \all{seine größten Ausgaben}, \all{am Ende des Monats} \\
\hline
3. Sparprojekt & Voyage en Autriche (450.000 CFA), calcul : 20.000 CFA/mois pendant 22 mois, réduction des dépenses, Sparkonto. & \all{zuerst}, \all{dann}, \all{schließlich}, \all{wenn …, wird …} \\
\hline
\end{tabular}
\end{exemplebox}

\courssub{Commentaire modèle}

\begin{textebox}[Commentaire modèle — Version allemande (8–10 lignes)]
Der Text „Lernen, mit Geld umzugehen“ handelt von Jonas, einem 17-jährigen Schüler in Yaoundé, und seiner Budgetplanung. Zuerst beschreibt der Autor Jonas' Einnahmen: Taschengeld von den Eltern und ein kleiner Job beim Onkel. Dann zeigt er seine größten Ausgaben: das Handy mit monatlichen Kosten und die Kleidung. Danach erklärt er Jonas' Projekt: Wenn er jeden Monat 20.000 CFA spart, wird er nach 22 Monaten genug Geld für den Flug nach Wien haben. Der Text zeigt, dass Planung und Disziplin wichtig sind, um ein Ziel zu erreichen. Meiner Meinung nach ist Jonas ein gutes Vorbild, weil viele Jugendliche ihr Geld nicht planen und am Monatsende nichts mehr haben.
\end{textebox}

\begin{exemplebox}[Traduction guidée phrase par phrase]
\begin{tabular}{|p{7cm}|p{7cm}|}
\hline
\rowcolor{popblueL}
\textbf{Allemand} & \textbf{Français} \\
\hline
Der Text „Lernen, mit Geld umzugehen“ handelt von Jonas, einem 17-jährigen Schüler in Yaoundé, und seiner Budgetplanung. & Le texte « Apprendre à gérer l'argent » traite de Jonas, un élève de 17 ans à Yaoundé, et de sa planification budgétaire. \\
\hline
Zuerst beschreibt der Autor Jonas' Einnahmen: Taschengeld von den Eltern und ein kleiner Job beim Onkel. & D'abord, l'auteur décrit les revenus de Jonas : l'argent de poche donné par les parents et un petit job chez l'oncle. \\
\hline
Dann zeigt er seine größten Ausgaben: das Handy mit monatlichen Kosten und die Kleidung. & Ensuite, il montre ses plus grosses dépenses : le portable avec ses coûts mensuels et les vêtements. \\
\hline
Danach erklärt er Jonas' Projekt: Wenn er jeden Monat 20.000 CFA spart, wird er nach 22 Monaten genug Geld für den Flug nach Wien haben. & Puis, il explique le projet de Jonas : s'il épargne 20.000 CFA par mois, il aura assez d'argent pour le vol vers Vienne après 22 mois. \\
\hline
Der Text zeigt, dass Planung und Disziplin wichtig sind, um ein Ziel zu erreichen. & Le texte montre que la planification et la discipline sont importantes pour atteindre un objectif. \\
\hline
Meiner Meinung nach ist Jonas ein gutes Vorbild, weil viele Jugendliche ihr Geld nicht planen und am Monatsende nichts mehr haben. & À mon avis, Jonas est un bon modèle, parce que beaucoup de jeunes ne planifient pas leur argent et n'ont plus rien à la fin du mois. \\
\hline
\end{tabular}
\end{exemplebox}

\begin{attention}[Le génitif des noms propres en -s]
Les noms propres terminés par \textbf{-s}, \textbf{-z}, \textbf{-x} forment le génitif avec une apostrophe, \textbf{sans s supplémentaire} : \all{Jonas' Einnahmen}, \all{Jonas' Plan}, \all{Franz' Buch}. La forme \all{Jonases} existe dans la langue familière, mais l'apostrophe est la norme à l'écrit. À l'oral, on prononce « Jonas-es ».
\end{attention}

\courssub{Banque de phrases toutes faites pour le commentaire}

\begin{propriete}[Phrases-types : Introduction, Résumé, Analyse, Opinion]
\begin{center}
\begin{tabularx}{\linewidth}{|X|X|}
\hline
\rowcolor{popblueL}
\textbf{Fonction} & \textbf{Formule allemande (à adapter)} \\
\hline
Présenter le texte & \all{Der Text „…“ handelt von …} \\
& \all{In dem Artikel „…“ geht es um …} \\
& \all{Der Autor beschreibt / zeigt / erklärt …} \\
\hline
Structurer le résumé & \all{Zuerst … / Zunächst …} \\
& \all{Dann … / Anschließend …} \\
& \all{Danach … / Im nächsten Teil …} \\
& \all{Schließlich … / Am Ende …} \\
\hline
Analyser le message & \all{Der Autor zeigt, dass …} \\
& \all{Das Beispiel von … verdeutlicht, dass …} \\
& \all{Wichtig ist die Botschaft: …} \\
& \all{Mit Zahlen und einem konkreten Beispiel macht der Autor klar, dass …} \\
\hline
Exprimer son avis & \all{Meiner Meinung nach …} \\
& \all{Ich finde, dass …} \\
& \all{Ich denke, … ist richtig / wichtig, weil …} \\
& \all{Meines Erachtens …} \\
\hline
Justifier avec \all{weil} & \all{…, weil …} (verbe conjugué à la fin) \\
& \all{…, da …} (verbe conjugué à la fin ; \all{da} en tête de phrase) \\
\hline
Conclure & \all{Zusammenfassend kann man sagen: …} \\
& \all{Der Text regt zum Nachdenken an.} \\
\hline
\end{tabularx}
\end{center}
\end{propriete}

\begin{attention}[Pièges fréquents des francophones]
\begin{enumerate}
\item \textbf{Résumé = copier-coller} : non. Il faut reformuler au présent de narration (\all{Präsens}) avec ses propres mots. Ne citer que 3–4 mots au maximum, entre guillemets allemands \all{„…“}.
\item \textbf{\all{Nach meiner Meinung} seul} : à éviter. La forme standard est \all{meiner Meinung nach} (génitif + \all{nach} postposé) ou \all{ich finde, dass …}. \all{Nach meiner Meinung} existe mais sonne moins naturel ; évitez-le au Bac.
\item \textbf{Oublier le verbe à la fin après \all{weil}} : \textcolor{popdark}{Faux :} \all{…, weil er ist diszipliniert.} \textcolor{popgreen}{Correct :} \all{…, weil er diszipliniert ist.}
\item \textbf{Confondre \all{dann} et \all{denn}} : \all{dann} = « ensuite, alors » (marqueur temporel, verbe en 2\textsuperscript{e} position : \all{Dann zeigt er …}) ; \all{denn} = « car » (coordination, verbe en 2\textsuperscript{e} position : \all{Er spart, denn er hat ein Ziel.}). Ne pas confondre \all{denn} avec \all{weil}, qui renvoie le verbe à la fin.
\item \textbf{Majuscules à tous les noms} : \all{das Budget, die Ausgaben, die Einnahmen, der Plan, das Projekt, das Geld, das Handy, die Kleidung, das Sparkonto, der Flug, die Tante, der Sommer, der Monat, das Jahr, das Ziel, die Disziplin, das Vorbild, das Monatsende}.
\item \textbf{Faux ami} : \all{das Kredit} est incorrect ; on dit \all{der Kredit} (le crédit). Pour le « crédit téléphonique », on utilise aussi \all{das Guthaben} ou \all{die Telefonkarte}.
\end{enumerate}
\end{attention}

\begin{popfigure}
\begin{tikzpicture}[
  node distance=0.55cm,
  box/.style={rectangle, rounded corners, draw=popblue, thick, fill=popblueL, text width=10.5cm, align=center, minimum height=1.1cm, font=\small},
  arrow/.style={-Stealth, thick, popdark}
]
\node[box] (einl) {\textbf{Einleitung} — type de texte, thème, personnage (1 phrase)};
\node[box, below=of einl, align=center] (zuss){\textbf{Zusammenfassung} — 3 parties : revenus, dépenses, projet\\Connecteurs : \all{zuerst – dann – danach – schließlich}};
\node[box, below=of zuss] (anal) {\textbf{Analyse} — message central + procédés : chiffres, citation, exemple concret};
\node[box, below=of anal] (mein) {\textbf{Eigene Meinung} — \all{Meiner Meinung nach …, weil …}};
\node[box, below=of mein] (schl) {\textbf{Schluss} — bilan + ouverture : \all{Zusammenfassend kann man sagen: …}};

\draw[arrow] (einl) -- (zuss);
\draw[arrow] (zuss) -- (anal);
\draw[arrow] (anal) -- (mein);
\draw[arrow] (mein) -- (schl);
\end{tikzpicture}
\legende{Organigramme de la structure du commentaire de texte}
\end{popfigure}

\courssub{Exercice résolu : rédiger l'introduction et le résumé}

\begin{exoresolu}[Rédiger l'introduction et les 3 phrases de résumé du texte « Lernen, mit Geld umzugehen »]
\textbf{Consigne :} en vous appuyant sur la méthode, écrivez en allemand : (1) la phrase d'introduction, (2) trois phrases de résumé (une par partie), en utilisant \all{zuerst}, \all{dann}, \all{schließlich}.

\tcblower
\textbf{Solution détaillée :}

\noindent\textbf{1. Introduction :}
\all{Der Text „Lernen, mit Geld umzugehen“ handelt von Jonas, einem 17-jährigen Gymnasiasten in Yaoundé, und seiner Budgetplanung.}

\noindent\textbf{2. Résumé (3 phrases) :}
\begin{itemize}
\item \all{Zuerst beschreibt der Autor Jonas' Einnahmen: Er erhält Taschengeld von seinen Eltern und verdient zusätzlich Geld bei seinem Onkel.}
\item \all{Dann stellt er die größten Ausgaben vor: das teure Handy mit monatlichen Kosten und die Kleidung, sodass am Monatsende oft nichts übrig bleibt.}
\item \all{Schließlich erklärt er Jonas' Sparprojekt: Wenn er monatlich 20.000 CFA zurücklegt, kann er nach 22 Monaten seine Tante in Wien besuchen.}
\end{itemize}

\noindent\textbf{Analyse de la correction :}
\begin{itemize}
\item Introduction : type implicite (texte), thème (\all{Budgetplanung}), personnage (Jonas, 17 ans, Yaoundé) — \textcolor{popgreen}{correct}.
\item Phrase 1 (revenus) : \all{zuerst} + reformulation (\all{Taschengeld}, \all{verdient zusätzlich Geld}) — \textcolor{popgreen}{correct}.
\item Phrase 2 (dépenses) : \all{dann} + synthèse (\all{das teure Handy}, \all{die Kleidung}, \all{nichts übrig bleibt} — particule \all{übrig} bien séparée et renvoyée à la fin) — \textcolor{popgreen}{correct}.
\item Phrase 3 (projet) : \all{schließlich} + subordonnée conditionnelle \all{wenn …, kann er …} (verbe à la fin dans la subordonnée, inversion dans la principale) + but (\all{seine Tante in Wien besuchen}) — \textcolor{popgreen}{correct}.
\item Vocabulaire clé réemployé : \all{die Einnahmen, die Ausgaben, das Sparprojekt, zurücklegen, monatlich}.
\end{itemize}
\end{exoresolu}

\begin{experience}[Atelier commentaire — Production orale guidée]
\textbf{Consigne :} en binôme, l'un joue le rôle de l'examinateur, l'autre celui du candidat.
\begin{enumerate}
\item Le candidat présente oralement le commentaire complet (5 phrases minimum) en 1 min 30.
\item L'examinateur note sur une grille : Introduction (type/thème/personnage) — Résumé en 3 parties (connecteurs) — Analyse (message/procédés) — Avis personnel (\all{weil} + verbe final) — Prononciation et fluidité.
\item Échange des rôles. Débriefing collectif : quelles phrases-types ont aidé ? Le verbe est-il bien passé à la fin après \all{weil} ?
\end{enumerate}
\textbf{Variante :} enregistrer sa production sur smartphone (application OnBuch+) et s'auto-évaluer avec la grille.
\end{experience}

\begin{savaistu}[Le commentaire au Bac camerounais et en Allemagne]
Au Cameroun (LV2 allemand), le commentaire de texte écrit est une tâche courte et guidée : il évalue à la fois la compréhension et l'expression. En Allemagne, l'équivalent au \all{Gymnasium} (Oberstufe) s'appelle \all{die Textanalyse} ou \all{die Erörterung} : on attend une introduction formelle avec une thèse, une analyse structurée par arguments (thématiques, linguistiques, stylistiques) et une conclusion critique. Notre méthode « 5 phrases minimum » est une adaptation pragmatique au niveau visé en Terminale : elle garantit une structure claire sans exiger l'analyse stylistique approfondie (figures de style, métaphores) réservée aux niveaux supérieurs.
\end{savaistu}

\begin{exercice}[Entraînement : à votre tour]
Rédigez un commentaire complet (5–7 phrases) sur le texte \all{Lernen, mit Geld umzugehen} en respectant la méthode.
\begin{enumerate}
\item Introduction (1 phrase).
\item Résumé en 3 phrases avec \all{zuerst – dann – schließlich}.
\item Analyse du message (1 phrase : \all{Der Autor zeigt, dass …}).
\item Avis personnel avec \all{Meiner Meinung nach …, weil …}.
\end{enumerate}
Utilisez la banque de phrases et le glossaire. Écrivez sur votre cahier, puis comparez avec le corrigé.
\end{exercice}

\begin{corrige}[Exercice — Commentaire complet]
\all{Der Text „Lernen, mit Geld umzugehen“ handelt von Jonas, einem 17-jährigen Schüler in Yaoundé, und seiner Budgetplanung. Zuerst beschreibt der Autor Jonas' Einnahmen: Taschengeld und ein Nebenjob beim Onkel. Dann zeigt er die hohen Ausgaben für das Handy und die Kleidung. Schließlich erklärt er Jonas' Sparplan: Wenn er monatlich 20.000 CFA spart, wird er nach 22 Monaten nach Wien fliegen können. Der Autor zeigt, dass Disziplin und Planung nötig sind, um ein Ziel zu erreichen. Meiner Meinung nach ist der Plan realistisch, weil Jonas konkrete Zahlen nennt und seine Ausgaben reduzieren will.}

\medskip
\textbf{Points de méthode vérifiés :} introduction complète ; 3 connecteurs de résumé ; subordonnée \all{wenn} avec verbe final et Futur I dans la principale (\all{wird … fliegen können}) ; phrase d'analyse avec \all{dass} ; avis personnel justifié avec \all{weil} + verbe à la fin (\all{nennt}, \all{will}).
\end{corrige}

\coursec{Production guidée : Mein Budget, mein Projekt}

Après avoir appris, dans la section précédente, à lire et à commenter un texte sur la gestion de l'argent, il est temps de passer à l'écriture : c'est à vous de parler de \emph{votre} budget et de \emph{votre} projet. Cette production guidée est l'aboutissement de la leçon : elle réutilise le lexique (\cle{das Budget}, \cle{die Einnahmen}, \cle{die Ausgaben}, \cle{sparen}, \cle{das Projekt}), le comparatif, le superlatif, la subordonnée avec \all{wenn} et le \cle{Futur I}. Suivez la trame pas à pas : une production réussie est une production simple, correcte et complète.

\courssub{La tâche : comprendre la consigne}

\begin{definition}[La tâche écrite]
Rédiger un paragraphe de \textbf{6 à 8 phrases} intitulé \all{„Mein Budget und mein Projekt“} (« Mon budget et mon projet ») en suivant une trame imposée, phrase par phrase. Le texte doit être personnel : parlez de votre argent de poche (\all{das Taschengeld}), de vos dépenses réelles et d'un projet que vous rêvez de financer.
\end{definition}

\begin{propriete}[La trame imposée et les contraintes linguistiques]
\begin{center}
\begin{tabularx}{\linewidth}{|l|X|l|}
\hline
\rowcolor{popblueL} \textbf{Phrase} & \textbf{Contenu attendu} & \textbf{Structure visée} \\
\hline
1 & Mes revenus : \all{Taschengeld}, \all{Geschenke} (cadeaux) & phrase simple \\
\hline
2 & Mes dépenses & \textbf{comparatif} \\
\hline
3 & Mon projet & phrase simple \\
\hline
4 & Une condition & subordonnée avec \all{wenn} \\
\hline
5 & Une annonce & \textbf{Futur I} \\
\hline
6 & Conclusion & \textbf{superlatif} \\
\hline
\end{tabularx}
\end{center}
Contraintes obligatoires : au moins \textbf{1 comparatif}, \textbf{1 superlatif}, \textbf{1 subordonnée avec} \all{wenn}, \textbf{1 verbe au Futur I} et \textbf{5 mots du lexique} de la leçon.
\end{propriete}

\courssub{Un modèle réussi : observer avant d'écrire}

Lisez d'abord ce texte modèle, rédigé par un élève de Terminale A de Yaoundé. Repérez les structures imposées : elles sont le squelette du texte.

\begin{textebox}[Modèle de production : „Mein Budget und mein Projekt“]
Mein Budget und mein Projekt

Jeden Monat bekomme ich 10 000 FCFA Taschengeld von meinen Eltern, und manchmal bekomme ich auch Geschenke von meiner Tante in Douala. Ich gebe mehr Geld für Transport und Essen aus als für Handyguthaben, denn die Schule ist weit von unserem Haus. Mein großes Projekt ist ein neues Fahrrad: es kostet 45 000 FCFA. Wenn ich jeden Monat 5 000 FCFA spare, werde ich das Fahrrad in neun Monaten kaufen. Ich werde weniger Snacks kaufen und mein Geld in einer kleinen Kasse zu Hause behalten. Das Fahrrad wird mein wichtigster Besitz sein, und das Sparen ist die beste Schule für das Leben.

\emph{Glossaire :} das Taschengeld (l'argent de poche, pas de pluriel usuel) ; das Geschenk, -e (le cadeau) ; ausgeben (dépenser) ; das Handyguthaben (le crédit téléphonique, pl. die Handyguthaben) ; das Fahrrad, -er (le vélo) ; sparen (économiser) ; die Kasse, -n (la caisse, la tirelire) ; behalten (garder) ; der Besitz (la possession, pl. die Besitze).
\end{textebox}

Analysons le modèle :

\begin{itemize}
\item \textbf{Phrase 1} (revenus) : \all{Jeden Monat bekomme ich 10 000 FCFA Taschengeld...} — verbe conjugué \all{bekomme} en position 2, après le complément de temps en tête.
\item \textbf{Phrase 2} (comparatif) : \all{Ich gebe mehr Geld für Transport und Essen aus als für Handyguthaben} — comparatif avec \all{mehr ... als} (« plus ... que »).
\item \textbf{Phrase 3} (projet) : \all{Mein großes Projekt ist ein neues Fahrrad}.
\item \textbf{Phrase 4} (condition) : \all{Wenn ich jeden Monat 5 000 FCFA spare, ...} — verbe \all{spare} en \textbf{fin} de subordonnée.
\item \textbf{Phrase 5} (Futur I) : \all{... werde ich das Fahrrad in neun Monaten kaufen} — \all{werden} conjugué en position 2 + infinitif \all{kaufen} en fin de phrase.
\item \textbf{Phrase 6} (superlatif) : \all{mein wichtigster Besitz}, \all{die beste Schule} — superlatifs avec accord et article.
\end{itemize}

\begin{exemplebox}[Modèles de phrases à réutiliser]
\all{„Wenn ich jeden Monat 5000 FCFA spare, werde ich mir ein neues Handy kaufen.“} = « Si j'économise 5000 FCFA par mois, je m'achèterai un nouveau téléphone. »

\all{„Meine größte Ausgabe ist das Schulgeld.“} = « Ma plus grande dépense, ce sont les frais de scolarité. »

\all{„Ich gebe weniger Geld für Spiele aus als mein Bruder.“} = « Je dépense moins d'argent pour les jeux que mon frère. »

\all{„Das Sparen ist am wichtigsten für mein Projekt.“} = « L'épargne est ce qu'il y a de plus important pour mon projet. »
\end{exemplebox}

\courssub{Étape 1 : préparation orale en binôme}

Avant d'écrire, on parle. L'interview en binôme permet de rassembler les idées et le vocabulaire : les réponses orales deviendront les phrases écrites.

\begin{experience}[Interview : „Mein Geld und ich“]
Par groupes de deux, posez-vous les questions suivantes et notez les réponses de votre partenaire. Changez ensuite les rôles.
\begin{enumerate}
\item \all{Wie viel Geld bekommst du jeden Monat?} (Combien d'argent reçois-tu chaque mois ?)
\item \all{Wofür gibst du Geld aus?} (Pour quoi dépenses-tu de l'argent ?)
\item \all{Was ist dein größter Wunsch?} (Quel est ton plus grand souhait ?)
\item \all{Was wirst du machen, wenn du sparst?} (Que feras-tu si tu économises ?)
\end{enumerate}
Exemple de réponse attendue : \all{Ich bekomme 5 000 FCFA. Ich gebe Geld für Essen und Transport aus. Wenn ich spare, werde ich mir Schuhe kaufen.}
\end{experience}

\courssub{Étape 2 : planifier avec le tableau-budget}

Un bon texte commence par un bon plan. Complétez ce tableau de planification avant de rédiger : chaque case vous donnera une phrase.

\begin{popfigure}
\begin{center}
\begin{tikzpicture}
\draw[popblue, thick] (0,0) rectangle (12,4.5);
\draw[popblue, thick] (0,3.4) -- (12,3.4);
\draw[popblue, thick] (3,0) -- (3,4.5);
\draw[popblue, thick] (6,0) -- (6,4.5);
\draw[popblue, thick] (9,0) -- (9,4.5);
\fill[popblueL] (0,3.4) rectangle (3,4.5);
\fill[poporangeL] (3,3.4) rectangle (6,4.5);
\fill[popgreenL] (6,3.4) rectangle (9,4.5);
\fill[poppurpleL] (9,3.4) rectangle (12,4.5);
\node[align=center] at (1.5,3.95) {\textbf{Einnahmen}\\ (revenus)};
\node[align=center] at (4.5,3.95) {\textbf{Ausgaben}\\ (dépenses)};
\node[align=center] at (7.5,3.95) {\textbf{Sparziel}\\ (objectif d'épargne)};
\node[align=center] at (10.5,3.95) {\textbf{Projekt}\\ (projet)};
\node[align=center, text=popmuted] at (1.5,1.7) {\all{Taschengeld:}\\ \all{...... FCFA}\\ \all{Geschenke:}\\ \all{......}};
\node[align=center, text=popmuted] at (4.5,1.7) {\all{Essen: ......}\\ \all{Transport:}\\ \all{......}\\ \all{mehr / weniger}\\ \all{als ...}};
\node[align=center, text=popmuted] at (7.5,1.7) {\all{Ich spare}\\ \all{...... FCFA}\\ \all{pro Monat.}};
\node[align=center, text=popmuted] at (10.5,1.7) {\all{Ich kaufe:}\\ \all{......}\\ \all{Preis:}\\ \all{...... FCFA}};
\end{tikzpicture}
\end{center}
\legende{Tableau de planification à compléter : vos réponses deviennent vos phrases.}
\end{popfigure}

\courssub{Étape 3 : la méthode de rédaction}

\begin{methode}[Réussir sa production guidée en quatre temps]
\begin{enumerate}
\item \textbf{Écrire d'abord des phrases simples et correctes.} Pour chaque ligne de la trame, rédigez une phrase de base au présent, avec le verbe en position 2 : \all{Ich bekomme Taschengeld. Ich gebe Geld aus. Ich habe ein Projekt.}
\item \textbf{Ajouter les structures imposées une par une.} Transformez la phrase 2 avec un comparatif (\all{mehr ... als}), la phrase 4 avec \all{wenn} (verbe en fin !), la phrase 5 avec \all{werden} + infinitif, la phrase 6 avec un superlatif (\all{der/die/das ...-ste} ou \all{am ...-sten}).
\item \textbf{Insérer le lexique de la leçon.} Vérifiez que 5 mots du thème apparaissent : \all{das Budget}, \all{die Einnahmen}, \all{die Ausgaben}, \all{sparen}, \all{das Geld}, \all{der Plan}, \all{das Projekt}.
\item \textbf{Relire avec la grille de correction} (voir ci-dessous), puis échanger sa copie avec un voisin pour la relecture croisée.
\end{enumerate}
\end{methode}

\begin{propriete}[Grille de relecture croisée]
\begin{center}
\begin{tabularx}{\linewidth}{|X|l|}
\hline
\rowcolor{popblueL} \textbf{Question de contrôle} & \textbf{Oui / Non} \\
\hline
Le verbe conjugué est-il en position 2 dans chaque phrase principale ? & \\
\hline
Le verbe est-il à la fin après \all{wenn} ? & \\
\hline
L'infinitif est-il à la fin avec \all{werden} (Futur I) ? & \\
\hline
Le superlatif est-il correctement accordé (\all{mein größter Wunsch}) ? & \\
\hline
Les noms ont-ils la majuscule, le bon article et le bon pluriel ? & \\
\hline
Y a-t-il au moins 5 mots du lexique de la leçon ? & \\
\hline
\end{tabularx}
\end{center}
\end{propriete}

\begin{attention}[Les quatre pièges à vérifier à la relecture]
\begin{enumerate}
\item \textbf{Verbe conjugué en position 2.} Après un complément en tête, le sujet se place après le verbe : \all{Jeden Monat \textbf{bekomme ich}...} et non \all{*Jeden Monat ich bekomme...}.
\item \textbf{Verbe en fin après \all{wenn}.} \all{Wenn ich jeden Monat 5000 FCFA \textbf{spare}, ...} — jamais \all{*wenn ich spare jeden Monat 5000 FCFA} sur le modèle français « si j'économise chaque mois... ».
\item \textbf{Infinitif en fin avec \all{werden}.} \all{Ich \textbf{werde} ein Fahrrad \textbf{kaufen}.} Ne dites jamais \all{*ich werde kaufen ein Fahrrad}.
\item \textbf{Article et pluriel des noms du lexique.} \all{die Ausgabe, -n} ; \all{die Einnahme, -n} ; \all{das Geschenk, -e} ; \all{das Projekt, -e}. Et n'oubliez pas la \textbf{majuscule} à tous les noms !
\end{enumerate}
\end{attention}

\courssub{Exercice résolu : corriger une production}

\begin{exoresolu}[Trouve et corrige les quatre erreurs]
Voici la production de Marie, élève de Tle A. Elle contient quatre erreurs (une par phrase). Retrouvez-les et corrigez-les.

\all{1. Jeden Monat ich bekomme 8000 FCFA von meinen Eltern. 2. Ich gebe mehr Geld für Essen aus als für Kleider. 3. Wenn ich spare jeden Monat 3000 FCFA, ich werde ein Handy kaufen. 4. Das Handy wird mein wichtigst Besitz.}

\tcblower
\textbf{Solution détaillée :}
\begin{enumerate}
\item \all{Jeden Monat \textbf{bekomme ich} 8000 FCFA von meinen Eltern.} — Erreur d'ordre des mots : le complément de temps en tête occupe la position 1, le verbe conjugué reste en position 2, le sujet suit (inversion).
\item Phrase 2 \textbf{correcte} : comparatif \all{mehr ... als} bien formé, particule \all{aus} du verbe \all{ausgeben} bien placée en fin.
\item \all{Wenn ich jeden Monat 3000 FCFA \textbf{spare}, \textbf{werde} ich ein Handy \textbf{kaufen}.} — Double erreur : le verbe de la subordonnée doit être en fin (\all{spare}), et la principale qui suit commence par le verbe conjugué (\all{werde}), pas par \all{ich}.
\item \all{Das Handy wird mein \textbf{wichtigster} Besitz.} — Superlatif attribut : il prend la terminaison adjectivale après le possessif \all{mein} : \all{wichtigst\textbf{er} Besitz} (masculin nominatif).
\end{enumerate}
\end{exoresolu}

\courssub{Étape 4 : la version orale avec affiche-budget}

L'écrit devient parole : chaque élève présente son projet devant la classe en \textbf{5 phrases}, en s'appuyant sur une affiche-budget simple (un grand tableau \all{Einnahmen / Ausgaben / Sparziel / Projekt} dessiné au tableau ou sur une feuille).

\begin{experience}[Présentation orale : „Mein Projekt in fünf Sätzen“]
\begin{enumerate}
\item Préparez votre affiche-budget : quatre cases, des chiffres, un dessin de votre projet.
\item Présentez-vous en 5 phrases : revenus, dépenses (comparatif), projet, condition avec \all{wenn}, annonce au Futur I.
\item La classe écoute et note : chaque auditeur doit poser \textbf{une question} au présentateur, par exemple \all{Wie lange musst du sparen?} (Combien de temps dois-tu économiser ?) ou \all{Warum ist dieses Projekt wichtig für dich?} (Pourquoi ce projet est-il important pour toi ?).
\item Répondez spontanément : c'est l'occasion de réutiliser \all{wenn} à l'oral : \all{Wenn ich nicht genug spare, kaufe ich ein billigeres Handy.}
\end{enumerate}
Critères d'évaluation orale : prononciation claire, verbes bien placés, regard vers le public, affiche lisible.
\end{experience}

\begin{exercice}[À vous d'écrire]
Rédigez maintenant votre propre texte \all{„Mein Budget und mein Projekt“} (6 à 8 phrases) en suivant la trame imposée. Utilisez le tableau de planification, la méthode en quatre temps et la grille de relecture. Échangez ensuite votre copie avec votre voisin : cochez sa grille et signalez-lui gentiment ses erreurs.
\end{exercice}

\begin{corrige}[Exercice : éléments d'appréciation]
Il n'existe pas une seule bonne réponse, mais tout texte réussi contient : (1) une phrase de revenus avec \all{bekommen} et \all{das Taschengeld} ; (2) un comparatif correct (\all{mehr/weniger ... als} ou \all{besser/größer als}) ; (3) une phrase de projet claire ; (4) une subordonnée \all{wenn} avec le verbe en fin ; (5) un Futur I avec \all{werden} + infinitif final ; (6) un superlatif accordé (\all{mein größter Wunsch}, \all{die beste Idee}, \all{am wichtigsten}) ; (7) au moins 5 mots du lexique avec articles et majuscules corrects.
\end{corrige}

\begin{savaistu}[La production écrite au Bac allemand Tle A]
Au baccalauréat, l'épreuve d'allemand comprend toujours une production écrite. Elle est notée sur trois critères : la \textbf{correction grammaticale} (ordre des mots, conjugaisons, déclinaisons), le \textbf{respect du sujet} (répondre exactement à la consigne, comme ici à la trame imposée) et la \textbf{richesse du vocabulaire}. Bonne nouvelle : réutiliser les mots de la leçon (\all{das Budget}, \all{sparen}, \all{die Ausgaben}...) rapporte des points ! Les correcteurs récompensent un texte simple et sans faute bien plus qu'un texte ambitieux et incorrect. La règle d'or : \all{„Lieber einfach und richtig als kompliziert und falsch!“} — « Plutôt simple et correct que compliqué et faux ! »
\end{savaistu}

\begin{aretenir}[L'essentiel de la production guidée]
\begin{itemize}
\item Suivre la trame phrase par phrase : revenus, dépenses (comparatif), projet, condition (\all{wenn}), annonce (Futur I), conclusion (superlatif).
\item Écrire simple d'abord, enrichir ensuite : les structures imposées s'ajoutent une par une.
\item Relire avec la grille : verbe en position 2, verbe final après \all{wenn}, infinitif final avec \all{werden}, articles et majuscules des noms.
\item Préparer l'oral par l'interview en binôme et présenter son projet avec une affiche-budget.
\end{itemize}
\end{aretenir}

\newpage

\coursec{Activité d'intégration}

\courssub{Objectifs de l'activité}

À l'issue de cette activité d'intégration, l'élève sera capable de :

\begin{itemize}
\item mobiliser dans une situation de communication authentique le \cle{lexique du budget et des projets} (\all{das Budget, die Einnahmen, die Ausgaben, sparen, das Projekt, der Plan, das Geld}) ;
\item comparer des options à l'aide du \cle{comparatif} et du \cle{superlatif} (\all{billiger als, die beste Lösung, am nützlichsten}) ;
\item formuler des hypothèses et des conditions avec la \cle{subordonnée conditionnelle} introduite par \all{wenn} (ordre des mots : verbe à la fin) ;
\item annoncer des actions à venir au \cle{Futur I} (\all{wir werden sammeln, ich werde organisieren}) ;
\item travailler en groupe, présenter un projet à l'oral et rédiger individuellement un court paragraphe de synthèse structuré.
\end{itemize}

\courssub{Situation : La note du trésorier}

La fin de l'année scolaire approche au lycée de Yaoundé. Le bureau des élèves de la classe de Terminale A souhaite organiser un \cle{projet de classe} avant les examens. Trois idées sont en discussion : une fête de fin d'année dans la salle des fêtes du lycée, une excursion d'une journée à Kribi, ou la création d'un club d'allemand avec l'achat de livres et de jeux. Le problème : l'argent est limité. Le trésorier de la classe a rédigé une note d'information.

\begin{textebox}[Note du trésorier de la classe]
Liebe Mitschülerinnen und Mitschüler,

unsere Klasse hat zurzeit 120\,000 FCFA auf dem Konto. Wir können noch Geld sammeln: Jeder Schüler kann 2\,000 FCFA beitragen, dann haben wir ungefähr 180\,000 FCFA. Wir müssen jetzt entscheiden: Was ist das beste Projekt für unsere Klasse?

\textbf{Option A — Die Abschlussfeier:} Die Feier ist billiger als der Ausflug. Wir brauchen Essen, Getränke und Musik. Kosten: circa 100\,000 FCFA. Mit dem Rest können wir Geld für nächstes Jahr sparen.

\textbf{Option B — Der Ausflug nach Kribi:} Der Ausflug ist interessanter als die Feier, aber er ist auch teurer: Bus, Essen und Eintritt kosten circa 170\,000 FCFA. Wenn wir 180\,000 FCFA sammeln, werden wir fast kein Geld mehr haben.

\textbf{Option C — Der Deutschclub:} Das ist die billigste Option: Bücher, Spiele und Plakate kosten nur 60\,000 FCFA. Das Projekt ist am nützlichsten für unsere Zukunft, aber vielleicht nicht am lustigsten.

Bitte diskutiert in Gruppen und präsentiert euren Plan!

Euer Schatzmeister
\end{textebox}

\textbf{Glossaire :} \all{der Schatzmeister, die Schatzmeister} (le trésorier) ; \all{das Konto, die Konten} (le compte) ; \all{beitragen} (contribuer, donner sa part) ; \all{der Ausflug, die Ausflüge} (l'excursion) ; \all{der Eintritt, -e} (le droit d'entrée) ; \all{nützlich} (utile) ; \all{lustig} (drôle, amusant).

\courssub{Compréhension du texte (Étape 1)}

\begin{exercice}[Questions de compréhension globale et détaillée]
Répondez en allemand par des phrases complètes.
\begin{enumerate}
\item Wie viel Geld hat die Klasse momentan auf dem Konto?
\item Welche drei Optionen schlägt der Schatzmeister vor?
\item Wie viel kostet Option A und wie viel Option B?
\item Finden Sie im Text zwei Komparative und zwei Superlative. Schreiben Sie sie auf.
\item Was ist der Vorteil von Option C laut Text? Was ist der Nachteil?
\end{enumerate}
\end{exercice}

\begin{corrige}[Exercice 1 — Compréhension]
\begin{enumerate}
\item \all{Die Klasse hat zurzeit 120\,000 FCFA auf dem Konto.}
\item \all{Option A ist die Abschlussfeier, Option B ist der Ausflug nach Kribi, Option C ist der Deutschclub.}
\item \all{Option A kostet circa 100\,000 FCFA, Option B kostet circa 170\,000 FCFA.}
\item \all{Komparative: billiger als, interessanter als, teurer als. Superlative: das beste Projekt, die billigste Option, am nützlichsten, am lustigsten.}
\item \all{Vorteil: Das Projekt ist am nützlichsten für die Zukunft. Nachteil: Vielleicht nicht am lustigsten.}
\end{enumerate}
\end{corrige}

\courssub{Lexique thématique : Geld, Budget und Projekte}

\begin{definition}[Vocabulaire de base — Noms, verbes, adjectifs]
\begin{center}
\begin{tabularx}{\linewidth}{|X|X|X|}
\hline
\rowcolor{popblueL} \textbf{Français} & \textbf{Deutsch (Artikel, Pluriel)} & \textbf{Remarque / Famille de mots} \\
\hline
le budget & \all{das Budget, die Budgets} & \all{budgetieren} (budgétiser) \\
la recette, l'entrée d'argent & \all{die Einnahme, -n} & \all{einnehmen} (encaisser) \\
la dépense, la sortie d'argent & \all{die Ausgabe, -n} & \all{ausgeben} (dépenser, verbe à particule séparable) \\
économiser, l'épargne & \all{sparen} ; \all{das Sparziel, -e} & \all{der Sparer, -} ; \all{die Spardose, -n} (tirelire) \\
coûter & \all{kosten} & \all{etw. kostet 100 FCFA} (accusatif) \\
le projet & \all{das Projekt, -e} & \all{planen, die Planung} \\
le plan & \all{der Plan, die Pläne} & \all{planen} \\
l'argent & \all{das Geld} (pas de pluriel usuel) & \all{geldwert} (qui a de la valeur) \\
collecter, réunir & \all{sammeln} & \all{die Sammlung, -en} \\
le trésorier & \all{der Schatzmeister, die Schatzmeister} & \all{die Schatzmeisterin, -nen} (fém.) \\
le compte (bancaire) & \all{das Konto, die Konten} & \all{kontoführend} \\
l'excursion & \all{der Ausflug, die Ausflüge} & \all{ausfliegen} (rare) \\
le droit d'entrée & \all{der Eintritt, -e} & \all{der Eintrittspreis, -e} \\
utile & \all{nützlich} & \all{der Nutzen} ; \all{nutzlos} (inutile) \\
cher / bon marché & \all{teuer / billig} & \all{preiswert} (bon rapport qualité/prix) \\
\hline
\end{tabularx}
\end{center}
\end{definition}

\begin{popfigure}
\begin{tikzpicture}[mindmap, grow cyclic, every node/.style=concept, concept color=popblue!60,
    level 1/.append style={level distance=3.5cm, sibling angle=90, concept color=popblue!40},
    level 2/.append style={level distance=3cm, sibling angle=45, concept color=poporange!40},
    text=popdark, font=\small]
\node {Budget \& Projekte}
    child { node {Einnahmen} child { node {sammeln} } child { node {beitragen} } child { node {Konto} } }
    child { node {Ausgaben} child { node {kosten} } child { node {ausgeben} } child { node {Eintritt} } }
    child { node {Planung} child { node {Projekt} } child { node {Plan} } child { node {Sparziel} } }
    child { node {Entscheidung} child { node {vergleichen} } child { node {wählen} } child { node {abstimmen} } };
\end{tikzpicture}
\legende{Carte mentale : le vocabulaire du budget et des projets}
\end{popfigure}

\courssub{Grammaire 1 : Le comparatif et le superlatif}

\begin{propriete}[Formation et emploi du comparatif et du superlatif]
\textbf{1. Comparatif d'inégalité (Comparativ)} \\
Structure : \cle{Adjectif/Adverbe + -er + als} (que). L'adjectif ne se décline pas s'il est attribut (après \all{sein}, \all{werden}, \all{bleiben}) ; il se décline s'il est épithète (devant le nom).
\all{Option A ist billiger als Option B.} (attribut, invariable) \\
\all{Eine billigere Feier als die Reise.} (épithète, décliné : \all{billiger-e} fém. nominatif/accusatif)

\textbf{2. Superlatif (Superlativ)} \\
Deux formes selon la fonction :
\begin{itemize}
\item \textbf{Prédicatif (après verbe d'état) :} \cle{am + Adjectif + -sten} (souvent avec Umlaut pour monosyllabes).
\all{Option C ist am billigsten.} \all{Das Projekt ist am nützlichsten.}
\item \textbf{Attributif (devant nom) :} \cle{Article défini + Adjectif + -ste + déclinaison}.
\all{Das ist die billigste Option.} \all{Der teuerste Ausflug.}
\end{itemize}

\textbf{3. Formes irrégulières importantes} (à connaître par cœur) :
\begin{center}
\begin{tabular}{|l|l|l|l|}
\hline
\rowcolor{popblueL} \textbf{Positif} & \textbf{Comparatif} & \textbf{Superlatif (préd.)} & \textbf{Superlatif (attr.)} \\
\hline
\all{gut} & \all{besser} & \all{am besten} & \all{die beste} \\
\all{viel} & \all{mehr} & \all{am meisten} & \all{die meisten} \\
\all{gerne / gern} & \all{lieber} & \all{am liebsten} & \all{—} \\
\hline
\end{tabular}
\end{center}

\textbf{4. Comparatif d'égalité :} \cle{so + Adjectif + wie} (\all{so teuer wie}). \\
\textbf{5. Comparatif de supériorité sans \all{als} :} \all{je mehr, desto besser} (plus..., plus...).
\end{propriete}

\begin{attention}[Pièges fréquents pour francophones]
\begin{itemize}
\item \textbf{\all{als} vs \all{wie}} : On dit \all{billiger \textbf{als}} (comparatif), jamais \all{billiger \textbf{wie}}. \all{wie} sert pour l'égalité (\all{so billig \textbf{wie}}).
\item \textbf{Umlaut au comparatif/superlatif} : Monosyllabes à voyelle \all{a, o, u} $\rightarrow$ \all{ä, ö, ü}. Ex. : \all{alt – älter – am ältesten} ; \all{groß – größer – am größten} ; \all{kalt – kälter – am kältesten}.
\item \textbf{Superlatif attributif} : N'oubliez pas la déclinaison de l'adjectif après l'article ! \all{die billigst\underline{e} Option} (fém.), \all{der billigst\underline{e} Ausflug} (masc.), \all{das billigst\underline{e} Projekt} (neutre).
\item \textbf{\all{am} = \all{an dem}} : Contraction obligatoire au superlatif prédicatif. \all{am nützlichsten}, pas \all{an dem nützlichsten}.
\end{itemize}
\end{attention}

\courssub{Grammaire 2 : Les subordonnées conditionnelles (\all{wenn}) et le Futur I}

\begin{propriete}[La subordonnée conditionnelle avec \all{wenn}]
\all{wenn} (si, quand) introduit une subordonnée. \textbf{Règle d'or : le verbe conjugué se place à la toute fin de la subordonnée.}
La principale qui suit (ou précède) commence par le verbe (inversion).
\begin{center}
\begin{tabularx}{\linewidth}{|X|X|}
\hline
\rowcolor{popblueL} Structure & Exemple \\
\hline
\all{Wenn} + Sujet + ... + \textbf{Verbe}, & \all{Wenn alle Schüler 2\,000 FCFA \textbf{geben}}, \\
Verbe + Sujet + ... & \all{\textbf{werden} wir 180\,000 FCFA \textbf{haben}}. \\
\hline
\all{Wenn wir mehr Geld \textbf{sammeln}}, & \all{\textbf{werden} wir einen Filmabend \textbf{organisieren}}. \\
\hline
\end{tabularx}
\end{center}
\emph{Traduction :} « Si tous les élèves donnent 2\,000 FCFA, nous aurons 180\,000 FCFA. » / « Si nous collectons plus d'argent, nous organiserons une soirée cinéma. »
\end{propriete}

\begin{propriete}[Le Futur I (Futur simple)]
Formation : \cle{\all{werden} (conjugué au présent) + Infinitif du verbe principal \textbf{à la fin de la phrase}}.
\begin{center}
\begin{tabular}{|l|l|l|}
\hline
\rowcolor{popblueL} \textbf{Personne} & \textbf{\all{werden}} & \textbf{Exemple avec \all{sammeln / organisieren}} \\
\hline
\all{ich} & \all{werde} & \all{ich werde sammeln / organisieren} \\
\all{du} & \all{wirst} & \all{du wirst sammeln / organisieren} \\
\all{er / sie / es} & \all{wird} & \all{er wird sammeln / organisieren} \\
\all{wir} & \all{werden} & \all{wir werden sammeln / organisieren} \\
\all{ihr} & \all{werdet} & \all{ihr werdet sammeln / organisieren} \\
\all{sie / Sie} & \all{werden} & \all{sie werden sammeln / organisieren} \\
\hline
\end{tabular}
\end{center}
\textbf{Particularités :}
\begin{itemize}
\item Pas de \all{zu} devant l'infinitif (\all{wir werden sammeln}, jamais \all{wir werden zu sammeln}).
\item Verbes à particule séparable : la particule reste collée à l'infinitif final. \all{wir werden die Plakate an die Wand \textbf{anhängen}} (et non \all{an... hängen}).
\item Verbes en \all{-ieren} / verbes composés longs : infinitif intact à la fin. \all{wir werden das Projekt \textbf{organisieren}}.
\end{itemize}
\end{propriete}

\begin{exercice}[Entraînement grammatical : Comparatif, \all{wenn}, Futur I]
\begin{enumerate}
\item \textbf{Comparatif / Superlatif.} Complétez avec la forme correcte (adjectif entre parenthèses).
\begin{itemize}
\item[a)] Der Ausflug ist \_\_\_\_\_\_\_\_\_\_\_\_ als die Feier. (\all{teuer})
\item[b)] Option C ist \_\_\_\_\_\_\_\_\_\_\_\_ für die Zukunft. (\all{nützlich}, superlatif prédicatif)
\item[c)] Wir wählen \_\_\_\_\_\_\_\_\_\_\_\_ Projekt. (\all{gut}, superlatif attributif accusatif neutre)
\item[d)] Das ist \_\_\_\_\_\_\_\_\_\_\_\_ Lösung für uns. (\all{realistisch}, superlatif attributif nominatif fém.)
\end{itemize}
\item \textbf{Subordonnée \all{wenn} + Futur I.} Mettez les verbes entre parenthèses à la forme correcte (ordre des mots !).
\begin{itemize}
\item[a)] \all{Wenn wir das Geld \_\_\_\_\_\_\_\_\_\_ (haben), \_\_\_\_\_\_\_\_\_\_ wir das Projekt \_\_\_\_\_\_\_\_\_\_ (starten).}
\item[b)] \all{\_\_\_\_\_\_\_\_\_\_ wir die Bücher \_\_\_\_\_\_\_\_\_\_ (kaufen), \_\_\_\_\_\_\_\_\_\_ wir \_\_\_\_\_\_\_\_\_\_ (lesen) jeden Mittwoch.}
\end{itemize}
\item \textbf{Futur I seul.} Conjuguez au Futur I : \all{ich (planen)}, \all{wir (sparen)}, \all{die Klasse (entscheiden)}, \all{der Schatzmeister (rechnen)}.
\end{enumerate}
\end{exercice}

\begin{corrige}[Exercice 2 — Entraînement grammatical]
\begin{enumerate}
\item \textbf{Comparatif / Superlatif.}
\begin{itemize}
\item[a)] \all{teurer als}
\item[b)] \all{am nützlichsten}
\item[c)] \all{das beste} (irrégulier : \all{gut – besser – am besten / der/die/das beste})
\item[d)] \all{die realistischste} (déclinaison fém. nominatif après \all{die})
\end{itemize}
\item \textbf{\all{wenn} + Futur I.}
\begin{itemize}
\item[a)] \all{Wenn wir das Geld \textbf{haben}, \textbf{werden} wir das Projekt \textbf{starten}.} (Subordonnée : verbe \all{haben} à la fin. Principale : \all{werden} en tête, infinitif \all{starten} à la fin).
\item[b)] \all{\textbf{Wenn} wir die Bücher \textbf{kaufen}, \textbf{werden} wir \textbf{lesen} jeden Mittwoch.} (Ou : \all{Wenn wir die Bücher kaufen, werden wir jeden Mittwoch lesen.} — adverbe de temps avant l'infinitif).
\end{itemize}
\item \textbf{Futur I :} \all{ich werde planen} ; \all{wir werden sparen} ; \all{die Klasse wird entscheiden} ; \all{der Schatzmeister wird rechnen}.
\end{enumerate}
\end{corrige}

\courssub{Méthode : Rédiger le paragraphe de synthèse « Unser Projekt »}

\begin{methode}[Étapes pour réussir la production écrite individuelle]
\begin{enumerate}
\item \textbf{Analyser les consignes :} 8–10 phrases. Titre obligatoire : \all{Unser Projekt}. Contraintes : présentation du projet, mini-budget (chiffres), 2 comparaisons (comparatif/superlatif), 1 condition \all{wenn}, 2 phrases au Futur I.
\item \textbf{Structurer le paragraphe (plan type) :}
\begin{itemize}
\item Phrase 1 : Choix du projet (\all{Unsere Gruppe hat Projekt X gewählt.}).
\item Phrases 2–3 : Justification par comparaisons (\all{... ist billiger als ... ; ... ist am nützlichsten.}).
\item Phrases 4–5 : Mini-budget (\all{Einnahmen, Ausgaben, Sparziel} avec chiffres).
\item Phrase 6 : Condition \all{wenn} + Futur I (\all{Wenn ..., werden wir ...}).
\item Phrases 7–8 : Actions futures détaillées (Futur I : \all{wir werden kaufen, organisieren, hängen...}).
\item Phrase 9–10 : Conclusion / Appel au vote (\all{Das ist die beste Lösung, weil ... ; Wir hoffen, dass ...}).
\end{itemize}
\item \textbf{Mobiliser la « Boîte à outils » :} Vérifier chaque phrase : majuscules aux noms, \all{als} au comparatif, \all{am ...sten} / \all{die ...ste} au superlatif, verbe à la fin après \all{wenn}, \all{werden} + infinitif final au Futur I, pas de \all{zu}.
\item \textbf{Relire et corriger :} Accords (genre/nombre/cas), ordre des mots (place du verbe), orthographe (Umlaute, ß).
\end{enumerate}
\end{methode}

\courssub{Activité d'intégration : Mein Budget, mein Projekt}

\begin{experience}[Travail de groupe — Présentation et vote]
\textbf{Durée :} 45–50 min. \textbf{Matériel :} Papier A3 pour afficher le mini-budget, feutres.

\begin{enumerate}
\item \textbf{Préparation (20 min).} Par groupes de 4. Choisissez une option (A, B, C) ou inventez la vôtre (ex. : tournoi de football, plantation d'arbres, spectacle de talents).
\begin{itemize}
\item Remplissez le \textbf{Tableau de budget} (3 lignes : \all{Einnahmen}, \all{Ausgaben}, \all{Sparziel}). Montants réalistes en FCFA.
\item Rédigez \textbf{3 phrases avec \all{wenn}} (conditions de réussite).
\item Préparez la \textbf{présentation orale : 6 phrases au Futur I} (réparties entre les 4 membres). Ex. : \all{Wir werden 2\,000 FCFA pro Schüler sammeln. Ich werde die Bücher kaufen. Wir werden einen Deutschclub gründen.}
\end{itemize}
\item \textbf{Présentations (15 min).} Chaque groupe affiche son budget et présente son projet (6 phrases). La classe prend des notes.
\item \textbf{Vote et débat (10 min).} \all{Wir stimmen für das beste Projekt.} Chaque élève justifie son vote à l'oral avec un superlatif : \all{Ich stimme für Projekt 3, weil es am realistischsten / am nützlichsten ist.}
\end{enumerate}
\end{experience}

\courssub{Production écrite individuelle : Unser Projekt}

\begin{exercice}[Rédaction finale (évaluation)]
Rédigez seul, sur votre cahier, le paragraphe \all{Unser Projekt} (8 à 10 phrases) en respectant scrupuleusement la méthode et les contraintes linguistiques. Utilisez votre tableau de budget et vos phrases \all{wenn} / Futur I préparés en groupe, mais écrivez votre propre texte.
\end{exercice}

\begin{savaistu}[Culture : Gestion de projet et vie associative en Allemagne]
En Allemagne, en Autriche et en Suisse, la \cle{Projektplanung} (planification de projet) est enseignée dès le lycée (\all{Gymnasium}, \all{Berufsschule}) via des \all{Schülerfirmen} (mini-entreprises d'élèves) ou des \all{Projektwochen} (semaines de projet). Les élèves gèrent un vrai budget (\all{Klassenkasse}), rédigent des \all{Businesspläne}, présentent devant un jury (\all{Pitch}) et rendent des comptes (\all{Abschlussbericht}). Au Cameroun, les \cle{Clubs d'allemand} (\all{Deutschclubs}) et les \cle{Associations d'anciens élèves} (\all{Ehemaligenvereine}) fonctionnent sur le même modèle : cotisations (\all{Mitgliedsbeiträge}), trésorier (\all{Schatzmeister}), assemblée générale (\all{Mitgliederversammlung}). Maîtriser ce vocabulaire et ces structures (\all{wenn}, Futur I, comparatif) vous prépare à la vie associative, entrepreneuriale et citoyenne, ici comme là-bas.
\end{savaistu}

\courssub{Production modèle et corrigé commenté}

\begin{exoresolu}[Unser Projekt — Production modèle (8–10 phrases)]
Rédigez le paragraphe de synthèse « Unser Projekt » en respectant les contraintes de la leçon.
\tcblower
\textbf{Modèle de production :}

\all{Unsere Gruppe hat das Projekt „Deutschclub" gewählt. Der Deutschclub ist billiger als der Ausflug nach Kribi, und er ist am nützlichsten für unsere Zukunft. Unser Budget ist überschaubar: Die Einnahmen betragen 120\,000 FCFA, und die Ausgaben sind nur 60\,000 FCFA für Bücher, Spiele und Plakate. Unser Sparziel beträgt 60\,000 FCFA für das nächste Jahr. Wenn alle Schüler mitmachen, werden wir jeden Mittwoch Deutsch spielen und lesen. Wir werden zuerst die Bücher kaufen, und dann werden wir die Plakate an die Wand anhängen. Wenn wir mehr Geld sammeln, werden wir auch einen Filmabend organisieren. Das ist die beste Lösung für unsere Klasse, weil wir dabei viel lernen. Wir hoffen, dass alle für unser Projekt stimmen!}

« Notre groupe a choisi le projet “club d'allemand”. Le club d'allemand est moins cher que l'excursion à Kribi, et c'est le plus utile pour notre avenir. Notre budget est modeste : les recettes s'élèvent à 120\,000 FCFA, et les dépenses à seulement 60\,000 FCFA pour des livres, des jeux et des affiches. Notre objectif d'épargne est de 60\,000 FCFA pour l'année prochaine. Si tous les élèves participent, nous jouerons et lirons en allemand chaque mercredi. Nous achèterons d'abord les livres, puis nous accrocherons les affiches au mur. Si nous collectons plus d'argent, nous organiserons aussi une soirée cinéma. C'est la meilleure solution pour notre classe, parce que nous apprenons beaucoup ainsi. Nous espérons que tous voteront pour notre projet ! »

\textbf{Commentaire des choix de langue :}
\begin{itemize}
\item \textbf{Comparatif et superlatif} : \all{billiger als} (comparatif + \all{als}), \all{am nützlichsten} (superlatif prédicatif \all{am} + adjectif + \all{-sten} avec Umlaut \all{ü}), \all{die beste Lösung} (superlatif attributif irrégulier de \all{gut} : \all{die beste} fém. nominatif/accusatif).
\item \textbf{Conditionnelles \all{wenn}} : \all{Wenn alle Schüler mitmachen, werden wir...} — verbe \all{mitmachen} à la fin de la subordonnée, inversion verbe-sujet (\all{werden wir}) dans la principale. Même structure : \all{Wenn wir mehr Geld sammeln, werden wir...}.
\item \textbf{Futur I} : \all{werden wir spielen, wir werden kaufen, wir werden anhängen, werden wir organisieren} — auxiliaire \all{werden} conjugué en 2e position, infinitif en fin de phrase. Noter le verbe à particule \all{anhängen} (accrocher) : l'infinitif reste soudé (\all{anhängen}) à la fin.
\item \textbf{Lexique du thème} : \all{das Budget, die Einnahmen, die Ausgaben, das Sparziel, sammeln, betragen, planen, gründen} — tous les mots-clés de la leçon sont intégrés naturellement.
\item \textbf{Connecteurs et structure} : \all{und} (coordination), \all{weil} (subordonnée causale, verbe à la fin), \all{dann} (adverbe de temps en tête de principale $\rightarrow$ inversion \all{werden wir}), \all{zuerst / dann} (chronologie).
\end{itemize}
\end{exoresolu}

\courssub{Bilan de la leçon}

Cette leçon a permis de réinvestir l'ensemble des acquis du Chapitre 2 « Vie économique » dans une tâche finale complexe et motivante, ancrée dans le contexte camerounais (Yaoundé, Kribi, FCFA). 

D'un point de vue \textbf{linguistique}, les élèves ont :
\begin{itemize}
\item consolidé le \cle{lexique économique} (\all{Budget, Einnahmen, Ausgaben, Sparziel, sammeln, kosten, Plan}) à travers un texte authentique et un tableau de budget réel ;
\item maîtrisé la \textbf{formation et l'emploi du comparatif et du superlatif} (réguliers, Umlaut, irréguliers \all{gut/viel/gerne}, distinction attributif/prédicatif) ;
\item automatisé la \textbf{structure de la subordonnée \all{wenn}} (verbe final) et l'\textbf{inversion en principale} (verbe en tête) ;
\item construit correctement le \textbf{Futur I} (\all{werden} + infinitif final), y compris avec verbes à particule séparable (\all{anhängen}).
\end{itemize}

D'un point de vue \textbf{méthodologique}, ils ont pratiqué les quatre compétences : \textbf{compréhension écrite} (note du trésorier, questions), \textbf{expression orale en interaction} (débat de groupe, présentation, vote justifié), \textbf{expression écrite guidée} (exercices d'application, paragraphe final structuré) et \textbf{traduction implicite} (glossaire, reformulation).

D'un point de vue \textbf{citoyen et transversal}, l'activité développe la \cle{compétence budgétaire} (réalisme des coûts, épargne), l'\cle{argumentation rationnelle} (comparaisons, superlatifs) et la \cle{prise de décision collective} (vote, justification), compétences transversales valorisées au Baccalauréat et dans l'enseignement supérieur (projets tutorés, associations, entrepreneuriat).

\textbf{Prolongement possible :} Le professeur conserve les meilleurs paragraphes \all{Unser Projekt} dans le portfolio de la classe. Le projet gagnant du vote peut être mis en œuvre réellement en fin d'année (création effective du \all{Deutschclub}, organisation de la fête ou de l'excursion), transformant ainsi la simulation pédagogique en action concrète.

\newpage

\coursec{Méthodes et exercices résolus}

\courssub{Méthode 1 : Comprendre un texte sur le budget et les projets}

\begin{methode}[Lire et exploiter un texte économique]
Pour comprendre un texte allemand sur \all{das Budget} ou \all{das Projekt}, procède en quatre étapes :
\begin{enumerate}
\item \cle{Lecture globale} : lis le texte une première fois sans dictionnaire. Repère le thème (argent, épargne, projet), le type de texte (article, reportage, témoignage) et le ton (informatif, critique).
\item \cle{Repérage des mots-clés} : souligne le vocabulaire du thème : \all{das Geld}, \all{die Einnahmen} (les recettes), \all{die Ausgaben} (les dépenses), \all{sparen} (économiser), \all{der Plan}, \all{das Projekt}. Ces mots portent le sens du texte.
\item \cle{Lecture détaillée} : pour chaque question, localise le passage correspondant, traduis-le mentalement, puis reformule la réponse avec tes propres mots en allemand — jamais une simple recopie de la phrase du texte.
\item \cle{Vérification} : relis ta réponse : verbe conjugué en 2\up{e} position, majuscule aux noms, cas corrects après les prépositions (\all{für das Projekt}, \all{mit dem Geld}).
\end{enumerate}
Réflexe : une question avec \all{Warum} appelle une réponse avec \all{weil} ou \all{denn} ; une question avec \all{Wie} appelle une réponse avec un complément de manière ; une question avec \all{Was} appelle une réponse avec un groupe nominal ou une phrase complète.
\end{methode}

\begin{exoresolu}[Compréhension de texte : questions type Bac]
\textbf{Texte (extrait) :} \all{Viele Jugendliche in Yaoundé bekommen jeden Monat ein kleines Taschengeld. Wer sein Geld gut plant, kann am Ende des Monats etwas sparen. Deshalb lernen die Schüler im Deutschclub, ein Budget zu machen: Sie schreiben ihre Einnahmen und Ausgaben in ein Heft.}\par
\textbf{Questions :} 1) \all{Was bekommen viele Jugendliche in Yaoundé?} 2) \all{Was kann man sparen, wenn man sein Geld gut plant?} 3) \all{Warum machen die Schüler ein Budget?}
\tcblower
\textbf{Raisonnement :} la question 1 porte sur la première phrase (\all{bekommen} = recevoir) ; la question 2 sur la deuxième (\all{sparen}) ; la question 3 demande la cause, signalée par \all{deshalb} (c'est pourquoi).\par
\textbf{Réponses rédigées :}\par
1) \all{Viele Jugendliche in Yaoundé bekommen jeden Monat ein kleines Taschengeld.} « Beaucoup de jeunes à Yaoundé reçoivent chaque mois un petit argent de poche. »\par
2) \all{Wenn man sein Geld gut plant, kann man am Ende des Monats etwas sparen.} « Si l'on planifie bien son argent, on peut économiser quelque chose à la fin du mois. »\par
3) \all{Die Schüler machen ein Budget, weil sie lernen wollen, ihre Einnahmen und Ausgaben zu kontrollieren.} « Les élèves font un budget parce qu'ils veulent apprendre à contrôler leurs recettes et leurs dépenses. »\par
\textbf{Justifications grammaticales :} à la question 3, la subordonnée avec \all{weil} envoie le verbe conjugué à la fin (\all{wollen}). À la question 2, la subordonnée \all{wenn} placée en tête fait passer le verbe de la principale avant le sujet (\all{kann man}) : c'est l'inversion obligatoire.
\end{exoresolu}

\courssub{Méthode 2 : Employer le vocabulaire du thème}

\begin{methode}[Maîtriser le lexique de l'argent et des projets]
\begin{enumerate}
\item Apprends chaque nom \cle{avec son article et son pluriel} : \all{das Budget, die Budgets} ; \all{die Ausgabe, -n} ; \all{die Einnahme, -n} ; \all{das Geld} (sans pluriel courant) ; \all{das Projekt, -e} ; \all{der Plan, die Pläne}.
\item Retiens les \cle{couples d'opposés} : \all{die Einnahmen} (recettes) / \all{die Ausgaben} (dépenses) ; \all{sparen} (économiser) / \all{ausgeben} (dépenser) ; \all{verdienen} (gagner) / \all{verlieren} (perdre).
\item Construis des \cle{familles de mots} : \all{sparen} $\rightarrow$ \all{der Sparer}, \all{das Sparen}, \all{sparsam} (économe) ; \all{planen} $\rightarrow$ \all{der Plan}, \all{die Planung} ; \all{ausgeben} $\rightarrow$ \all{die Ausgabe} ; \all{einnehmen} $\rightarrow$ \all{die Einnahme}.
\item Mémorise les \cle{expressions figées} : \all{mit Geld umgehen} (gérer l'argent), \all{Geld ausgeben für + accusatif} (dépenser de l'argent pour), \all{ein Projekt planen / durchführen} (planifier / réaliser un projet), \all{ein Budget machen / aufstellen} (faire / établir un budget).
\end{enumerate}
\end{methode}

\begin{exoresolu}[Lexique : texte à trous]
\textbf{Énoncé :} Complète avec : \all{Budget, Ausgaben, sparen, Projekt, Einnahmen}.\par
\all{Für unser Schulprojekt machen wir ein (1) \dots. Unsere (2) \dots sind klein: nur 10\,000 FCFA. Aber unsere (3) \dots sind größer: das Material kostet 12\,000 FCFA. Deshalb müssen wir (4) \dots. Am Ende wollen wir das (5) \dots erfolgreich beenden.}
\tcblower
\textbf{Raisonnement :} on cherche d'abord les oppositions logiques : recettes petites / dépenses grandes $\rightarrow$ il faut économiser.\par
\textbf{Solution :} (1) \all{Budget} (on « fait un budget » : \all{ein Budget machen}) ; (2) \all{Einnahmen} (les recettes, opposé de dépenses) ; (3) \all{Ausgaben} (les dépenses, confirmé par \all{das Material kostet}) ; (4) \all{sparen} (conséquence logique : \all{deshalb}) ; (5) \all{Projekt} (\all{ein Projekt beenden} = terminer un projet).\par
\textbf{Texte complet :} \all{Für unser Schulprojekt machen wir ein Budget. Unsere Einnahmen sind klein: nur 10\,000 FCFA. Aber unsere Ausgaben sind größer: das Material kostet 12\,000 FCFA. Deshalb müssen wir sparen. Am Ende wollen wir das Projekt erfolgreich beenden.}\par
« Pour notre projet scolaire, nous faisons un budget. Nos recettes sont petites : seulement 10\,000 FCFA. Mais nos dépenses sont plus grandes : le matériel coûte 12\,000 FCFA. C'est pourquoi nous devons économiser. À la fin, nous voulons terminer le projet avec succès. »
\end{exoresolu}

\courssub{Méthode 3 : Former et employer le comparatif et le superlatif}

\begin{methode}[Le comparatif et le superlatif de l'adjectif]
\begin{enumerate}
\item \cle{Formes} : positif \all{billig} $\rightarrow$ comparatif \all{billiger} $\rightarrow$ superlatif \all{am billigsten} (adverbe) / \all{der billigste} (attribut).
\item \cle{Umlaut} aux voyelles a, o, u des adjectifs monosyllabiques : \all{alt $\rightarrow$ älter $\rightarrow$ am ältesten} ; \all{groß $\rightarrow$ größer $\rightarrow$ am größten} ; \all{jung $\rightarrow$ jünger $\rightarrow$ am jüngsten}.
\item \cle{Irréguliers à connaître par cœur} : \all{gut $\rightarrow$ besser $\rightarrow$ am besten} ; \all{viel $\rightarrow$ mehr $\rightarrow$ am meisten} ; \all{wenig $\rightarrow$ weniger $\rightarrow$ am wenigsten} ; \all{gern $\rightarrow$ lieber $\rightarrow$ am liebsten} ; \all{hoch $\rightarrow$ höher $\rightarrow$ am höchsten} ; \all{nah $\rightarrow$ näher $\rightarrow$ am nächsten}.
\item \cle{Comparaison} : égalité \all{so \dots\ wie} (\all{Paul ist so fleißig wie Anna.} « Paul est aussi travailleur qu'Anna. ») ; supériorité ou différence avec \all{als} (\all{Das Projekt ist teurer als der Plan.} « Le projet est plus cher que le plan. »).
\item \cle{Superlatif attribut} : devant un nom, il prend l'article défini et se décline comme un adjectif : \all{das beste Projekt}, \all{die größten Ausgaben}, \all{mit dem kleinsten Budget}.
\end{enumerate}
\end{methode}

\begin{propriete}[Tableau récapitulatif des formes]
\begin{center}
\begin{tabular}{|l|l|l|l|}
\hline
\rowcolor{popblueL} Positif & Comparatif & Superlatif (adverbe) & Superlatif (attribut) \\
\hline
\all{billig} (bon marché) & \all{billiger} & \all{am billigsten} & \all{der/die/das billigste} \\
\hline
\all{teuer} (cher) & \all{teurer} & \all{am teuersten} & \all{der/die/das teuerste} \\
\hline
\all{groß} (grand) & \all{größer} & \all{am größten} & \all{der/die/das größte} \\
\hline
\all{alt} (vieux, âgé) & \all{älter} & \all{am ältesten} & \all{der/die/das älteste} \\
\hline
\all{gut} (bon) & \all{besser} & \all{am besten} & \all{der/die/das beste} \\
\hline
\all{viel} (beaucoup) & \all{mehr} & \all{am meisten} & \all{das meiste} \\
\hline
\end{tabular}
\end{center}
Remarque : les adjectifs terminés par \all{-d}, \all{-t}, \all{-s}, \all{-ß}, \all{-sch} prennent \all{-esten} au superlatif : \all{alt $\rightarrow$ am ältesten}, \all{groß $\rightarrow$ am größten}. L'adjectif \all{teuer} perd son \all{-e-} au comparatif : \all{teurer} (et non \all{teuerer}).
\end{propriete}

\begin{exoresolu}[Grammaire : comparatifs et superlatifs]
\textbf{Énoncé :} Mets l'adjectif entre parenthèses à la forme correcte.\par
1) \all{Ein Fahrrad ist (billig) \dots\ als ein Auto.} 2) \all{Das Sparen ist (gut) \dots\ als das Ausgeben.} 3) \all{Der Markt in Douala ist (groß) \dots\ Markt der Stadt.} 4) \all{Mein Bruder ist (alt) \dots\ als ich.}
\tcblower
\textbf{Raisonnement :} les phrases 1, 2 et 4 contiennent \all{als} $\rightarrow$ comparatif ; la phrase 3 désigne le plus grand de tous $\rightarrow$ superlatif attribut avec article défini.\par
\textbf{Solutions :}\par
1) \all{Ein Fahrrad ist billiger als ein Auto.} « Un vélo est moins cher qu'une voiture. » (comparatif régulier : \all{billig + er}).\par
2) \all{Das Sparen ist besser als das Ausgeben.} « Économiser est mieux que dépenser. » (irrégulier : \all{gut $\rightarrow$ besser}).\par
3) \all{Der Markt in Douala ist der größte Markt der Stadt.} « Le marché de Douala est le plus grand marché de la ville. » (Umlaut : \all{groß $\rightarrow$ größ-} ; superlatif attribut : \all{der größte}, avec déclinaison \all{-e} après l'article défini).\par
4) \all{Mein Bruder ist älter als ich.} « Mon frère est plus âgé que moi. » (Umlaut : \all{alt $\rightarrow$ älter} ; après \all{als}, le nominatif \all{ich}).\par
\textbf{Conclusion :} toujours vérifier : voyelle a/o/u monosyllabique $\rightarrow$ Umlaut ; adjectif irrégulier $\rightarrow$ forme apprise par cœur ; superlatif attribut $\rightarrow$ article défini + déclinaison.
\end{exoresolu}

\courssub{Méthode 4 : Construire une conditionnelle avec wenn et le Futur I}

\begin{methode}[La phrase conditionnelle : wenn + Futur I]
\begin{enumerate}
\item \cle{Structure de base} : \all{Wenn} + sujet + compléments + \textbf{verbe conjugué à la fin} (subordonnée), puis virgule, puis la principale avec le \textbf{verbe en 1\iere{} position} : \all{Wenn ich Geld spare, kaufe ich ein Fahrrad.} « Si j'économise de l'argent, j'achète un vélo. »
\item \cle{Futur I} : auxiliaire \all{werden} conjugué + infinitif à la fin : \all{ich werde sparen}, \all{wir werden das Projekt beginnen}.
\item \cle{Condition + futur} : pour un projet futur, combine les deux : \all{Wenn ich genug Geld habe, werde ich ein Projekt starten.} « Si j'ai assez d'argent, je lancerai un projet. »
\item \cle{Ordre inversé possible} : la principale d'abord, sans changement de sens : \all{Ich werde ein Projekt starten, wenn ich genug Geld habe.}
\item Attention : \all{wenn} sert aussi pour le temps (\all{quand}, au présent ou au futur) et l'habitude ; le contexte décide. Pour un événement passé unique, on utilise \all{als} : \all{Als ich Kind war, \dots} « Quand j'étais enfant, \dots »
\end{enumerate}
\end{methode}

\begin{propriete}[Conjugaison de werden au Futur I]
\begin{center}
\begin{tabular}{|l|l|l|}
\hline
\rowcolor{popblueL} Personne & werden + Infinitiv & Exemple \\
\hline
\all{ich} & \all{werde \dots\ sparen} & \all{Ich werde Geld sparen.} \\
\hline
\all{du} & \all{wirst \dots\ sparen} & \all{Du wirst Geld sparen.} \\
\hline
\all{er/sie/es} & \all{wird \dots\ sparen} & \all{Sie wird Geld sparen.} \\
\hline
\all{wir} & \all{werden \dots\ sparen} & \all{Wir werden ein Projekt starten.} \\
\hline
\all{ihr} & \all{werdet \dots\ sparen} & \all{Ihr werdet ein Budget machen.} \\
\hline
\all{sie/Sie} & \all{werden \dots\ sparen} & \all{Sie werden sparen.} \\
\hline
\end{tabular}
\end{center}
L'auxiliaire \all{werden} occupe la 2\up{e} position de la phrase (ou la 1\iere{} si la subordonnée précède) ; l'infinitif se place à la toute fin : \all{Ich werde nächstes Jahr mein Projekt beginnen.} « Je commencerai mon projet l'année prochaine. »
\end{propriete}

\begin{exoresolu}[Thème : traduction français $\rightarrow$ allemand]
\textbf{Énoncé :} Traduis en allemand :\par
1) « Si j'économise mon argent de poche, j'achèterai un ordinateur. »\par
2) « Nous réaliserons notre projet, si nous faisons un bon budget. »\par
3) « Quand les élèves planifient bien, ils dépensent moins. »
\tcblower
\textbf{Raisonnement :} phrases 1 et 2 : condition réelle au futur $\rightarrow$ \all{wenn} + Präsens dans la subordonnée, \all{werden} + infinitif dans la principale. Phrase 3 : habitude $\rightarrow$ \all{wenn} + Präsens des deux côtés ; « moins » = \all{weniger} (comparatif de \all{wenig}).\par
\textbf{Solutions :}\par
1) \all{Wenn ich mein Taschengeld spare, werde ich einen Computer kaufen.} Subordonnée : verbe \all{spare} à la fin ; principale : \all{werde} en 1\iere{} position, infinitif \all{kaufen} à la fin ; \all{einen Computer} : accusatif masculin.\par
2) \all{Wir werden unser Projekt durchführen, wenn wir ein gutes Budget machen.} Principale d'abord : \all{werden} en 2\up{e} position, infinitif \all{durchführen} à la fin ; \all{ein gutes Budget} : accusatif neutre ; verbe de la subordonnée \all{machen} à la fin.\par
3) \all{Wenn die Schüler gut planen, geben sie weniger aus.} Attention au verbe à particule : \all{ausgeben} $\rightarrow$ \all{geben \dots\ aus}, la particule \all{aus} va à la fin de la principale.\par
\textbf{Conclusion :} dans une conditionnelle, deux pièges à éviter : la place du verbe dans la subordonnée et l'inversion sujet-verbe dans la principale.
\end{exoresolu}

\courssub{Méthode 5 : Rédiger un court paragraphe : Mein Budget, mein Projekt}

\begin{methode}[Produire un texte guidé sur son budget et son projet]
\begin{enumerate}
\item \cle{Plan en 4 étapes} : (a) présenter le projet (\all{Ich habe ein Projekt: \dots}) ; (b) les moyens (\all{Ich bekomme / verdiene \dots}) ; (c) le budget (\all{Meine Einnahmen sind \dots, meine Ausgaben sind \dots}) ; (d) la condition et le futur (\all{Wenn ich genug spare, werde ich \dots}).
\item \cle{Connecteurs} : \all{zuerst} (d'abord), \all{dann} (ensuite), \all{deshalb} (c'est pourquoi), \all{am Ende} (à la fin).
\item \cle{Varier les structures} : au moins un comparatif (\all{besser als}), une conditionnelle \all{wenn}, un Futur I.
\item \cle{Relecture} : majuscules aux noms, verbe en 2\up{e} position, particules séparables à la fin, accords des articles et des adjectifs.
\end{enumerate}
\end{methode}

\begin{exoresolu}[Expression écrite guidée : rédaction]
\textbf{Énoncé :} Rédige un paragraphe de 5 à 7 phrases intitulé \all{Mein Projekt} : présente un projet, ton budget, et dis ce que tu feras si tu économises assez d'argent. Utilise au moins un comparatif, une phrase avec \all{wenn} et le Futur I.
\tcblower
\textbf{Raisonnement :} suivre le plan en 4 étapes, intégrer les trois structures exigées, puis relire.\par
\textbf{Production modèle :}\par
\all{Ich habe ein Projekt: Ich möchte einen kleinen Laden für Schulmaterial in Yaoundé öffnen. Jeden Monat bekomme ich 15\,000 FCFA Taschengeld. Meine Ausgaben sind kleiner als meine Einnahmen, deshalb kann ich jeden Monat 5\,000 FCFA sparen. Zuerst schreibe ich mein Budget in ein Heft, dann kaufe ich nur das Wichtigste. Wenn ich ein Jahr lang genug spare, werde ich mein Projekt beginnen. Am Ende werde ich mein eigenes Geld verdienen, und das ist besser als Taschengeld!}\par
« J'ai un projet : je voudrais ouvrir une petite boutique de fournitures scolaires à Yaoundé. Chaque mois, je reçois 15\,000 FCFA d'argent de poche. Mes dépenses sont plus petites que mes recettes, c'est pourquoi je peux économiser 5\,000 FCFA par mois. D'abord, j'écris mon budget dans un cahier, ensuite je n'achète que l'essentiel. Si j'économise assez pendant un an, je commencerai mon projet. À la fin, je gagnerai mon propre argent, et c'est mieux que l'argent de poche ! »\par
\textbf{Vérification :} comparatifs présents (\all{kleiner als}, \all{besser als}) ; conditionnelle \all{wenn} avec verbe conjugué final (\all{spare}) et inversion dans la principale (\all{werde ich}) ; Futur I correct (\all{werde ich mein Projekt beginnen}, \all{werde ich mein eigenes Geld verdienen}) ; \all{öffnen} est un infinitif après \all{möchte}, placé à la fin ; majuscules à tous les noms (\all{Projekt, Laden, Taschengeld, Ausgaben, Einnahmen, Budget, Heft, Ende, Geld}).
\end{exoresolu}

\begin{attention}[Les 5 erreurs qui coûtent des points au Bac]
\begin{enumerate}
\item \cle{Oublier la majuscule aux noms} : \all{das geld}, \all{das projekt} sont FAUX. Écris toujours \all{das Geld}, \all{das Projekt}, \all{das Budget}.
\item \cle{Mal placer le verbe après wenn} : \all{Wenn ich spare Geld, \dots} est un calque du français. Correct : \all{Wenn ich Geld spare, \dots} — le verbe conjugué va à la \textbf{fin} de la subordonnée, et la principale commence par le verbe : \all{\dots, werde ich sparen.}
\item \cle{Confondre als et wie} : après un comparatif, c'est toujours \all{als}, jamais \all{wie} : \all{besser als}, \all{mehr als}. \all{wie} ne sert que pour l'égalité : \all{so groß wie}. Et ne confonds pas \all{wenn} (condition, habitude, futur) avec \all{als} (passé unique).
\item \cle{Faux amis et calques} : « faire un budget » se dit \all{ein Budget machen} (pas \all{ein Budget tun}) ; « dépenser de l'argent » = \all{Geld ausgeben} (particule séparable : \all{ich gebe Geld aus}) ; \all{sparen} signifie « économiser », pas « épargner quelqu'un » ; \all{bekommen} signifie « recevoir », pas « devenir » (\all{werden}).
\item \cle{Oublier l'Umlaut au comparatif} : \all{alter}, \all{grosser} sont faux. Correct : \all{älter}, \all{größer}. Et au superlatif : \all{am größten}, \all{der beste} (jamais \all{der guteste}).
\end{enumerate}
\end{attention}

\newpage

\coursec{Exercices}

\courssub{Je vérifie mes connaissances}

\begin{exercice}[QCM : vocabulaire du budget]
Entoure la bonne réponse.
\begin{enumerate}
\item \all{Die Einnahmen} signifie : a) les dépenses \quad b) les recettes \quad c) les dettes
\item Le pluriel de \all{das Budget} est : a) \all{die Budgete} \quad b) \all{die Budgets} \quad c) \all{die Budgeten}
\item \all{sparen} veut dire : a) dépenser \quad b) emprunter \quad c) économiser
\item Le contraire de \all{die Ausgabe} est : a) \all{der Plan} \quad b) \all{die Einnahme} \quad c) \all{das Geld}
\item \all{mit Geld umgehen} signifie : a) gagner de l'argent \quad b) gérer l'argent \quad c) perdre de l'argent
\end{enumerate}
\end{exercice}

\begin{exercice}[Vrai ou faux — justifie ta réponse]
Réponds par \all{richtig} ou \all{falsch} et justifie en une phrase en français.
\begin{enumerate}
\item Le comparatif de \all{gut} est \all{guter}.
\item Dans une subordonnée introduite par \all{wenn}, le verbe conjugué se place à la fin.
\item Le Futur I se forme avec \all{haben} + participe II.
\item Le superlatif de \all{hoch} est \all{am höchsten}.
\item On dit \all{mehr billiger} pour « moins cher ».
\end{enumerate}
\end{exercice}

\begin{exercice}[Vocabulaire : trouve le mot]
Complète chaque phrase avec le mot qui convient : \all{das Budget}, \all{die Ausgaben}, \all{sparen}, \all{das Projekt}, \all{der Plan}, \all{die Einnahmen}.
\begin{enumerate}
\item Am Ende des Monats sind meine \ldots\ höher als meine \ldots : ich habe ein Problem !
\item Unser \ldots\ für die Klassenfahrt beträgt 50\,000 Franken.
\item Ich muss jeden Monat etwas Geld \ldots, um ein Handy zu kaufen.
\item Unser größtes \ldots\ dieses Jahr ist die Reise nach Deutschland.
\item Wir machen einen \ldots\ für die nächsten drei Monate.
\end{enumerate}
\end{exercice}

\begin{exercice}[Conjugaison : Futur I à compléter]
Complète avec la forme correcte du verbe \all{werden} au présent.
\begin{enumerate}
\item Ich \ldots\ nächstes Jahr mehr sparen.
\item Wir \ldots\ ein neues Projekt beginnen.
\item \ldots\ du dein Geld besser planen ?
\item Die Schüler \ldots\ eine Reise nach Wien organisieren.
\item Mein Bruder \ldots\ weniger ausgeben.
\end{enumerate}
\end{exercice}

\courssub{Je m'entraîne}

\begin{exercice}[Comparatif et superlatif : texte à trous]
Complète les 8 phrases avec le comparatif ou le superlatif de l'adjectif entre parenthèses.
\begin{enumerate}
\item Ein Fahrrad ist \ldots\ als ein Auto. (billig)
\item Das Leben in Douala ist \ldots\ als auf dem Land. (teuer)
\item Mein Bruder spart \ldots\ Geld als ich. (viel)
\item Ich kaufe \ldots\ auf dem Markt als im Supermarkt. (gern)
\item Der \ldots\ Monat für unsere Familie ist der Dezember. (teuer)
\item Yaoundé ist eine der \ldots\ Städte Kameruns. (groß)
\item Das \ldots\ Projekt der Klasse ist die Deutschlandreise. (interessant)
\item Die Preise in der Stadt sind am \ldots. (hoch)
\end{enumerate}
\end{exercice}

\begin{exercice}[Transformation : phrases avec \all{wenn}]
Transforme les 6 phrases suivantes en subordonnées conditionnelles avec \all{wenn}. Attention à la place du verbe !
\begin{enumerate}
\item Er spart Geld. Er reist nach Wien.
\item Ich habe Zeit. Ich helfe meinen Eltern.
\item Wir arbeiten gut. Wir bekommen ein besseres Budget.
\item Das Wetter ist schön. Wir verkaufen mehr auf dem Markt.
\item Du planst deine Ausgaben. Du hast am Ende des Monats Geld.
\item Die Schüler lernen Deutsch. Sie können in Deutschland studieren.
\end{enumerate}
\end{exercice}

\begin{exercice}[Futur I : mise au temps et traduction]
Mets les 6 phrases suivantes au Futur I, puis traduis-les en français.
\begin{enumerate}
\item Ich spare jeden Monat 5\,000 Franken.
\item Wir beginnen unser Projekt im Januar.
\item Meine Schwester kauft einen neuen Computer.
\item Die Familie plant ihr Budget besser.
\item Ihr reist nächstes Jahr nach Deutschland.
\item Er gibt weniger Geld aus.
\end{enumerate}
\end{exercice}

\begin{exercice}[Appariement : mots et définitions]
Relie chaque mot allemand à sa définition française.

\begin{center}
\begin{tabularx}{\linewidth}{|X|X|}
\hline
\rowcolor{popblueL} \textbf{Deutsch} & \textbf{Français} \\
\hline
1. \all{das Taschengeld} & a) somme d'argent disponible pour un projet \\
\hline
2. \all{das Budget} & b) argent que l'on reçoit \\
\hline
3. \all{die Einnahme} & c) argent de poche \\
\hline
4. \all{die Ausgabe} & d) mettre de l'argent de côté \\
\hline
5. \all{sparen} & e) argent que l'on dépense \\
\hline
6. \all{der Plan} & f) projet organisé, programme d'actions \\
\hline
\end{tabularx}
\end{center}
\end{exercice}

\courssub{Je me prépare au Bac}

\begin{exercice}[Texte et questions de compréhension — type Bac]
Lis le texte suivant, puis réponds aux questions \emph{en allemand, par des phrases complètes}.

\begin{textebox}[Lernen, mit Geld umzugehen]
Viele Jugendliche in Kamerun bekommen von ihren Eltern ein kleines Taschengeld. Aber wer gut mit Geld umgehen will, muss ein Budget planen. Zuerst schreibt man alle Einnahmen auf, dann alle Ausgaben. Wenn die Ausgaben höher sind als die Einnahmen, muss man sparen. Das ist oft schwieriger, als man denkt, denn das Leben in der Stadt wird immer teurer. Familie Ngo Bell aus Yaoundé hat ein gutes System: Am Anfang jedes Monats macht die Familie einen Plan. Das wichtigste Projekt der Familie ist die Ausbildung der Kinder. „Wenn wir jeden Monat etwas sparen, werden wir nächstes Jahr einen Computer für unsere Tochter kaufen“, sagt der Vater. Die Tochter Marie findet dieses Projekt am interessantesten, denn sie möchte später Informatik studieren. Experten sagen: Wer früh sparen lernt, wird später weniger Probleme mit Geld haben.
\end{textebox}

\textbf{Fragen zum Text :}
\begin{enumerate}
\item Was muss man zuerst aufschreiben, wenn man ein Budget plant ?
\item \all{Richtig oder falsch ?} Begründe mit einem Satz aus dem Text : \emph{Das Leben in der Stadt wird immer billiger.}
\item Was ist das wichtigste Projekt der Familie Ngo Bell ?
\item \all{Richtig oder falsch ?} Begründe mit einem Satz aus dem Text : \emph{Die Familie will nächstes Jahr ein Handy kaufen.}
\item Was sagen die Experten über das Sparen ?
\end{enumerate}

\textbf{Fragen zur Sprache :}
\begin{enumerate}
\item Finde im Text einen Komparativ und einen Superlativ und bilde mit jedem einen eigenen Satz.
\item Setze den Satz ins Futur I : \all{Die Familie macht einen Plan.}
\end{enumerate}
\end{exercice}

\begin{exercice}[Thème et version — type Bac]
\textbf{A. Thème.} Traduis les 5 phrases suivantes en allemand.
\begin{enumerate}
\item Si j'ai assez d'argent, j'achèterai un nouvel ordinateur.
\item Mon projet le plus important est le voyage en Allemagne.
\item La vie est plus chère en ville qu'au village.
\item Quand je serai grand, je gérerai mieux mon budget. (utiliser \all{wenn} + Futur I)
\item Ma sœur économise plus d'argent que moi parce qu'elle dépense moins.
\end{enumerate}

\textbf{B. Version.} Traduis le texte suivant en français.

\begin{textebox}[Das Budget einer Familie]
Familie Müller aus München plant ihr Budget jeden Monat genau. Die Einnahmen sind höher als die Ausgaben, aber das Leben in München ist teurer als in vielen anderen deutschen Städten. Das größte Projekt der Familie ist der Kauf eines Hauses. Wenn die Familie jeden Monat 500 Euro spart, wird sie in fünf Jahren genug Geld haben. „Sparen ist manchmal schwierig, aber es ist das Wichtigste für unsere Zukunft“, sagt Frau Müller.
\end{textebox}
\end{exercice}

\begin{exercice}[Production écrite — type Bac]
\textbf{Sujet :} \all{Beschreibe dein Budget und dein wichtigstes Projekt.}

Rédige un texte de \textbf{8 à 10 lignes} en allemand. Tu répondras aux questions suivantes : D'où vient ton argent ? Quelles sont tes dépenses principales ? Quel est ton projet le plus important et comment vas-tu le financer ?

\textbf{Contraintes obligatoires} — ton texte doit contenir :
\begin{enumerate}
\item un \textbf{comparatif} (ex. \all{teurer als}, \all{mehr als}) ;
\item un \textbf{superlatif} (ex. \all{am wichtigsten}, \all{das größte Projekt}) ;
\item une \textbf{subordonnée avec \all{wenn}} (attention à la place du verbe !) ;
\item une phrase au \textbf{Futur I} (\all{werden} + infinitif).
\end{enumerate}

Souligne dans ta copie chacune des quatre structures exigées.
\end{exercice}



\newpage

\coursec{Corrigés des exercices}

\coursec{Corrigés des exercices — La gestion du budget et des projets}

\begin{corrige}[Exercice 1 — QCM : vocabulaire du budget]
\begin{enumerate}
\item \textbf{b)} les recettes. \all{Die Einnahme, -n} = la recette, le revenu (du verbe \all{einnehmen}, encaisser) ; le contraire est \all{die Ausgabe, -n} (la dépense).
\item \textbf{b)} \all{die Budgets}. Les noms étrangers récents prennent souvent un \textbf{-s} au pluriel : \all{das Budget, die Budgets}. Attention, \all{das Projekt} fait \all{die Projekte} (pluriel en \textbf{-e}).
\item \textbf{c)} économiser. \all{sparen} = économiser, épargner ; « dépenser » se dit \all{ausgeben} (verbe à particule séparable : \all{Ich gebe Geld aus.}).
\item \textbf{b)} \all{die Einnahme}. Le couple à retenir : \all{die Einnahme} (recette) / \all{die Ausgabe} (dépense).
\item \textbf{b)} gérer l'argent. \all{mit Geld umgehen} = gérer, manier l'argent (verbe à particule séparable : \all{Er geht gut mit Geld um.}).
\end{enumerate}
\textbf{Points de vigilance :} ne pas confondre \all{sparen} (\cle{économiser}) et \all{spenden} (\cle{faire un don}) ; retenir chaque nom avec son article et son pluriel (\all{das Budget, die Budgets} ; \all{das Projekt, die Projekte} ; \all{der Plan, die Pläne}).
\end{corrige}

\begin{corrige}[Exercice 2 — Vrai ou faux]
\begin{enumerate}
\item \all{Falsch}. Le comparatif de \all{gut} est irrégulier : \all{gut -- besser -- am besten}. \all{Guter} est simplement l'adjectif \all{gut} décliné (ex. \all{ein guter Plan}).
\item \all{Richtig}. Après \all{wenn}, le verbe conjugué va en dernière position : \all{Wenn ich Geld habe, kaufe ich ein Buch.}
\item \all{Falsch}. Le Futur I se forme avec \all{werden} (conjugué au présent) + l'infinitif placé à la fin : \all{Ich werde sparen.} La structure \all{haben} + participe II correspond au Perfekt (\all{Ich habe gespart.}).
\item \all{Richtig}. \all{hoch -- höher -- am höchsten} : l'adjectif perd son \textbf{c} au comparatif et au superlatif.
\item \all{Falsch}. On ne peut pas accumuler deux marques de comparaison. « Moins cher » se dit \all{weniger teuer} ; « meilleur marché » se dit simplement \all{billiger}.
\end{enumerate}
\textbf{Points de vigilance :} l'erreur \cle{mehr billiger} est un calque du français « plus moins cher » : en allemand, une seule marque de comparaison par adjectif (soit \cle{-er}, soit \cle{mehr/weniger}).
\end{corrige}

\begin{corrige}[Exercice 3 — Vocabulaire : trouve le mot]
\begin{enumerate}
\item \all{Am Ende des Monats sind meine Ausgaben höher als meine Einnahmen: ich habe ein Problem!} (mes dépenses sont plus élevées que mes recettes ; les deux noms sont au pluriel, d'où le verbe \all{sind})
\item \all{Unser Budget für die Klassenfahrt beträgt 50\,000 Euro.} (\all{betragen} = s'élever à ; \all{die Klassenfahrt} = le voyage scolaire)
\item \all{Ich muss jeden Monat etwas Geld sparen, um ein Handy zu kaufen.} (infinitif après \all{muss} ; \all{um ... zu} = pour)
\item \all{Unser größtes Projekt dieses Jahr ist die Reise nach Deutschland.} (accord du superlatif avec \all{das Projekt} : \all{größtes})
\item \all{Wir machen einen Plan für die nächsten drei Monate.} (accusatif masculin après \all{machen} : \all{einen Plan})
\end{enumerate}
\textbf{Points de vigilance :} vérifier le genre du nom avant d'accorder l'adjectif ou le déterminant ; \all{das Projekt} est neutre, \all{der Plan} est masculin. Noter \all{betragen} (verbe intransitif, pas de complément d'objet direct).
\end{corrige}

\begin{corrige}[Exercice 4 — Conjugaison : Futur I]
Rappel de la conjugaison de \all{werden} au présent : \all{ich werde, du wirst, er/sie/es wird, wir werden, ihr werdet, sie/Sie werden}.
\begin{enumerate}
\item \all{Ich werde nächstes Jahr mehr sparen.}
\item \all{Wir werden ein neues Projekt beginnen.}
\item \all{Wirst du dein Geld besser planen?} (dans une question directe sans mot interrogatif, le verbe conjugué est en première position)
\item \all{Die Schüler werden eine Reise nach Wien organisieren.}
\item \all{Mein Bruder wird weniger ausgeben.} (\all{er wird} ; l'infinitif \all{ausgeben} reste à la fin, la particule n'est pas séparée à l'infinitif)
\end{enumerate}
\textbf{Points de vigilance :} ne jamais écrire \cle{ich werde sparen werde} ni placer l'infinitif ailleurs qu'en fin de phrase ; \all{du wirst} avec \textbf{-st}, sans \textbf{t} supplémentaire. Le verbe \all{werden} sert aussi de verbe plein (\all{Er wird Arzt} = il devient médecin) : ne pas confondre avec \all{bekommen} (recevoir).
\end{corrige}

\begin{corrige}[Exercice 5 — Comparatif et superlatif : texte à trous]
\begin{enumerate}
\item \all{Ein Fahrrad ist billiger als ein Auto.} (comparatif régulier : adjectif + \textbf{-er} + \all{als})
\item \all{Das Leben in Douala ist teurer als auf dem Land.} (\all{teuer} perd son \textbf{e} : \all{teurer}, jamais \emph{teuerer})
\item \all{Mein Bruder spart mehr Geld als ich.} (irrégulier : \all{viel -- mehr -- am meisten})
\item \all{Ich kaufe lieber auf dem Markt als im Supermarkt.} (irrégulier : \all{gern -- lieber -- am liebsten})
\item \all{Der teuerste Monat für unsere Familie ist der Dezember.} (superlatif attributif décliné : \all{der teuerste Monat})
\item \all{Yaoundé ist eine der größten Städte Kameruns.} (\all{groß -- größer -- am größten} ; datif pluriel après \all{eine der} : \all{der größten Städte})
\item \all{Das interessanteste Projekt der Klasse ist die Deutschlandreise.} (accord avec \all{das Projekt} : \all{das interessanteste})
\item \all{Die Preise in der Stadt sind am höchsten.} (superlatif adverbial : \all{am höchsten} ; \all{hoch} perd son \textbf{c})
\end{enumerate}
\textbf{Points de vigilance :} distinguer le superlatif adverbial (\all{am höchsten}, après le verbe \all{sein}) du superlatif attributif décliné (\all{der teuerste Monat}, devant le nom) ; les adjectifs en \all{-el}, \all{-er}, \all{-en} perdent souvent leur \textbf{e} au comparatif (\all{teuer -- teurer}, \all{dunkel -- dunkler}, \all{trocken -- trockener}).
\end{corrige}

\begin{corrige}[Exercice 6 — Transformation : phrases avec \all{wenn}]
Rappel : quand la subordonnée en \all{wenn} est placée en tête, elle occupe la position 1 ; le verbe conjugué de la principale vient donc juste après la virgule, avant le sujet.
\begin{enumerate}
\item \all{Wenn er Geld spart, reist er nach Wien.}
\item \all{Wenn ich Zeit habe, helfe ich meinen Eltern.} (\all{helfen} + datif : \all{meinen Eltern})
\item \all{Wenn wir gut arbeiten, bekommen wir ein besseres Budget.}
\item \all{Wenn das Wetter schön ist, verkaufen wir mehr auf dem Markt.}
\item \all{Wenn du deine Ausgaben planst, hast du am Ende des Monats Geld.}
\item \all{Wenn die Schüler Deutsch lernen, können sie in Deutschland studieren.}
\end{enumerate}
\textbf{Points de vigilance :} l'erreur typique du francophone est de garder l'ordre sujet-verbe après la virgule (\emph{Wenn er Geld spart, er reist...}) : c'est faux, le verbe de la principale doit être en position 2, donc immédiatement après la virgule. Dans la subordonnée, le verbe est à la fin.
\end{corrige}

\begin{corrige}[Exercice 7 — Futur I : mise au temps et traduction]
\begin{enumerate}
\item \all{Ich werde jeden Monat 5\,000 Euro sparen.} — « J'économiserai 5\,000 euros chaque mois. »
\item \all{Wir werden unser Projekt im Januar beginnen.} — « Nous commencerons notre projet en janvier. »
\item \all{Meine Schwester wird einen neuen Computer kaufen.} — « Ma sœur achètera un nouvel ordinateur. »
\item \all{Die Familie wird ihr Budget besser planen.} — « La famille planifiera mieux son budget. »
\item \all{Ihr werdet nächstes Jahr nach Deutschland reisen.} — « Vous voyagerez en Allemagne l'année prochaine. »
\item \all{Er wird weniger Geld ausgeben.} — « Il dépensera moins d'argent. » (l'infinitif \all{ausgeben} se place à la fin, sans séparation de la particule)
\end{enumerate}
\textbf{Points de vigilance :} le verbe \all{werden} est le seul élément conjugué (position 2) ; tout le reste du verbal va à la fin sous forme d'infinitif. Ne pas confondre \all{werden} (auxiliaire du futur / devenir) avec \all{bekommen} (recevoir), faux ami de « devenir » : \all{Er wird Arzt.} = il deviendra médecin ; \all{Er bekommt ein Geschenk.} = il reçoit un cadeau.
\end{corrige}

\begin{corrige}[Exercice 8 — Appariement : mots et définitions]
\textbf{1 -- c)} : \all{das Taschengeld, -er} = l'argent de poche (littéralement « l'argent de la poche », \all{die Tasche, -n}).

\textbf{2 -- a)} : \all{das Budget, -s} = le budget, la somme disponible pour un projet.

\textbf{3 -- b)} : \all{die Einnahme, -n} = la recette, l'argent que l'on reçoit (de \all{einnehmen}, encaisser).

\textbf{4 -- e)} : \all{die Ausgabe, -n} = la dépense, l'argent que l'on dépense (de \all{ausgeben}, dépenser).

\textbf{5 -- d)} : \all{sparen} = économiser, mettre de l'argent de côté.

\textbf{6 -- f)} : \all{der Plan, die Pläne} = le plan, le programme d'actions organisé.

\textbf{Points de vigilance :} mémoriser les couples par familles de mots : \all{einnehmen -- die Einnahme} et \all{ausgeben -- die Ausgabe} ; cela aide à deviner le sens de mots nouveaux (\all{die Einnahmen steigen} = les recettes augmentent ; \all{die Ausgaben senken} = réduire les dépenses).
\end{corrige}

\begin{corrige}[Exercice 9 — Texte et questions de compréhension — type Bac]
\textit{Se reporter au texte support „Lernen, mit Geld umzugehen“ (non reproduit ici).}
\textbf{Fragen zum Text :}
\begin{enumerate}
\item \all{Zuerst muss man alle Einnahmen aufschreiben, dann alle Ausgaben.} (« Il faut d'abord noter toutes les recettes, puis toutes les dépenses. »)
\item \all{Falsch. Im Text steht: „Das Leben in der Stadt wird immer teurer.“} (le texte dit le contraire : la vie devient de plus en plus chère)
\item \all{Das wichtigste Projekt der Familie Ngo Bell ist die Ausbildung der Kinder.} (« Le projet le plus important de la famille Ngo Bell est l'éducation des enfants. »)
\item \all{Falsch. Der Vater sagt: „Wenn wir jeden Monat etwas sparen, werden wir nächstes Jahr einen Computer für unsere Tochter kaufen.“} (c'est un ordinateur, pas un téléphone portable)
\item \all{Die Experten sagen: Wer früh sparen lernt, wird später weniger Probleme mit Geld haben.} (« Celui qui apprend tôt à économiser aura plus tard moins de problèmes d'argent. »)
\end{enumerate}
\textbf{Fragen zur Sprache :}
\begin{enumerate}
\item Comparatifs dans le texte : \all{höher}, \all{schwieriger}, \all{teurer}, \all{weniger}. Superlatifs : \all{das wichtigste Projekt}, \all{am interessantesten}. Phrases personnelles possibles : \all{Das Leben in Yaoundé ist teurer als in meinem Dorf.} / \all{Mein wichtigstes Projekt ist das Abitur.}
\item \all{Die Familie wird einen Plan machen.}
\end{enumerate}
\textbf{Barème indicatif (10 points) :} questions de compréhension : 1 point par réponse correcte en phrase complète (5 points) ; questions de langue : 2 points chacune (comparatif + superlatif correctement identifiés et réemployés ; Futur I correct). \textbf{Points de vigilance :} répondre par des phrases complètes avec verbe en position 2 ; pour \all{richtig/falsch}, la justification doit citer le texte (guillemets „…“) ; ne pas recopier une phrase qui ne répond pas à la question.
\end{corrige}

\begin{corrige}[Exercice 10 — Thème et version — type Bac]
\textbf{A. Thème (Français $\rightarrow$ Allemand).}
\begin{enumerate}
\item \all{Wenn ich genug Geld habe, werde ich einen neuen Computer kaufen.} (subordonnée en \all{wenn} au présent, verbe à la fin ; principale au Futur I avec \all{werde} juste après la virgule)
\item \all{Mein wichtigstes Projekt ist die Reise nach Deutschland.} (superlatif attributif : \all{wichtigstes}, accord avec \all{das Projekt})
\item \all{Das Leben ist in der Stadt teurer als auf dem Land.} (« au village » = \all{auf dem Land}, datif ; « que » de comparaison = \all{als})
\item \all{Wenn ich groß bin, werde ich mein Budget besser verwalten.} (l'allemand emploie le présent après \all{wenn} là où le français met un futur ; « gérer » = \all{verwalten} ; on accepte aussi \all{... werde ich besser mit Geld umgehen.})
\item \all{Meine Schwester spart mehr Geld als ich, weil sie weniger ausgibt.} (\all{weil} renvoie le verbe conjugué à la fin ; \all{ausgeben} : \all{sie gibt ... aus}, particule détachée en fin de subordonnée)
\end{enumerate}

\textbf{B. Version (Allemand $\rightarrow$ Français).}
\textit{Texte source :}
\all{Die Familie Müller aus München plant ihr Budget jeden Monat genau. Die Einnahmen sind höher als die Ausgaben, aber das Leben in München ist teurer als in vielen anderen deutschen Städten. Das größte Projekt der Familie ist der Kauf eines Hauses. Wenn die Familie jeden Monat 500 Euro spart, wird sie in fünf Jahren genug Geld haben. „Sparen ist manchmal schwer, aber es ist das Wichtigste für unsere Zukunft“, sagt Frau Müller.}

\textit{Traduction proposée :}
« La famille Müller, de Munich, planifie son budget avec précision chaque mois. Les recettes sont plus élevées que les dépenses, mais la vie à Munich est plus chère que dans beaucoup d'autres villes allemandes. Le plus grand projet de la famille est l'achat d'une maison. Si la famille économise 500 euros chaque mois, elle aura assez d'argent dans cinq ans. "Économiser est parfois difficile, mais c'est le plus important pour notre avenir", dit Mme Müller. »

\textbf{Barème indicatif (10 points) :} thème : 1 point par phrase (5 points) ; version : 5 points (respect du sens 3 points, qualité du français 2 points). \textbf{Points de vigilance :} en thème, soigner la place du verbe dans chaque type de phrase (position 2 en principale, fin en subordonnée) ; en version, rendre \all{genau} par « avec précision » et \all{die Zukunft} par « l'avenir », sans calquer l'ordre allemand.
\end{corrige}

\begin{corrige}[Exercice 11 — Production écrite — type Bac]
Voici un texte modèle (les quatre structures exigées sont indiquées entre parenthèses) :

\all{Ich bekomme von meinen Eltern jeden Monat ein kleines Taschengeld. Meine größten Ausgaben sind der Transport mit dem Bendskin und das Mittagessen in der Schule. Das Leben in Yaoundé ist teurer als früher (comparatif), deshalb muss ich mein Budget gut planen. Am Anfang jedes Monats schreibe ich alle Einnahmen und alle Ausgaben auf. Mein wichtigstes Projekt (superlatif) ist der Kauf eines Computers, denn ich möchte nach dem Abitur Informatik studieren. Wenn ich jeden Monat 5\,000 Francs spare (subordonnée avec wenn), werde ich nächstes Jahr genug Geld haben (Futur I). Sparen ist nicht immer leicht, aber es ist das Wichtigste für meine Zukunft.}

Traduction du modèle : « Je reçois de mes parents une petite somme d'argent de poche chaque mois. Mes plus grandes dépenses sont le transport en bendskin et le déjeuner à l'école. La vie à Yaoundé est plus chère qu'avant, c'est pourquoi je dois bien planifier mon budget. Au début de chaque mois, je note toutes les recettes et toutes les dépenses. Mon projet le plus important est l'achat d'un ordinateur, car je veux étudier l'informatique après le baccalauréat. Si j'économise 5\,000 francs chaque mois, j'aurai assez d'argent l'année prochaine. Économiser n'est pas toujours facile, mais c'est le plus important pour mon avenir. »

\textbf{Barème indicatif (10 points) :} respect des 4 contraintes grammaticales (4 points, 1 par structure correcte et soulignée) ; vocabulaire du thème (2 points) ; correction grammaticale générale : place du verbe, accords, cas (2 points) ; cohérence, longueur et présentation (2 points).

\textbf{Points de vigilance :} vérifier que la subordonnée en \all{wenn} a bien son verbe à la fin (\all{spare}) et que la principale commence par le verbe conjugué (\all{werde}) ; au Futur I, un seul verbe conjugué (\all{werden}) et l'infinitif à la fin (\all{haben}) ; superlatif attributif décliné devant le nom (\all{mein wichtigstes Projekt}) ; penser à la majuscule à tous les noms (\all{Geld, Budget, Projekt, Ausgaben, Einnahmen, Bendskin, Yaoundé, Abitur, Informatik, Francs, Zukunft}).
\end{corrige}

\newpage

\coursec{Fiche bilan}

\begin{popfigure}
\begin{center}
\begin{tikzpicture}[mindmap, grow cyclic, every node/.style={concept, align=center, font=\small},
root concept/.append style={concept color=popblue, font=\bfseries\small, minimum size=3.2cm},
level 1 concept/.append style={sibling angle=60, level distance=3.6cm, minimum size=2.4cm, font=\footnotesize}]
\node[root concept, align=center]{Budget und Projekte\\ (AL06)}
child[concept color=poporange]{node[align=center]{Lexique\\ Geld, Budget\\ Einnahmen,\\ Ausgaben}}
child[concept color=popgreen]{node[align=center]{Comparatif\\ billig -- billiger\\ -- am billigsten}}
child[concept color=poppurple]{node[align=center]{Superlatif\\ der beste\\ am besten}}
child[concept color=poppink]{node[align=center]{Conditionnelles\\ wenn + verbe\\ en fin de phrase}}
child[concept color=popgold]{node[align=center]{Futur I\\ werden +\\ Infinitiv}}
child[concept color=popblue!60]{node[align=center]{Méthode\\ lire -- repérer\\ -- résumer}};
\end{tikzpicture}
\end{center}
\legende{Carte mentale de la leçon : la gestion du budget et des projets}
\end{popfigure}

\begin{aretenir}[L'essentiel de la leçon AL06]
\textbf{1. Lexique clé (à connaître avec article et pluriel)}
\begin{itemize}
\item \all{das Geld} (Sg.) : l'argent ; \all{das Budget, -s} : le budget ; \all{der Plan, \"-e} : le plan, le projet
\item \all{die Einnahme, -n} : la recette, le revenu ; \all{die Ausgabe, -n} : la dépense
\item \all{das Projekt, -e} : le projet ; \all{die Rechnung, -en} : la facture ; \all{das Konto, -en} (ou \all{die Konten}) : le compte
\item \all{sparen} : économiser ; \all{ausgeben (gibt aus, gab aus, hat ausgegeben)} : dépenser ; \all{verdienen} : gagner ; \all{planen} : planifier ; \all{leihen (leiht, lieh, hat geliehen)} : prêter ou emprunter
\item Expressions : \all{mit Geld umgehen} : gérer l'argent ; \all{ein Budget aufstellen} : établir un budget ; \all{Geld f\"ur ein Projekt sparen} : économiser pour un projet
\end{itemize}

\textbf{2. Grammaire à connaître par c\oe ur}
\begin{center}
\begin{tabularx}{\linewidth}{|l|X|X|}
\hline
\rowcolor{popblueL} Point & Règle & Exemple \\
\hline
Comparatif & adjectif + \textbf{-er} (+ \all{als}) ; umlaut sur a/o/u des monosyllabes fréquents & \all{gut -- besser} ; \all{alt -- \"alter als} \\
\hline
Superlatif & \textbf{am} + adjectif + \textbf{-sten} (attribut : der/die/das + adj. + \textbf{-ste}, décliné) & \all{am billigsten} ; \all{das beste Projekt} \\
\hline
Irréguliers & \all{gut -- besser -- am besten} ; \all{viel -- mehr -- am meisten} ; \all{gern -- lieber -- am liebsten} & \all{Ich spare lieber.} \\
\hline
Conditionnelle & \all{wenn} + verbe conjugué \textbf{en fin} de subordonnée ; si la subordonnée ouvre la phrase, le verbe principal suit immédiatement (inversion) & \all{Wenn ich Geld spare, kaufe ich ein Fahrrad.} \\
\hline
Futur I & \all{werden} (conjugué) + infinitif \textbf{à la fin} & \all{Ich werde ein Budget aufstellen.} \\
\hline
\end{tabularx}
\end{center}

\textbf{3. Méthodes express}
\begin{itemize}
\item \textbf{Compréhension} : lire le titre et la source, repérer le thème, souligner les chiffres et les mots du champ lexical de l'argent, répondre en phrases complètes.
\item \textbf{Résumé} : une phrase par idée principale, au présent, avec ses propres mots, sans citation longue.
\item \textbf{Production} : 5 à 8 phrases ; varier les connecteurs (\all{zuerst, dann, danach, schlie\ss lich}) ; utiliser au moins un comparatif, une conditionnelle avec \all{wenn} et une phrase au Futur I.
\end{itemize}
\end{aretenir}

\begin{propriete}[Rappels de conjugaison : werden au présent (auxiliaire du Futur I)]
\begin{center}
\begin{tabular}{|l|l|}
\hline
\rowcolor{popgoldL} Personne & Forme \\
\hline
ich & werde \\
\hline
du & wirst \\
\hline
er/sie/es & wird \\
\hline
wir & werden \\
\hline
ihr & werdet \\
\hline
sie/Sie & werden \\
\hline
\end{tabular}
\end{center}
Exemple : \all{Wir werden n\"achstes Jahr ein Projekt planen.} « Nous planifierons un projet l'année prochaine. » Attention : \all{werden} conjugué se place en position 2, l'infinitif à la toute fin de la phrase.
\end{propriete}

\begin{attention}[Pièges fréquents des francophones]
\begin{itemize}
\item Ne pas confondre \all{wenn} (condition ou répétition : « si, quand ») et \all{als} (moment unique au passé : « quand »). Ici, pour la condition : toujours \all{wenn}.
\item Après \all{wenn} en tête de phrase, le verbe principal vient juste après la virgule : \all{Wenn ich Geld habe, spare ich.} et non « \all{Wenn ich Geld habe, ich spare.} ».
\item Le comparatif se forme avec \all{als} (« que »), jamais avec \all{wie} : \all{Das Fahrrad ist billiger als das Auto.}
\item Majuscule obligatoire à tous les noms : \all{das Geld, der Plan, die Ausgabe}.
\item \all{leihen} signifie à la fois « prêter » et « emprunter » ; le sens dépend du contexte : \all{Ich leihe mir Geld.} « J'emprunte de l'argent. »
\end{itemize}
\end{attention}

\begin{exemplebox}[Mini-production modèle : Mein Budget, mein Projekt]
\all{Zuerst stelle ich jeden Monat ein Budget auf. Dann schreibe ich meine Einnahmen und meine Ausgaben auf. Wenn ich weniger ausgebe, kann ich mehr sparen. Mein Projekt ist ein Fahrrad, denn es ist billiger als ein Motorrad und besser f\"ur die Gesundheit. N\"achstes Jahr werde ich das Fahrrad kaufen. Schlie\ss lich bin ich stolz, weil ich gut mit Geld umgehen kann.}\par
« D'abord, j'établis un budget chaque mois. Ensuite, je note mes recettes et mes dépenses. Si je dépense moins, je peux économiser davantage. Mon projet est un vélo, car il est moins cher qu'une moto et meilleur pour la santé. L'année prochaine, j'achèterai le vélo. Enfin, je suis fier de savoir bien gérer l'argent. »
\end{exemplebox}

\begin{aretenir}[Checklist avant le Bac]
\begin{itemize}
\item[$\square$] Je connais le lexique du budget avec articles et pluriels (\all{die Einnahme, -n} ; \all{die Ausgabe, -n}).
\item[$\square$] Je sais former le comparatif : adjectif + \all{-er}, avec umlaut (\all{alt -- \"alter}).
\item[$\square$] Je sais former le superlatif : \all{am -sten} et \all{der/die/das -ste}.
\item[$\square$] Je connais par c\oe ur les formes irrégulières : \all{gut -- besser -- am besten}, \all{viel -- mehr -- am meisten}, \all{gern -- lieber -- am liebsten}.
\item[$\square$] Je place le verbe conjugué à la fin d'une subordonnée introduite par \all{wenn}.
\item[$\square$] Quand la subordonnée \all{wenn} ouvre la phrase, je mets le verbe principal juste après la virgule (inversion).
\item[$\square$] Je conjugue \all{werden} sans erreur : \all{ich werde, du wirst, er/sie/es wird, wir werden, ihr werdet, sie/Sie werden}.
\item[$\square$] Au Futur I, je place l'infinitif à la fin de la phrase.
\item[$\square$] Je mets une majuscule à tous les noms allemands (\all{das Geld, der Plan}).
\item[$\square$] Je sais répondre à une question de compréhension par une phrase complète en allemand.
\item[$\square$] Je sais rédiger un court paragraphe « \all{Mein Budget, mein Projekt} » avec connecteurs et au moins trois points de grammaire de la leçon.
\item[$\square$] Je relis ma production : place du verbe, accords, orthographe des Umlaute et du \all{\ss}.
\end{itemize}
\end{aretenir}

\findecours

\end{document}
